% ThinkAlgo — Hình học và Đo lường — Lớp 7
% Edition: teacher (with solutions), VI.  Generated by build_grade7.py — edit body/ and rebuild.
\documentclass[11pt]{article}
\newif\ifsol \soltrue
\usepackage{iftex}
\ifPDFTeX
  \usepackage[utf8]{inputenc}\usepackage[T5]{fontenc}
\else
  \usepackage{fontspec}
\fi
\usepackage[a4paper,top=2.95cm,headheight=1.5cm,headsep=0.45cm,bottom=3.1cm,left=1.9cm,right=1.9cm]{geometry}
\usepackage{amsmath,amssymb}
\usepackage{tikz}\usetikzlibrary{calc,shadows,shapes.geometric}
\usepackage[most]{tcolorbox}
\usepackage{enumitem}
\usepackage{array}
\usepackage{tabularx}
\usepackage{colortbl}
\usepackage{comment}
\usepackage{needspace}
\usepackage{fancyhdr}
\usepackage{fontawesome5}
\usepackage{graphicx}
\usepackage{microtype}
% PDF outline (bookmarks) so each chapter shows in the reader's outline — does not change the page layout
\usepackage[hidelinks,unicode,bookmarksnumbered=false,bookmarksopen=true,bookmarksopenlevel=0]{hyperref}
\usepackage{bookmark}
\usetikzlibrary{decorations.pathreplacing,arrows.meta,angles,quotes,patterns,positioning,through,intersections}

% ---------------- palette: navy family + subtle light tint; orange ONLY in the logo ----------------
% ----- EXACT palette from the user's original main (36).tex (vivid, not pale) -----
\definecolor{tmBlue}{RGB}{66,133,244}   % problem / header / section
\definecolor{tmGreen}{RGB}{52,168,83}   % solution
\definecolor{tmOrange}{RGB}{255,152,0}  % example
\colorlet{Navy}{tmBlue}                 % header rule / section heading / badge / divider
\definecolor{NavyD}{HTML}{222222}       % solution body text (dark, readable)
\colorlet{Gold}{Navy}
\colorlet{ProbBlue}{tmBlue}             % problem boxes
\colorlet{SolGreen}{tmGreen}            % solution boxes
\definecolor{AnsRed}{HTML}{E02424}      % vivid red for answers / totals
\definecolor{CardGreen}{HTML}{34A853}   % digit-card green
\definecolor{StarGold}{HTML}{F2B134}    % difficulty star (gold)
\definecolor{KidOrange}{HTML}{DE8A2C}   % logo accent only
\definecolor{Ink}{HTML}{1F2733}       % body text
\definecolor{Slate}{HTML}{5B6A82}     % muted labels
\definecolor{Line}{HTML}{C6D0DF}      % light hairlines / student boxes
\definecolor{Tint}{HTML}{EDF1F7}      % pale panel fill (worked solution / answer)
% aliases so older names keep working
\colorlet{TextDark}{Ink}\colorlet{TextGray}{Slate}\colorlet{Soft}{Slate}
\colorlet{CardLine}{Line}\colorlet{Teal}{Navy}\colorlet{ProbB}{Navy}\colorlet{SolG}{Navy}
% bright, kid-friendly colours for bar-model (sơ đồ đoạn thẳng) blocks
\definecolor{tmPurple}{RGB}{124,77,226}
\definecolor{tmTeal}{RGB}{0,150,136}
\definecolor{tmPink}{RGB}{233,30,99}
\AtBeginDocument{\color{Ink}}
\raggedbottom  % natural flow, even spacing (no stretched gaps)
\setlength{\parindent}{0pt}\setlength{\parskip}{0.18em}
\renewcommand{\arraystretch}{1.35}
\providecommand{\AnsWord}{Đáp số}

% ---------------- brand logo (original mark) ----------------
\newcommand{\logomark}[1]{\begin{tikzpicture}[baseline=-0.55ex,scale=#1]
  \draw[KidOrange,line width=1.5pt] (-0.27,0) .. controls (-0.27,0.23) and (0,0.23) .. (0,0)
        .. controls (0,-0.23) and (-0.27,-0.23) .. (-0.27,0);
  \draw[Navy,line width=1.5pt] (0.27,0) .. controls (0.27,0.23) and (0,0.23) .. (0,0)
        .. controls (0,-0.23) and (0.27,-0.23) .. (0.27,0);
\end{tikzpicture}}
% real brand logo image (gold infinity + "think algo" wordmark)
\newcommand{\logohead}{\IfFileExists{logo.png}{\includegraphics[height=1.35cm]{logo.png}}%
  {\raisebox{0.35cm}{\logomark{1.5}\hspace{0.5em}{\Large\bfseries\color{KidOrange}think\,\color{Navy}algo}}}}

% ---------------- running header (logo + Name/Date + topic) on EVERY page, single navy rule ----------------
\pagestyle{fancy}\fancyhf{}
\renewcommand{\headrulewidth}{1.4pt}\renewcommand{\footrulewidth}{0pt}
\renewcommand{\headrule}{{\color{Navy}\hrule height 1.3pt}}
\fancyhead[L]{\raisebox{-0.95cm}{\logohead}}
\fancyhead[R]{\setlength{\tabcolsep}{0pt}%
  \begin{tabular}{@{}r@{}}
    {\small\color{Ink}{\color{tmBlue}\faUserEdit}~\NameWord:~\rule{3.0cm}{0.5pt}\quad{\color{tmBlue}\faCalendarDay}~\DateWord:~\rule{1.9cm}{0.5pt}}\\[5pt]
    {\large\bfseries\color{Navy}\CurTopic}
  \end{tabular}}

% ---------------- footer (rounded page border + Topic/Grade + thinkalgo.com + page + tagline) ----------------
% NOTE: a remember-picture overlay must be placed DIRECTLY in \fancyfoot, NOT inside \def/\newcommand.
\fancyfoot[C]{\begin{tikzpicture}[remember picture, overlay]
  \draw[Navy!40, line width=1.4pt, rounded corners=12pt]
     ($(current page.north west)+(0.8cm,-0.8cm)$) rectangle ($(current page.south east)+(-0.8cm,0.45cm)$);
  \coordinate (FL) at ($(current page.south west)+(1.2cm, 2.18cm)$);
  \coordinate (FR) at ($(current page.south east)+(-1.2cm, 2.18cm)$);
  \draw[Navy!55, line width=1pt] (FL) -- (FR);
  \node[fill=white, draw=Navy, line width=1.2pt, rounded corners=8pt, font=\bfseries\small,
        text=Navy, inner xsep=11pt, inner ysep=4pt] at ($(current page.south)+(0,2.18cm)$) {{\color{tmBlue}\faGlobe}\ \ thinkalgo.com};
  \node[text=Slate, font=\footnotesize, anchor=north west, align=left, inner sep=0pt, text width=7.6cm]
        at ($(FL)+(0.1cm,-0.30cm)$) {{\color{tmBlue}\faTag}\ \TopicWord: \textbf{\color{Ink}\CurTopic}\\[1.5pt]{\color{tmBlue}\faGraduationCap}\ \DocGrade};
  \node[text=Slate, font=\footnotesize] at ($(current page.south)+(0,1.20cm)$) {\thepage};
  \node[text=Slate, font=\scriptsize, anchor=north east, align=right, inner sep=0pt, text width=7.4cm]
        at ($(FR)+(-0.1cm,-0.34cm)$) {{\normalsize\bfseries\color{Navy}ThinkAlgo}\ — From Math Thinking to Code};
\end{tikzpicture}}

% ---------------- page-1 lead-in (topic is already in the header; keep a slim spacer) ----------------
\newcommand{\titlehead}{\vspace*{-0.2cm}}
\providecommand{\CurTopic}{}

% chapter banner + sets the running-header topic. #1 = number, #2 = title
\newcommand{\chapterhead}[2]{\renewcommand{\CurTopic}{#2}%
  \setcounter{seccnt}{0}% each chapter is independent: section badges restart at 1
  \phantomsection\pdfbookmark[0]{#1. #2}{chap#1}\label{chap:#1}% add this chapter to the PDF outline + TOC ref
  \par\vspace{0.1cm}\noindent\begin{tikzpicture}[baseline=-0.5ex]
    \node[circle,fill=tmBlue,text=white,font=\Large\bfseries,minimum size=0.95cm,inner sep=0pt](cn){#1};
    \node[anchor=west,font=\LARGE\bfseries,text=tmBlue] at ([xshift=12pt]cn.east){#2};
  \end{tikzpicture}\par\vspace{4pt}\noindent\textcolor{tmBlue}{\rule{\linewidth}{2pt}}\par\vspace{0.5cm}}

% ---------------- section heading (outline navy number badge, no fill) ----------------
\newcounter{seccnt}
\newcommand{\shead}[1]{\par\vspace{0.7em}\refstepcounter{seccnt}%
  \noindent\begin{tikzpicture}[baseline=-0.6ex]
    \node[circle,draw=Gold,line width=1.2pt,text=Gold,font=\bfseries,minimum size=1.05em,inner sep=1pt]{\theseccnt};
  \end{tikzpicture}\hspace{0.4em}{\large\bfseries\color{Navy} #1}\par\vspace{0.25em}}

% ============ boxes EXACTLY per the original main (36).tex: solid-blue PROBLEM, solid-orange EXAMPLE,
%              outline-green SOLUTION, all with drop shadow + rounded corners ============
% ---------------- knowledge box: white, tmBlue border + left accent + drop shadow ----------------
\newtcolorbox{knowbox}{enhanced,breakable,colback=white,colframe=tmBlue!70,boxrule=0.9pt,arc=5pt,
  left=12pt,right=11pt,top=6pt,bottom=6pt,drop shadow,
  borderline west={3.5pt}{0pt}{tmBlue}}

% ---------------- EXAMPLE (Ví dụ): SOLID ORANGE tab, white text ----------------
\newtcolorbox{exbox}[1]{enhanced,breakable,colback=white,colframe=tmOrange,boxrule=1pt,arc=5pt,
  left=12pt,right=11pt,top=14pt,bottom=7pt,drop shadow,
  attach boxed title to top left={xshift=10pt,yshift=-\tcboxedtitleheight/2},
  boxed title style={colback=tmOrange,colframe=tmOrange,arc=4pt,boxrule=0pt,drop shadow},
  title={#1},fonttitle=\bfseries\color{white}}
\newcommand{\guide}{\par\vspace{2pt}{\bfseries\color{tmBlue}\GuideWord}\par\nopagebreak}

% ---------------- PROBLEM (Bài): SOLID BLUE tab, white text ----------------
\newtcolorbox{probbox}[2]{enhanced,breakable,colback=white,colframe=tmBlue,boxrule=1pt,arc=5pt,
  left=12pt,right=10pt,top=14pt,bottom=7pt,drop shadow,
  attach boxed title to top left={xshift=10pt,yshift=-\tcboxedtitleheight/2},
  boxed title style={colback=tmBlue,colframe=tmBlue,arc=4pt,boxrule=0pt,drop shadow},
  fonttitle=\bfseries\color{white},title={\ProbWord\ #1\ifx\relax#2\relax\else#2\fi}}

% ---------------- SOLUTION (Lời giải): OUTLINE GREEN tab (white fill, green border + text) ----------------
% Teacher: a green tcolorbox. Student: the body is dropped entirely (comment pkg — robust with tables/\\).
\ifsol
  \newtcolorbox{solbox}{enhanced,breakable,colback=white,colframe=tmGreen,boxrule=1pt,arc=5pt,
    left=12pt,right=10pt,top=14pt,bottom=7pt,drop shadow,
    attach boxed title to top left={xshift=10pt,yshift=-\tcboxedtitleheight/2},
    boxed title style={colback=white,colframe=tmGreen,boxrule=1pt,arc=4pt,drop shadow},
    fonttitle=\bfseries\color{tmGreen!88!black},title={\SolTitle},before upper={\color{NavyD}}}
\else
  \excludecomment{solbox}
\fi

% ---------------- answer highlight: label green, value RED ----------------
\newcommand{\ans}[1]{\par\smallskip{\bfseries\color{tmGreen!88!black}\AnsWord:}\ {\bfseries\color{AnsRed}#1}\par}

% ---------------- digit cards (exact original style: fill=c!10, border=c!60!black, drop shadow) ----------------
% \digitcards{0/gray,1/cyan,2/green,3/cyan} : value/colour pairs
\newcommand{\digitcards}[1]{\begin{center}\begin{tikzpicture}[scale=0.7,transform shape,font=\sffamily]
  \def\dx{0}\foreach \v/\c in {#1}{
    \node[draw=\c!60!black,line width=1.5pt,fill=\c!10,rounded corners=4pt,
          minimum width=1cm,minimum height=1.4cm,drop shadow] at (\dx,0){\Huge\bfseries \v};
    \pgfmathsetmacro\dx{\dx+1.5}\xdef\dx{\dx}}
\end{tikzpicture}\end{center}}

% ---------------- student working space (white, light border) ----------------
\newcommand{\worksp}[1]{\ifsol\else\begin{tcolorbox}[colback=white,colframe=Line,boxrule=0.8pt,
  arc=5pt,height=\dimexpr#1+0.05cm\relax,left=8pt,right=8pt,top=4pt,bottom=4pt]{\footnotesize\color{Slate}\itshape\WorkWord}\end{tcolorbox}\fi}

% ---------------- \probsol{solution} : Teacher shows the green solution box; Student shows a BLANK
%   "Bài làm" box whose height matches that solution's height (writing room = work needed), so both
%   versions have identical per-problem height and paginate the same. ----------------
\newsavebox{\solmeasure}\newlength{\solht}
\newcommand{\probsol}[1]{%
  \begin{lrbox}{\solmeasure}\begin{minipage}{\dimexpr\linewidth-24pt\relax}#1\end{minipage}\end{lrbox}%
  \ifsol
    \begin{solbox}#1\end{solbox}%
  \else
    \setlength{\solht}{\dimexpr\ht\solmeasure+\dp\solmeasure+0.85cm\relax}%
    \ifdim\solht<1.35cm\setlength{\solht}{1.35cm}\fi
    \ifdim\solht>\StudentMaxWork\setlength{\solht}{\StudentMaxWork}\fi
    \begin{tcolorbox}[colback=white,colframe=Line,boxrule=0.8pt,arc=5pt,
      left=8pt,right=8pt,top=4pt,bottom=4pt,height=\solht]{\footnotesize\color{Slate}\itshape\WorkWord}\end{tcolorbox}%
  \fi}

% ---------------- practice divider: blue left accent bar, no fill ----------------
\newcommand{\dividerbar}[1]{\par\vspace{0.6em}\noindent
  \begin{tcolorbox}[enhanced,frame hidden,interior hidden,borderline west={5pt}{0pt}{ProbBlue},
    colback=white,left=12pt,right=12pt,top=5pt,bottom=5pt,arc=4pt]
  {\large\bfseries\color{ProbBlue} #1}\end{tcolorbox}\vspace{0.4em}}

% place-value choice-count diagram (Cambridge style): #1 = "label/count/note/color" list, #2 = result
% notes must NOT contain commas or slashes (they separate \foreach fields)
\newcommand{\countboxes}[2]{\begin{center}\begin{tikzpicture}[scale=0.62,transform shape,font=\sffamily]
  \def\xx{0}%
  \foreach \lab/\cnt/\note/\col [count=\i] in {#1}{%
    \ifnum\i>1\relax \node[text=gray,font=\Huge] at (\xx-1.8,0){$\times$};\fi
    \node[draw=\col!65,fill=\col!8,rounded corners=3pt,line width=1pt,minimum width=2.7cm,minimum height=1.7cm] at (\xx,0){};
    \node[text=\col!85,font=\bfseries] at (\xx,0.58){\lab};
    \node[text=\col!88,font=\Huge\bfseries] at (\xx,-0.1){\cnt};
    \node[text=gray,font=\scriptsize] at (\xx,-1.05){\note};
    \pgfmathsetmacro\xx{\xx+3.6}\xdef\xx{\xx}}%
  \node[text=gray,font=\Huge] at (\xx-1.8,0){$=$};
  \node[text=red,font=\Huge\bfseries] at (\xx,0){#2};
\end{tikzpicture}\end{center}}

% ======== bar model (sơ đồ đoạn thẳng): lightly-rounded blocks, darker border, tiny gaps, bright colours ========
% one block: \blk{x1}{x2}{yc}{height}{colour}{label}  (a small gap is left on each side)
\newcommand{\blk}[6]{%
  \draw[rounded corners=1.3pt,fill=#5!16,draw=#5!80!black,line width=1pt]
    (#1+0.045,#3-#4/2) rectangle (#2-0.045,#3+#4/2);
  \node[font=\small\bfseries,text=#5!52!black] at ({(#1+#2)/2},#3){#6};}
% measurement with a double-headed ARROW (not a brace) + label above the span: \brcup{x1}{x2}{y}{label}
\newcommand{\brcup}[4]{%
  \draw[<->,>=stealth,tmBlue!72!black,line width=0.7pt] (#1+0.05,#3)--(#2-0.05,#3);
  \node[above=-1.5pt,font=\footnotesize,text=tmBlue!58!black] at ({(#1+#2)/2},#3){#4};}
% measurement with a double-headed ARROW + label below the span: \brcdn{x1}{x2}{y}{label}
\newcommand{\brcdn}[4]{%
  \draw[<->,>=stealth,tmBlue!72!black,line width=0.7pt] (#1+0.05,#3)--(#2-0.05,#3);
  \node[below=-1.5pt,font=\footnotesize,text=tmBlue!58!black] at ({(#1+#2)/2},#3){#4};}
% row label at the left of a bar: \rowlab{x}{y}{text}
\newcommand{\rowlab}[3]{\node[anchor=east,font=\small\bfseries,text=Ink] at (#1,#2){#3};}
% vertical right-pointing CURLY BRACE spanning two stacked bars: \vbrace{x}{ytop}{ybot}{label}
\newcommand{\vbrace}[4]{%
  \pgfmathsetmacro\ym{(#2+#3)/2}%
  \draw[tmBlue!72!black,line width=0.9pt,line cap=round]
    (#1,#2) .. controls (#1+0.18,#2) and (#1+0.18,#2-0.18) .. (#1+0.18,#2-0.22)
    -- (#1+0.18,{\ym+0.22})
    .. controls (#1+0.18,\ym) and (#1+0.30,{\ym+0.08}) .. (#1+0.42,\ym)
    .. controls (#1+0.30,{\ym-0.08}) and (#1+0.18,\ym) .. (#1+0.18,{\ym-0.22})
    -- (#1+0.18,#3+0.22)
    .. controls (#1+0.18,#3+0.18) and (#1+0.18,#3) .. (#1,#3);
  \node[right=4pt,font=\footnotesize,text=tmBlue!58!black] at (#1+0.42,\ym){#4};}
% blank 3x3 grid for a problem statement (student fills it in)
\newcommand{\blankgrid}{\begin{tikzpicture}[font=\sffamily]
  \foreach \r in {0,1,2}{\foreach \c in {0,1,2}{
    \draw[rounded corners=1.3pt,fill=Tint,draw=Slate!55,line width=0.9pt]
      (\c+0.05,\r+0.05) rectangle (\c+0.95,\r+0.95);}}
\end{tikzpicture}}
% wrapper: LEFT-aligned (not centred), a little air above/below
\newenvironment{barfig}{\par\vspace{3pt}\noindent\kern0.35cm\begin{tikzpicture}[font=\sffamily,>=stealth]}%
  {\end{tikzpicture}\par\vspace{3pt}}
% wrapper: CENTRED (for illustrative diagrams in theory / examples)
\newenvironment{barfigC}{\par\vspace{3pt}\begin{center}\begin{tikzpicture}[font=\sffamily,>=stealth]}%
  {\end{tikzpicture}\end{center}\vspace{2pt}}

% ---------------- difficulty stars (nicely drawn) : \stars{1}=hard, \stars{2}=very hard ----------------
\newcommand{\diffstar}{\tikz[baseline=-0.55ex]{\node[star,star points=5,star point ratio=0.46,
  draw=StarGold!72!black,fill=StarGold,line width=0.3pt,inner sep=1.75pt]{};}}
\newcommand{\stars}[1]{\ifcase#1\or\kern2pt\diffstar\or\kern2pt\diffstar\kern1.5pt\diffstar\fi}

\newcommand{\nd}[1]{\overline{#1}}

% number line: #1 last index, #2 ticks "i/label", #3 open (excluded) indices, #4 filled (solution) indices
\newcommand{\numlinex}[4]{\begin{center}\begin{tikzpicture}[x=1.55cm,>=stealth]
  \draw[->,thick,Navy](-0.45,0)--(#1+0.45,0);
  \foreach \i/\lab in {#2}{\draw[Navy,line width=0.8pt](\i,0.08)--(\i,-0.08);
     \node[below,font=\small,text=Ink]at(\i,-0.1){\lab};}
  \foreach \i in {#3}{\draw[Navy,fill=white,line width=1pt](\i,0)circle(2.3pt);}
  \foreach \i in {#4}{\fill[KidOrange](\i,0)circle(2.7pt);}
\end{tikzpicture}\end{center}}


% ======================================================================
% Grade 7 additions (shared by all four editions)
% ======================================================================
% student edition: blank "work" box height = min(solution height, this cap)
\newlength{\StudentMaxWork}\setlength{\StudentMaxWork}{7.5cm}

% ThinkAlgo tip box (label set per language by \TipWord)
\newtcolorbox{tipbox}{enhanced,breakable,colback=StarGold!8,colframe=StarGold!75!black,boxrule=1pt,arc=4pt,
  left=11pt,right=10pt,top=6pt,bottom=6pt,
  before upper={\textbf{\color{StarGold!65!black}\faLightbulb\ \TipWord}\ }}

% "Recall" box: re-states a piece of theory inside a solution / example
\newtcolorbox{recallbox}{enhanced,breakable,colback=tmBlue!5,colframe=tmBlue!45,boxrule=0.6pt,arc=3pt,
  left=8pt,right=8pt,top=4pt,bottom=4pt,
  before upper={\textbf{\color{tmBlue!80!black}\faBookOpen\ \RecallWord}\ }}

% "Why is it true? / Proof" box
\newtcolorbox{proofbox}{enhanced,breakable,colback=tmPurple!4,colframe=tmPurple!45,boxrule=0.6pt,arc=3pt,
  left=8pt,right=8pt,top=4pt,bottom=4pt,
  before upper={\textbf{\color{tmPurple!85!black}\faSearch\ \ProofWord}\ },
  after upper={\hfill$\square$}}

% "Common mistake" box
\newtcolorbox{warnbox}{enhanced,breakable,colback=AnsRed!4,colframe=AnsRed!45,boxrule=0.6pt,arc=3pt,
  left=8pt,right=8pt,top=4pt,bottom=4pt,
  before upper={\textbf{\color{AnsRed!85!black}\faExclamationTriangle\ \WarnWord}\ }}

% "Note for parents" box (short hint on how to explain the idea at home)
\newtcolorbox{parentbox}{enhanced,breakable,colback=tmTeal!5,colframe=tmTeal!50,boxrule=0.6pt,arc=3pt,
  left=8pt,right=8pt,top=4pt,bottom=4pt,
  before upper={\textbf{\color{tmTeal!80!black}\faUsers\ \ParentWord}\ }}

% numbered step inside a solution:  \step{1}{Short title}  -> "Bước 1. Short title"
\newcommand{\step}[2]{\par\smallskip\noindent{\bfseries\color{tmBlue!85!black}\StepWord~#1.}\ {\bfseries #2}\par\nopagebreak}
% verification line:  \verify{...}
\newcommand{\verify}[1]{\par\smallskip\noindent{\bfseries\color{tmGreen!80!black}\faCheckCircle\ \CheckWord}\ #1\par}

% multiple-choice options
% \choices{..}{..}{..}{..}{..}      : five short options on one line
% \choicesII{..}{..}{..}{..}{..}    : two/three per line (medium length)
% \choicesL{..}{..}{..}{..}{..}     : one option per line (long text)
\newcommand{\choices}[5]{\par\smallskip\noindent
  \begin{tabularx}{\linewidth}{@{}XXXXX@{}}(A)\ #1&(B)\ #2&(C)\ #3&(D)\ #4&(E)\ #5\end{tabularx}\par}
\newcommand{\choicesII}[5]{\par\smallskip\noindent
  \begin{tabularx}{\linewidth}{@{}XXX@{}}(A)\ #1&(B)\ #2&(C)\ #3\\(D)\ #4&(E)\ #5&\end{tabularx}\par}
\newcommand{\choicesL}[5]{\par\smallskip\noindent
  \begin{tabular}{@{}l@{}}(A)\ #1\\(B)\ #2\\(C)\ #3\\(D)\ #4\\(E)\ #5\end{tabular}\par}

% centred figure inside a problem / solution
\newenvironment{figC}{\par\vspace{2pt}\begin{center}\begin{tikzpicture}[font=\sffamily\small,>=stealth,line join=round]}%
  {\end{tikzpicture}\end{center}\vspace{1pt}}

\def\DocGrade{Lớp 7}
\def\ProbWord{Bài}
\def\SolWord{Lời giải.}
\def\GuideWord{Hướng dẫn giải.}
\def\PageWord{Trang}
\def\NameWord{Họ và tên}
\def\DateWord{Ngày}
\def\SolTitle{Lời giải chi tiết}
\def\WorkWord{Bài làm:}
\def\TopicWord{Chủ đề}
\def\AnsWord{Đáp số}
\def\TipWord{Mẹo ThinkAlgo.}
\def\RecallWord{Nhắc lại kiến thức.}
\def\ProofWord{Vì sao lại đúng?}
\def\WarnWord{Lỗi thường gặp.}
\def\ParentWord{Gợi ý cho phụ huynh.}
\def\StepWord{Bước}
\def\CheckWord{Thử lại:}

\def\DocTopic{Hình học và Đo lường}
\begin{document}
\titlehead
\chapterhead{2}{Hình học và Đo lường}

% ---------------- part a ----------------
% ======================================================================
% Chủ đề 2 -- Hình học và Đo lường -- phần a (lí thuyết, ví dụ, bài 1--15)
% ======================================================================
% hình minh hoạ căn giữa, nhận tuỳ chọn TikZ (giống figC nhưng có [scale=..])
\ifdefined\gmFig\else
\newenvironment{gmFig}[1][]{\par\vspace{2pt}\begin{center}\begin{tikzpicture}[font=\sffamily\small,>=stealth,line join=round,#1]}%
  {\end{tikzpicture}\end{center}\vspace{1pt}}
\fi
% vẽ khối lập phương nhìn từ phía trước: #1 = dịch ngang, #2 = nhãn, #3 = lệnh tô trên mặt trước (toạ độ 0..1)
\providecommand{\gmCube}[3]{%
  \begin{scope}[shift={(#1,0)},x=1.05cm,y=1.05cm]
    \fill[white] (0,0) rectangle (1,1);
    #3
    \draw[Ink,line width=0.8pt] (0,0) rectangle (1,1);
    \draw[Ink,line width=0.8pt] (0,1)--(0.4,1.25)--(1.4,1.25)--(1.4,0.25)--(1,0);
    \draw[Ink,line width=0.8pt] (1,1)--(1.4,1.25);
    \node[font=\small\bfseries,anchor=east] at (-0.05,1.1) {#2};
  \end{scope}}

\begin{knowbox}
\textbf{Ý tưởng lớn.} Hình học là môn học về \emph{hình dạng}, còn đo lường là việc gắn cho hình dạng ấy những \emph{con số}: độ dài, số đo góc, chu vi, diện tích, thể tích. Hầu hết các bài toán trong chủ đề này đều được giải bằng một vài ý tưởng rất đơn giản nhưng dùng đi dùng lại nhiều lần:
\begin{itemize}[leftmargin=1.3em,itemsep=1pt,topsep=2pt]
  \item \textbf{Góc có ``quỹ'' cố định:} hai góc kề bù luôn có tổng $180^\circ$, ba góc của một tam giác luôn có tổng $180^\circ$, một vòng quay trọn vẹn là $360^\circ$. Biết một số góc, ta suy ra các góc còn lại.
  \item \textbf{Diện tích có thể cắt -- ghép -- bù:} một hình phức tạp được chia thành các hình quen thuộc (chữ nhật, tam giác, hình thang), hoặc được tính bằng ``hình lớn trừ đi phần thừa''.
  \item \textbf{Tỉ lệ thay cho con số:} ``phần tô màu chiếm bao nhiêu phần trăm'' chỉ cần so sánh phần tô màu với toàn bộ hình, nhiều khi không cần biết kích thước thật.
  \item \textbf{Hình dạng không đổi khi quay hay lật:} quay, đối xứng, lăn một khối lập phương chỉ làm thay đổi \emph{vị trí}, không làm thay đổi \emph{kích thước} và \emph{hình dạng}.
  \item \textbf{Khối hộp được xếp thành từng lớp:} thể tích $=$ diện tích đáy $\times$ chiều cao.
\end{itemize}
Khi gặp một bài hình, hãy luôn \textbf{vẽ hình, ghi số liệu lên hình}, rồi tự hỏi: ``Mình đã biết gì? Còn thiếu gì? Công thức nào nối hai điều đó với nhau?''.
\end{knowbox}
\vspace{0.3cm}

% ----------------------------------------------------------------------
\shead{Mặt phẳng toạ độ}
\begin{knowbox}
\textbf{Hệ trục toạ độ $Oxy$} gồm hai trục số vuông góc với nhau tại gốc $O$:
\begin{itemize}[leftmargin=1.3em,itemsep=1pt,topsep=2pt]
  \item trục nằm ngang $Ox$ gọi là \textbf{trục hoành}, chiều dương hướng sang \emph{phải};
  \item trục thẳng đứng $Oy$ gọi là \textbf{trục tung}, chiều dương hướng lên \emph{trên}.
\end{itemize}
Mỗi điểm $M$ trong mặt phẳng ứng với đúng một cặp số $(x_0;\,y_0)$, viết $M(x_0;\,y_0)$:
\begin{itemize}[leftmargin=1.3em,itemsep=1pt,topsep=2pt]
  \item $x_0$ là \textbf{hoành độ}: từ $M$ kẻ đường vuông góc xuống trục $Ox$, chân đường vuông góc ứng với số $x_0$;
  \item $y_0$ là \textbf{tung độ}: từ $M$ kẻ đường vuông góc sang trục $Oy$, chân đường vuông góc ứng với số $y_0$.
\end{itemize}
Hai trục chia mặt phẳng thành bốn \textbf{góc phần tư}, đánh số ngược chiều kim đồng hồ:
\begin{center}
\begin{tabular}{|c|c|c|c|c|}
\hline
 & Góc phần tư I & Góc phần tư II & Góc phần tư III & Góc phần tư IV \\ \hline
Dấu của $(x;\,y)$ & $(+;\,+)$ & $(-;\,+)$ & $(-;\,-)$ & $(+;\,-)$ \\ \hline
\end{tabular}
\end{center}
\textbf{Khoảng cách trên đường ngang, đường dọc.} Hai điểm cùng tung độ (cùng nằm trên một đường nằm ngang) thì cách nhau một khoảng bằng hiệu hai hoành độ (lấy số lớn trừ số nhỏ). Tương tự, hai điểm cùng hoành độ cách nhau một khoảng bằng hiệu hai tung độ.
\end{knowbox}

\begin{gmFig}[scale=0.55]
  \fill[tmBlue!7] (0,0) rectangle (4.6,4.6);
  \fill[tmOrange!8] (0,0) rectangle (-4.6,4.6);
  \fill[tmGreen!8] (0,0) rectangle (-4.6,-4.6);
  \fill[tmPurple!7] (0,0) rectangle (4.6,-4.6);
  \draw[Line,line width=0.3pt] (-4.5,-4.5) grid (4.5,4.5);
  \draw[->,Navy,thick] (-4.8,0)--(5,0) node[right,text=Ink]{$x$};
  \draw[->,Navy,thick] (0,-4.8)--(0,5) node[above,text=Ink]{$y$};
  \foreach \t in {-4,-3,-2,-1,1,2,3,4}{
     \draw[Navy] (\t,0.12)--(\t,-0.12); \draw[Navy] (0.12,\t)--(-0.12,\t);}
  \foreach \t in {-4,-2,2,4}{
     \node[below,font=\scriptsize] at (\t,-0.12){$\t$}; \node[left,font=\scriptsize] at (-0.12,\t){$\t$};}
  \node[below left,font=\scriptsize] at (0,0){$O$};
  \node[font=\bfseries,text=tmBlue,fill=white,inner sep=1pt] at (3.2,3.9){I $(+;+)$};
  \node[font=\bfseries,text=tmOrange!85!black,fill=white,inner sep=1pt] at (-3.2,3.9){II $(-;+)$};
  \node[font=\bfseries,text=tmGreen!70!black,fill=white,inner sep=1pt] at (-3.2,-3.9){III $(-;-)$};
  \node[font=\bfseries,text=tmPurple,fill=white,inner sep=1pt] at (3.2,-3.9){IV $(+;-)$};
  \draw[dashed,AnsRed] (3,0)--(3,2)--(0,2);
  \fill[AnsRed] (3,2) circle (3.5pt) node[above right,text=AnsRed,fill=white,inner sep=1pt]{$M(3;\,2)$};
  \draw[dashed,tmPurple] (-2,0)--(-2,-3)--(0,-3);
  \fill[tmPurple] (-2,-3) circle (3.5pt) node[below right,text=tmPurple,fill=white,inner sep=1pt]{$N(-2;\,-3)$};
\end{gmFig}

\begin{tipbox}
Đọc toạ độ theo thứ tự ``\textbf{đi ngang trước, đi dọc sau}'': bắt đầu từ $O$, đi sang phải (số dương) hoặc sang trái (số âm) để được hoành độ; rồi đi lên (số dương) hoặc đi xuống (số âm) để được tung độ. Chẳng hạn, để đến $N(-2;\,-3)$: sang trái $2$ ô, rồi xuống $3$ ô.
\end{tipbox}

\begin{warnbox}
Đảo thứ tự hai số là lỗi rất hay gặp: $(3;\,-1)$ và $(-1;\,3)$ là hai điểm \emph{khác nhau} (một điểm ở góc phần tư IV, một điểm ở góc phần tư II). Hoành độ \emph{luôn} viết trước.
\end{warnbox}

\textbf{Ví dụ nhỏ.} Cho $A(1;\,1)$, $B(5;\,1)$, $C(5;\,4)$. Vì $A$ và $B$ cùng tung độ $1$ nên $AB=5-1=4$. Vì $B$ và $C$ cùng hoành độ $5$ nên $BC=4-1=3$. Góc $B$ vuông (một cạnh nằm ngang, một cạnh thẳng đứng), nên diện tích tam giác $ABC$ là $\dfrac{4\cdot 3}{2}=6$ (đơn vị diện tích).

\begin{parentbox}
Có thể biến mặt phẳng toạ độ thành trò chơi ``tìm kho báu'' trên giấy kẻ ô: bố mẹ đọc ``sang trái 2, lên 3'', con đánh dấu điểm và viết toạ độ $(-2;\,3)$. Đổi vai cho con đọc, bố mẹ vẽ. Chỉ sau vài lượt, thứ tự ``hoành trước, tung sau'' sẽ thành phản xạ.
\end{parentbox}

% ----------------------------------------------------------------------
\shead{Góc: kề bù, đối đỉnh, tổng ba góc của tam giác}
\begin{knowbox}
\begin{itemize}[leftmargin=1.3em,itemsep=2pt,topsep=2pt]
  \item \textbf{Góc vuông} bằng $90^\circ$, \textbf{góc bẹt} (nửa vòng) bằng $180^\circ$, \textbf{một vòng trọn} bằng $360^\circ$.
  \item \textbf{Hai góc kề bù} (có chung một cạnh, hai cạnh còn lại là hai tia đối nhau) có tổng bằng $180^\circ$.
  \item \textbf{Hai góc đối đỉnh} (mỗi cạnh của góc này là tia đối của một cạnh của góc kia) thì \textbf{bằng nhau}.
  \item \textbf{Tổng ba góc của một tam giác bằng $180^\circ$.} Do đó: góc thứ ba $=180^\circ-$ (tổng hai góc đã biết).
  \item \textbf{Góc ngoài của tam giác} (góc kề bù với một góc trong) bằng \textbf{tổng hai góc trong không kề với nó}.
  \item Trong tam giác vuông, hai góc nhọn có tổng bằng $90^\circ$.
\end{itemize}
\end{knowbox}

\begin{gmFig}[scale=0.9]
  % hai đường thẳng cắt nhau
  \coordinate (O) at (0,0);
  \draw[tmBlue,thick] (-2,-0.9)--(2,0.9);
  \draw[tmOrange,thick] (-1.6,1.2)--(1.6,-1.2);
  \coordinate (A) at (2,0.9); \coordinate (B) at (-1.6,1.2);
  \coordinate (C) at (-2,-0.9); \coordinate (D) at (1.6,-1.2);
  \pic[draw,AnsRed,thick,angle radius=0.45cm,"$1$",angle eccentricity=1.6]{angle=A--O--B};
  \pic[draw,tmGreen,thick,angle radius=0.45cm,"$2$",angle eccentricity=1.6]{angle=B--O--C};
  \pic[draw,AnsRed,thick,angle radius=0.45cm,"$3$",angle eccentricity=1.6]{angle=C--O--D};
  \node[font=\footnotesize,align=center] at (0,-1.8){$\widehat{1}+\widehat{2}=180^\circ$ (kề bù)\\ $\widehat{1}=\widehat{3}$ (đối đỉnh)};
  % tam giác với đường song song
  \begin{scope}[shift={(4.3,-1)}]
    \coordinate (P) at (0,0); \coordinate (Q) at (3.4,0); \coordinate (R) at (1.2,2.2);
    \draw[Ink,thick] (P)--(Q)--(R)--cycle;
    \draw[tmPurple,thick,dashed] (-0.6,2.2)--(3.2,2.2);
    \coordinate (L) at (-0.6,2.2); \coordinate (M) at (3.2,2.2);
    \pic[draw,AnsRed,thick,angle radius=0.45cm,"$a$",angle eccentricity=1.55]{angle=Q--P--R};
    \pic[draw,tmGreen!70!black,thick,angle radius=0.45cm,"$c$",angle eccentricity=1.55]{angle=R--Q--P};
    \pic[draw,tmBlue,thick,angle radius=0.4cm,"$b$",angle eccentricity=1.6]{angle=P--R--Q};
    \pic[draw,AnsRed,thick,angle radius=0.45cm,"$a$",angle eccentricity=1.55]{angle=L--R--P};
    \pic[draw,tmGreen!70!black,thick,angle radius=0.45cm,"$c$",angle eccentricity=1.55]{angle=Q--R--M};
    \node[font=\footnotesize] at (1.6,-0.6){$a+b+c=180^\circ$};
  \end{scope}
\end{gmFig}

\begin{proofbox}
\textbf{Hai góc đối đỉnh bằng nhau.} Ở hình bên trái, $\widehat{1}+\widehat{2}=180^\circ$ (kề bù) và $\widehat{3}+\widehat{2}=180^\circ$ (cũng kề bù). Hai góc $\widehat{1}$ và $\widehat{3}$ cùng ``bù'' với một góc $\widehat{2}$, nên $\widehat{1}=180^\circ-\widehat{2}=\widehat{3}$.

\textbf{Tổng ba góc tam giác bằng $180^\circ$.} Qua đỉnh trên của tam giác, kẻ đường thẳng song song với cạnh đáy (nét đứt). Theo tính chất hai đường thẳng song song, các cặp góc \emph{so le trong} bằng nhau: góc bên trái ở đỉnh trên bằng $a$, góc bên phải bằng $c$. Ba góc $a$, $b$, $c$ ở đỉnh trên ghép lại thành một góc bẹt, nên $a+b+c=180^\circ$.

\textbf{Góc ngoài.} Góc ngoài tại một đỉnh kề bù với góc trong tại đỉnh đó: góc ngoài $=180^\circ-$ góc trong. Mà $180^\circ-$ góc trong lại đúng bằng tổng hai góc trong còn lại. Vậy góc ngoài bằng tổng hai góc trong không kề với nó.
\end{proofbox}

\begin{tipbox}
Trong bài ``tìm góc'', hãy viết số đo của mọi góc tìm được \textbf{lên hình} ngay lập tức. Một góc vừa tìm ra thường là chìa khoá để tìm góc tiếp theo (kề bù, đối đỉnh, tổng ba góc,...).
\end{tipbox}

\begin{warnbox}
Đừng nhầm \emph{đối đỉnh} (bằng nhau) với \emph{kề bù} (tổng $180^\circ$). Hai góc đối đỉnh nằm ``đối diện'' nhau qua giao điểm, không có cạnh chung; hai góc kề bù nằm ``cạnh nhau'', có một cạnh chung.
\end{warnbox}

% ----------------------------------------------------------------------
\shead{Tam giác cân, tam giác đều và bất đẳng thức tam giác}
\begin{knowbox}
\begin{itemize}[leftmargin=1.3em,itemsep=2pt,topsep=2pt]
  \item \textbf{Tam giác cân} là tam giác có hai cạnh bằng nhau (gọi là hai \emph{cạnh bên}); cạnh còn lại là \emph{cạnh đáy}, góc xen giữa hai cạnh bên là \emph{góc ở đỉnh}.
  \item Trong tam giác cân, \textbf{hai góc ở đáy bằng nhau}. Do đó, nếu góc ở đỉnh là $d$ và mỗi góc ở đáy là $g$ thì
  \[ d+2g=180^\circ,\qquad d=180^\circ-2g,\qquad g=\frac{180^\circ-d}{2}. \]
  \item Góc ở đáy luôn là \textbf{góc nhọn} (nhỏ hơn $90^\circ$), vì hai góc ở đáy bằng nhau mà tổng của chúng phải nhỏ hơn $180^\circ$.
  \item \textbf{Tam giác đều} có ba cạnh bằng nhau, ba góc bằng nhau, mỗi góc bằng $180^\circ:3=60^\circ$.
  \item \textbf{Bất đẳng thức tam giác:} trong một tam giác, mỗi cạnh luôn \emph{nhỏ hơn} tổng hai cạnh còn lại. Để kiểm tra ba độ dài có tạo thành tam giác không, chỉ cần kiểm tra: \emph{cạnh dài nhất $<$ tổng hai cạnh kia}.
\end{itemize}
\end{knowbox}

\begin{gmFig}[scale=0.85]
  \coordinate (A) at (1.5,2.6); \coordinate (B) at (0,0); \coordinate (C) at (3,0);
  \fill[tmBlue!6] (A)--(B)--(C)--cycle;
  \draw[Ink,thick] (A)--(B)--(C)--cycle;
  \draw[tmBlue,thick] ($(A)!0.5!(B)+(-0.1,0.06)$)--++(0.2,-0.12);
  \draw[tmBlue,thick] ($(A)!0.5!(C)+(-0.1,-0.06)$)--++(0.2,0.12);
  \pic[draw,tmOrange,thick,angle radius=0.5cm,"$g$",angle eccentricity=1.5]{angle=C--B--A};
  \pic[draw,tmOrange,thick,angle radius=0.5cm,"$g$",angle eccentricity=1.5]{angle=A--C--B};
  \pic[draw,tmGreen!70!black,thick,angle radius=0.5cm,"$d$",angle eccentricity=1.5]{angle=B--A--C};
  \node[above,font=\footnotesize] at (A){đỉnh};
  \node[below,font=\footnotesize] at (1.5,0){cạnh đáy};
  \begin{scope}[shift={(5,0)}]
    \draw[tmGreen!70!black,line width=2pt] (0,0)--(4,0);
    \draw[tmBlue,line width=2pt] (0,0)--(1.6,1.6);
    \draw[tmOrange,line width=2pt] (1.6,1.6)--(4,0);
    \node[below,font=\footnotesize] at (2,0){$8$};
    \node[left,font=\footnotesize] at (0.8,0.9){$4$};
    \node[right,font=\footnotesize] at (2.9,0.9){$5$};
    \node[font=\footnotesize,align=center] at (2,-0.85){$8<4+5$: tạo được tam giác};
  \end{scope}
\end{gmFig}

\begin{proofbox}
\textbf{Vì sao hai góc ở đáy bằng nhau?} Gọi $M$ là trung điểm của cạnh đáy $BC$ trong tam giác cân $ABC$ ($AB=AC$). Hai tam giác $ABM$ và $ACM$ có $AB=AC$, $BM=CM$ và cạnh chung $AM$, nên chúng bằng nhau (trường hợp cạnh -- cạnh -- cạnh). Suy ra hai góc tương ứng $\widehat{B}$ và $\widehat{C}$ bằng nhau.

\textbf{Vì sao có bất đẳng thức tam giác?} Đường ngắn nhất nối hai điểm là đoạn thẳng. Đi từ $B$ đến $C$ ``vòng qua'' $A$ (tức là đi $BA+AC$) luôn dài hơn đi thẳng $BC$. Nếu cạnh dài nhất bằng hoặc lớn hơn tổng hai cạnh kia thì hai cạnh ngắn ``không với tới nhau'': chúng chỉ nằm bẹp trên cạnh dài hoặc không chạm được nhau, không tạo thành tam giác.
\end{proofbox}

\begin{warnbox}
Khi đề chỉ nói ``tam giác cân có \emph{một góc} bằng $50^\circ$'', ta \textbf{chưa biết} góc đó là góc ở đỉnh hay góc ở đáy. Phải xét \textbf{cả hai trường hợp}:
\begin{itemize}[leftmargin=1.3em,itemsep=1pt,topsep=2pt]
  \item $50^\circ$ là góc ở đỉnh: hai góc đáy là $(180^\circ-50^\circ):2=65^\circ$ và $65^\circ$;
  \item $50^\circ$ là góc ở đáy: góc đáy kia cũng là $50^\circ$, góc ở đỉnh là $180^\circ-2\cdot 50^\circ=80^\circ$.
\end{itemize}
Riêng khi góc đã cho \emph{lớn hơn hoặc bằng} $90^\circ$ thì nó chỉ có thể là góc ở đỉnh.
\end{warnbox}

% ----------------------------------------------------------------------
\shead{Phép quay, phép đối xứng và khối lập phương lăn}
\begin{knowbox}
\textbf{Phép quay} quanh một điểm $Z$ (gọi là \emph{tâm quay}) một góc $\alpha$ theo một chiều nhất định (cùng chiều hoặc ngược chiều kim đồng hồ) đưa mỗi điểm $M$ đến điểm $M'$ sao cho $ZM'=ZM$ và góc $\widehat{MZM'}=\alpha$.
\begin{itemize}[leftmargin=1.3em,itemsep=2pt,topsep=2pt]
  \item Mọi điểm của hình đều quay \textbf{cùng một góc}, cùng một chiều; khoảng cách đến tâm quay không đổi.
  \item Phép quay giữ nguyên hình dạng và kích thước của hình; hình chỉ ``xoay'' đi chứ không bị ``lật mặt''.
  \item Quay $\alpha$ theo chiều kim đồng hồ cho cùng kết quả với quay $360^\circ-\alpha$ ngược chiều kim đồng hồ. Quay $360^\circ$ thì hình trở về chỗ cũ.
  \item \textbf{Cách đo góc quay:} chọn một đoạn thẳng xuất phát từ tâm quay $Z$ trên hình ban đầu, tìm đoạn tương ứng trên hình mới, rồi đo góc mà đoạn đó đã quét qua \emph{theo đúng chiều quay}.
\end{itemize}
\textbf{Phép đối xứng trục} (lật qua một đường thẳng, giống soi gương) cũng giữ nguyên kích thước, nhưng \textbf{làm ngược chiều} hình: chữ \textsf{P} soi gương sẽ thành chữ \textsf{P} quay mặt ngược lại.

\textbf{Trong mặt phẳng toạ độ} (tâm quay là gốc $O$):
\begin{center}
\begin{tabular}{|l|c|}
\hline
Quay $90^\circ$ ngược chiều kim đồng hồ & $(x;\,y)\ \to\ (-y;\,x)$ \\ \hline
Quay $180^\circ$ & $(x;\,y)\ \to\ (-x;\,-y)$ \\ \hline
Quay $90^\circ$ cùng chiều kim đồng hồ & $(x;\,y)\ \to\ (y;\,-x)$ \\ \hline
Đối xứng qua trục $Ox$ \;/\; qua trục $Oy$ & $(x;\,y)\to(x;\,-y)$ \;/\; $(x;\,y)\to(-x;\,y)$ \\ \hline
\end{tabular}
\end{center}
\end{knowbox}

\begin{gmFig}[scale=1.05]
  \coordinate (Z) at (0,0);
  \fill[tmBlue!35,draw=tmBlue!80!black,thick] (0,0)--(2.2,0)--(2.2,0.9)--(0.9,0.9)--cycle;
  \fill[white,draw=tmOrange!90!black,thick] (0,0)--(0,2.2)--(-0.9,2.2)--(-0.9,0.9)--cycle;
  \fill (Z) circle (2.2pt) node[below left]{$Z$};
  \draw[->,AnsRed,thick] (1.4,-0.15) arc[start angle=-6,end angle=-264,radius=1.4];
  \node[AnsRed,font=\footnotesize] at (0,-1.75){chiều kim đồng hồ: $270^\circ$};
  \draw[->,tmGreen!70!black,thick] (2.4,0.25) arc[start angle=6,end angle=84,radius=2.4];
  \node[tmGreen!60!black,font=\footnotesize,align=left] at (3.5,2.1){ngược chiều\\kim đồng hồ: $90^\circ$};
  \node[font=\footnotesize] at (1.4,0.4){ban đầu};
  \node[font=\footnotesize,rotate=90] at (-0.45,1.6){ảnh};
\end{gmFig}

\begin{proofbox}
\textbf{Vì sao quay $270^\circ$ theo chiều kim đồng hồ lại giống quay $90^\circ$ ngược chiều?} Một vòng trọn là $360^\circ$. Đi $270^\circ$ theo một chiều thì chỉ còn thiếu $360^\circ-270^\circ=90^\circ$ nữa là về chỗ cũ; nghĩa là điểm cuối cũng chính là điểm ta đến được nếu đi $90^\circ$ theo chiều ngược lại. Trên hình: cạnh nằm ngang (hướng sang phải) muốn chuyển thành cạnh thẳng đứng (hướng lên), ta có thể quay ngược chiều kim đồng hồ $90^\circ$, hoặc quay theo chiều kim đồng hồ $270^\circ$ (qua hướng xuống, hướng sang trái, rồi mới lên trên).
\end{proofbox}

\begin{knowbox}
\textbf{Khối lập phương lăn.} Khi lăn một khối lập phương, mỗi mặt được mang theo nguyên vẹn. Vì vậy, một mặt đã sơn họa tiết:
\begin{itemize}[leftmargin=1.3em,itemsep=1pt,topsep=2pt]
  \item có thể hiện ra ở vị trí khác (mặt trước, mặt trên, mặt bên,...);
  \item có thể bị \emph{xoay} đi $90^\circ$, $180^\circ$ hoặc $270^\circ$ trong mặt phẳng của nó;
  \item \textbf{nhưng không bao giờ} bị lật thành ảnh gương, thay đổi kích thước, hay ``thêm'' phần sơn mới.
\end{itemize}
\end{knowbox}

\begin{tipbox}
Với bài khối lập phương, hãy ``đánh dấu'' một chi tiết đặc biệt của họa tiết (ví dụ: góc vuông của hình tam giác sơn đen nằm ở \emph{một đỉnh} của mặt) rồi kiểm tra chi tiết ấy ở từng phương án. Cắt một mẩu giấy hình vuông, tô họa tiết lên và xoay thử cũng là một cách kiểm tra rất hiệu quả.
\end{tipbox}

% ----------------------------------------------------------------------
\shead{Đơn vị đo và trục số}
\begin{knowbox}
\textbf{Đơn vị đo độ dài.} Mỗi đơn vị gấp $10$ lần đơn vị liền sau nó:
\[ 1\text{ km}=10\text{ hm},\quad 1\text{ hm}=10\text{ dam},\quad 1\text{ dam}=10\text{ m}, \]
\[ 1\text{ m}=10\text{ dm},\quad 1\text{ dm}=10\text{ cm},\quad 1\text{ cm}=10\text{ mm}. \]
Các mối liên hệ hay dùng nhất: $1\text{ m}=100\text{ cm}=1000\text{ mm}$; $\;1\text{ km}=1000\text{ m}$.

\textbf{Đơn vị đo diện tích} gấp nhau $100$ lần: $1\text{ m}^2=100\text{ dm}^2=10\,000\text{ cm}^2$.

\textbf{Đơn vị đo thể tích} gấp nhau $1000$ lần: $1\text{ m}^3=1000\text{ dm}^3$, $\;1\text{ dm}^3=1000\text{ cm}^3=1$ lít.

\textbf{Quy tắc đổi:} đổi từ đơn vị \emph{lớn} sang đơn vị \emph{nhỏ} thì \textbf{nhân}, từ đơn vị \emph{nhỏ} sang đơn vị \emph{lớn} thì \textbf{chia}.

\textbf{Trục số có vạch chia đều.} Nếu biết số ứng với hai vạch và đếm được giữa chúng có $k$ \emph{khoảng} (không phải $k$ vạch), thì
\[ \text{giá trị một khoảng}=\frac{\text{hiệu hai số}}{k}. \]
\end{knowbox}

\begin{proofbox}
\textbf{Vì sao $1\text{ m}^2=10\,000\text{ cm}^2$ chứ không phải $100\text{ cm}^2$?} Một mét vuông là diện tích hình vuông cạnh $1$ m, tức cạnh $100$ cm. Hình vuông đó chia được thành $100$ hàng, mỗi hàng $100$ ô vuông cạnh $1$ cm, nên có $100\times 100=10\,000$ ô vuông $1\text{ cm}^2$. Tương tự, khối lập phương cạnh $1$ dm $=10$ cm gồm $10\times10\times10=1000$ khối lập phương nhỏ $1\text{ cm}^3$.
\end{proofbox}

\begin{gmFig}[x=1.15cm]
  \draw[->,thick,Navy] (-0.4,0)--(10.6,0);
  \foreach \i in {0,...,10}{\draw[Navy,line width=0.8pt] (\i,0.13)--(\i,-0.13);}
  \node[below=3pt] at (0,0){$0$};
  \node[below=3pt] at (8,0){$2$};
  \foreach \i in {0,...,7}{\draw[<->,tmOrange,thin] (\i+0.06,0.32)--(\i+0.94,0.32);}
  \node[above=6pt,font=\footnotesize,text=tmOrange!85!black] at (4,0.32){$8$ khoảng bằng nhau từ $0$ đến $2$ $\Rightarrow$ mỗi khoảng $=2:8=0{,}25$};
  \fill[AnsRed] (3,0) circle (2.6pt) node[below=4pt,text=AnsRed]{$A$};
  \fill[AnsRed] (10,0) circle (2.6pt) node[below=4pt,text=AnsRed]{$B$};
\end{gmFig}
\textbf{Ví dụ nhỏ.} Trên trục số trên, $A$ cách $0$ ba khoảng nên $A=3\cdot 0{,}25=0{,}75$; $B$ cách $0$ mười khoảng nên $B=10\cdot0{,}25=2{,}5$.

\begin{warnbox}
Đếm \emph{vạch} thay vì đếm \emph{khoảng}: từ vạch $0$ đến vạch $2$ ở hình trên có $9$ vạch nhưng chỉ có $8$ khoảng. Giữa hai số luôn có số khoảng ít hơn số vạch $1$ đơn vị (giống như hàng cây: $9$ cây thì có $8$ khoảng cách).
\end{warnbox}

% ----------------------------------------------------------------------
\shead{Chu vi và diện tích các hình phẳng}
\begin{knowbox}
\begin{center}
\renewcommand{\arraystretch}{1.5}
\begin{tabular}{|l|c|c|}
\hline
\textbf{Hình} & \textbf{Chu vi} & \textbf{Diện tích} \\ \hline
Hình chữ nhật (dài $a$, rộng $b$) & $P=2\cdot(a+b)$ & $S=a\cdot b$ \\ \hline
Hình vuông (cạnh $a$) & $P=4\cdot a$ & $S=a\cdot a=a^2$ \\ \hline
Tam giác (đáy $a$, chiều cao $h$ ứng với đáy đó) & tổng ba cạnh & $S=\dfrac{a\cdot h}{2}$ \\ \hline
Hình bình hành (đáy $a$, chiều cao $h$) & $2\cdot$(tổng hai cạnh kề) & $S=a\cdot h$ \\ \hline
Hình thang (hai đáy $a$, $b$; chiều cao $h$) & tổng bốn cạnh & $S=\dfrac{(a+b)\cdot h}{2}$ \\ \hline
\end{tabular}
\end{center}
\textbf{Chiều cao} của tam giác (hình bình hành, hình thang) là đoạn thẳng \emph{vuông góc} kẻ từ đỉnh (hoặc từ cạnh đối diện) xuống đường thẳng chứa đáy.

\textbf{Tính ngược:} từ $S=\dfrac{a\cdot h}{2}$ suy ra $h=\dfrac{2\cdot S}{a}$ và $a=\dfrac{2\cdot S}{h}$; từ chu vi hình vuông suy ra cạnh $=P:4$.
\end{knowbox}

\begin{gmFig}[scale=0.8]
  % tam giác trong hình chữ nhật
  \fill[tmBlue!25] (0,0)--(4,0)--(1.4,2.4)--cycle;
  \draw[Slate,dashed] (0,0) rectangle (4,2.4);
  \draw[Slate,dashed] (1.4,0)--(1.4,2.4);
  \draw[Ink,thick] (0,0)--(4,0)--(1.4,2.4)--cycle;
  \draw[Ink] (1.4,0) rectangle (1.62,0.22);
  \node[font=\footnotesize] at (0.55,1.55){$1$};
  \node[font=\footnotesize] at (3.25,1.75){$2$};
  \node[font=\footnotesize] at (0.85,0.55){$1'$};
  \node[font=\footnotesize] at (2.25,0.55){$2'$};
  \node[left,font=\footnotesize] at (1.4,1.2){$h$};
  \node[below,font=\footnotesize] at (2,0){$a$};
  % hình thang ghép đôi
  \begin{scope}[shift={(5.2,0)}]
    \fill[tmGreen!25,draw=Ink,thick] (0,0)--(3,0)--(2.2,2)--(0.8,2)--cycle;
    \fill[tmOrange!20,draw=Ink,thick,dashed] (3,0)--(4.4,0)--(5.2,2)--(2.2,2)--cycle;
    \node[below,font=\footnotesize] at (1.5,0){$a$};
    \node[above,font=\footnotesize] at (1.5,2){$b$};
    \node[below,font=\footnotesize] at (3.7,0){$b$};
    \node[above,font=\footnotesize] at (3.7,2){$a$};
    \draw[Slate] (0.8,2)--(0.8,0); \draw[Slate] (0.8,0) rectangle (1.0,0.2);
    \node[right,font=\footnotesize] at (0.8,1.1){$h$};
    \node[font=\footnotesize,align=center] at (2.6,-0.85){hai hình thang ghép thành\\hình bình hành đáy $a+b$};
  \end{scope}
\end{gmFig}

\begin{proofbox}
\textbf{Diện tích tam giác.} Đặt tam giác vào trong hình chữ nhật có cùng đáy $a$ và cùng chiều cao $h$ (hình trái). Đường cao chia hình chữ nhật thành hai hình chữ nhật nhỏ; mỗi hình chữ nhật nhỏ bị đường chéo chia thành hai tam giác bằng nhau: phần trong tam giác ($1'$, $2'$) bằng phần ngoài tam giác ($1$, $2$). Vậy tam giác chiếm đúng một nửa hình chữ nhật: $S=\dfrac{a\cdot h}{2}$.

\textbf{Diện tích hình bình hành.} Cắt một tam giác vuông ở một đầu, dời sang đầu kia, ta được hình chữ nhật có cùng đáy $a$, cùng chiều cao $h$: $S=a\cdot h$.

\textbf{Diện tích hình thang.} Hai hình thang giống hệt nhau, một hình quay ngược lại, ghép thành một hình bình hành có đáy $a+b$ và chiều cao $h$ (hình phải). Một hình thang bằng một nửa: $S=\dfrac{(a+b)\cdot h}{2}$.
\end{proofbox}

\begin{tipbox}
Hai tam giác có \textbf{cùng đáy} (hoặc hai đáy bằng nhau) và \textbf{cùng chiều cao} thì có diện tích bằng nhau, dù hình dạng trông rất khác nhau. Hệ quả: đường nối một đỉnh với trung điểm cạnh đối diện chia tam giác thành hai phần có diện tích bằng nhau.
\end{tipbox}

\begin{warnbox}
\begin{itemize}[leftmargin=1.3em,itemsep=1pt,topsep=0pt]
  \item Chiều cao phải \textbf{vuông góc} với đáy; cạnh bên xiên \emph{không phải} là chiều cao.
  \item Quên chia $2$ trong công thức diện tích tam giác, hình thang.
  \item Nhầm đơn vị: chu vi đo bằng cm, diện tích đo bằng $\text{cm}^2$. Một hình vuông chu vi $36$ cm có cạnh $9$ cm và diện tích $81\text{ cm}^2$ -- không phải $36\text{ cm}^2$.
\end{itemize}
\end{warnbox}

% ----------------------------------------------------------------------
\shead{Phần diện tích: phân số, phần trăm, hình tròn và cắt ghép}
\begin{knowbox}
\begin{itemize}[leftmargin=1.3em,itemsep=2pt,topsep=2pt]
  \item Nếu một hình được chia thành $n$ phần \textbf{bằng nhau} và tô màu $m$ phần thì phần tô màu chiếm $\dfrac{m}{n}$ hình, tức là $\dfrac{m}{n}\cdot 100\%$.
  \item \textbf{Hình tròn và hình quạt.} Một vòng quanh tâm là $360^\circ$. Hình quạt có góc ở tâm $n^\circ$ chiếm $\dfrac{n}{360}$ diện tích hình tròn. Ví dụ: nửa hình tròn ứng với $180^\circ$, một phần tư hình tròn ứng với $90^\circ$ (một góc vuông ở tâm).
  \item \textbf{Vòng quay may mắn:} kim quay dừng ngẫu nhiên, mọi hướng như nhau, nên khả năng dừng ở một ô bằng đúng phần diện tích (tức phần góc) của ô đó.
  \item \textbf{Cắt -- ghép:} chia hình thành các mảnh quen thuộc rồi cộng diện tích. \textbf{Phần bù:} diện tích phần cần tìm $=$ diện tích hình lớn $-$ diện tích các phần thừa.
  \item \textbf{Đếm hình có hệ thống:} khi phải đếm bao nhiêu hình (tam giác, hình bình hành...) trong một hình vẽ, ta chia theo \emph{kích thước} (gồm $1$ mảnh, $2$ mảnh, $3$ mảnh,...) và đếm từng loại để không sót, không trùng.
\end{itemize}
\end{knowbox}

\begin{gmFig}[scale=0.9]
  \fill[tmBlue!30] (0,0)--(0:1.6) arc[start angle=0,end angle=90,radius=1.6]--cycle;
  \fill[tmOrange!30] (0,0)--(90:1.6) arc[start angle=90,end angle=135,radius=1.6]--cycle;
  \draw[Ink,thick] (0,0) circle (1.6);
  \draw[Ink,thick] (0,0)--(0:1.6); \draw[Ink,thick] (0,0)--(90:1.6);
  \draw[Ink,thick] (0,0)--(135:1.6); \draw[Ink,thick] (0,0)--(270:1.6);
  \draw[Ink] (0,0) rectangle (0.2,0.2);
  \node[font=\footnotesize] at (45:0.95){$90^\circ$};
  \node[font=\footnotesize] at (112:1.05){$45^\circ$};
  \node[font=\footnotesize] at (202:0.85){$135^\circ$};
  \node[font=\footnotesize] at (315:0.9){$90^\circ$};
  \node[font=\footnotesize,align=left,anchor=west] at (2.1,0.4){xanh: $\frac{90}{360}=\frac14=25\%$};
  \node[font=\footnotesize,align=left,anchor=west] at (2.1,-0.3){cam: $\frac{45}{360}=\frac18=12{,}5\%$};
  \begin{scope}[shift={(6.6,-1.2)},scale=0.6]
    \foreach \r/\c in {0/0,0/1,1/1,2/0,2/2,3/1,3/3,1/3}{\fill[tmGreen!45] (\c,\r) rectangle ++(1,1);}
    \draw[Ink] (0,0) grid (4,4);
    \node[font=\footnotesize,align=center] at (2,-0.9){$8$ ô trong $16$ ô: $\frac{8}{16}=50\%$};
  \end{scope}
\end{gmFig}

\begin{proofbox}
\textbf{Vì sao hình quạt góc $n^\circ$ chiếm $\frac{n}{360}$ hình tròn?} Hãy tưởng tượng chia hình tròn thành $360$ hình quạt nhỏ, mỗi hình quạt có góc ở tâm $1^\circ$. Các hình quạt nhỏ này giống hệt nhau (quay hình tròn quanh tâm $1^\circ$ thì hình quạt này trùng khít hình quạt bên cạnh), nên có diện tích bằng nhau, mỗi cái bằng $\frac{1}{360}$ hình tròn. Hình quạt góc $n^\circ$ gồm $n$ hình quạt nhỏ như thế.
\end{proofbox}

\begin{warnbox}
Chỉ được đếm ``$m$ phần trong $n$ phần'' khi các phần \textbf{bằng nhau}. Nếu các mảnh có kích thước khác nhau, phải quy về cùng một loại mảnh (ví dụ chia nhỏ thêm cho đều) hoặc tính diện tích từng mảnh.
\end{warnbox}

\begin{parentbox}
Cắt một chiếc bánh (hoặc tờ giấy tròn) thành $4$, rồi thành $8$ phần bằng nhau để con thấy: $1$ phần tư $=90^\circ$, $1$ phần tám $=45^\circ$, và nửa của một phần tư là một phần tám. Ý tưởng ``chia đôi góc thì chia đôi diện tích'' sẽ trở nên rất trực quan.
\end{parentbox}

% ----------------------------------------------------------------------
\shead{Hình khối: thể tích và diện tích bề mặt}
\begin{knowbox}
\begin{itemize}[leftmargin=1.3em,itemsep=2pt,topsep=2pt]
  \item \textbf{Hình hộp chữ nhật} có kích thước $a$, $b$, $c$:
  \[ V=a\cdot b\cdot c,\qquad S_{\text{xung quanh}}=2\cdot(a+b)\cdot c,\qquad S_{\text{toàn phần}}=2\cdot(ab+bc+ca). \]
  \item \textbf{Hình lập phương} cạnh $a$: $V=a^3$, $\;S_{\text{toàn phần}}=6\cdot a^2$.
  \item \textbf{Hình lăng trụ đứng} (đáy là tam giác, tứ giác,...): $V=S_{\text{đáy}}\cdot h$, với $h$ là chiều cao lăng trụ (độ dài cạnh bên).
  \item Số khối lập phương $1\text{ cm}\times1\text{ cm}\times1\text{ cm}$ cần để xếp kín một hình hộp đúng bằng thể tích của hình hộp tính bằng $\text{cm}^3$.
\end{itemize}
\end{knowbox}

\begin{gmFig}[scale=0.62]
  % hình hộp 4x2x3 xếp lớp
  \foreach \z/\c in {0/tmBlue,1/tmGreen,2/tmOrange}{
    \fill[\c!25,draw=\c!80!black] (0,\z) rectangle (4,\z+1);
    \fill[\c!15,draw=\c!80!black] (4,\z)--(5,\z+0.7)--(5,\z+1.7)--(4,\z+1)--cycle;
    \foreach \x in {1,2,3}{\draw[\c!80!black] (\x,\z)--(\x,\z+1);}
    \draw[\c!80!black] (4.5,\z+0.35)--(4.5,\z+1.35);
  }
  \fill[tmOrange!10,draw=tmOrange!80!black] (0,3)--(4,3)--(5,3.7)--(1,3.7)--cycle;
  \foreach \x in {1,2,3}{\draw[tmOrange!80!black] (\x,3)--(\x+1,3.7);}
  \draw[tmOrange!80!black] (0.5,3.35)--(4.5,3.35);
  \node[below,font=\footnotesize] at (2,0){$4$};
  \node[below right,font=\footnotesize] at (4.5,0.35){$2$};
  \node[right,font=\footnotesize] at (5,2.2){$3$};
  \node[font=\footnotesize,align=center] at (2.5,-1.3){mỗi lớp $4\cdot2=8$ khối, $3$ lớp: $8\cdot3=24$ khối};
  % lăng trụ tam giác
  \begin{scope}[shift={(8.5,0)}]
    \fill[tmPurple!15,draw=Ink,thick] (0,0)--(2,0)--(0,1.6)--cycle;
    \draw[Ink,thick] (0,0)--(1.6,0.9); \draw[Ink,thick] (2,0)--(3.6,0.9); \draw[Ink,thick] (0,1.6)--(1.6,2.5);
    \draw[Ink,thick,dashed] (1.6,0.9)--(3.6,0.9); \draw[Ink,thick,dashed] (1.6,0.9)--(1.6,2.5);
    \draw[Ink,thick] (3.6,0.9)--(1.6,2.5);
    \node[font=\footnotesize] at (0.55,0.45){đáy};
    \node[below right,font=\footnotesize] at (2.8,0.45){$h$};
    \node[font=\footnotesize,align=center] at (1.8,-1.3){$V=S_{\text{đáy}}\cdot h$};
  \end{scope}
\end{gmFig}

\begin{proofbox}
\textbf{Vì sao $V=a\cdot b\cdot c$?} Xếp các khối lập phương $1\text{ cm}^3$ thành từng lớp. Lớp dưới cùng là một hình chữ nhật $a$ hàng $\times$ $b$ cột, gồm $a\cdot b$ khối (đúng bằng diện tích đáy). Chồng $c$ lớp như vậy, ta có $a\cdot b\cdot c$ khối. Lập luận ``mỗi lớp $=$ diện tích đáy, số lớp $=$ chiều cao'' đúng cho mọi hình lăng trụ đứng, nên $V=S_{\text{đáy}}\cdot h$. Riêng với lăng trụ có đáy là tam giác vuông với hai cạnh góc vuông $a$, $b$: ghép hai lăng trụ như thế ta được hình hộp chữ nhật $a\times b\times h$, nên mỗi lăng trụ có thể tích $\dfrac{a\cdot b}{2}\cdot h$.
\end{proofbox}

\begin{tipbox}
Diện tích bề mặt của một hình hộp là tổng diện tích của $6$ mặt, chia thành $3$ cặp mặt đối diện bằng nhau. Khi ghép nhiều khối nhỏ thành một khối lớn, những mặt bị ``áp vào nhau'' sẽ không còn nằm ở bề mặt nữa, nên diện tích bề mặt của khối lớn \emph{nhỏ hơn} tổng diện tích bề mặt của các khối nhỏ.
\end{tipbox}

% ======================================================================
% VÍ DỤ MẪU
% ======================================================================
\begin{exbox}{Ví dụ 1}
Trong mặt phẳng toạ độ $Oxy$, cho ba điểm $A(-2;\,3)$, $B(4;\,3)$ và $C(4;\,-1)$.
\begin{enumerate}[label=\alph*),leftmargin=1.6em,itemsep=0pt,topsep=2pt]
  \item Mỗi điểm $A$, $B$, $C$ nằm ở góc phần tư nào?
  \item Tính độ dài các đoạn $AB$, $BC$ và diện tích tam giác $ABC$.
  \item Tìm toạ độ điểm $D$ để $ABCD$ là hình chữ nhật.
\end{enumerate}
\guide
\begin{gmFig}[scale=0.55]
  \draw[Line,line width=0.3pt] (-3.5,-2.5) grid (5.5,4.5);
  \draw[->,Navy,thick] (-3.7,0)--(5.8,0) node[right,text=Ink]{$x$};
  \draw[->,Navy,thick] (0,-2.7)--(0,4.8) node[above,text=Ink]{$y$};
  \foreach \t in {-3,-2,-1,1,2,3,4,5}{\node[below,font=\scriptsize] at (\t,0){$\t$};}
  \foreach \t in {-2,-1,1,2,3,4}{\node[left,font=\scriptsize] at (0,\t){$\t$};}
  \fill[tmBlue!20] (-2,3)--(4,3)--(4,-1)--cycle;
  \draw[tmBlue,thick] (-2,3)--(4,3)--(4,-1)--cycle;
  \draw[tmOrange,thick,dashed] (-2,3)--(-2,-1)--(4,-1);
  \draw[Ink] (4,3) rectangle (3.75,2.75);
  \fill[AnsRed] (-2,3) circle (3pt) node[above left]{$A$};
  \fill[AnsRed] (4,3) circle (3pt) node[above right]{$B$};
  \fill[AnsRed] (4,-1) circle (3pt) node[below right]{$C$};
  \fill[tmOrange] (-2,-1) circle (3pt) node[below left]{$D$};
  \node[above,font=\footnotesize,text=tmBlue] at (1,3){$6$};
  \node[right,font=\footnotesize,text=tmBlue] at (4,1){$4$};
\end{gmFig}
\step{1}{Xét dấu toạ độ để tìm góc phần tư}
$A(-2;\,3)$: hoành độ âm, tung độ dương $\Rightarrow$ góc phần tư II.
$B(4;\,3)$: cả hai dương $\Rightarrow$ góc phần tư I.
$C(4;\,-1)$: hoành độ dương, tung độ âm $\Rightarrow$ góc phần tư IV.

\step{2}{Tính độ dài các cạnh nằm ngang và thẳng đứng}
$A$ và $B$ cùng tung độ $3$, nên $AB$ nằm ngang và $AB=4-(-2)=4+2=6$ (đơn vị độ dài).
$B$ và $C$ cùng hoành độ $4$, nên $BC$ thẳng đứng và $BC=3-(-1)=3+1=4$.

\step{3}{Tính diện tích tam giác}
Cạnh $AB$ nằm ngang, cạnh $BC$ thẳng đứng nên góc $B$ vuông. Tam giác vuông có diện tích bằng nửa tích hai cạnh góc vuông:
\[ S_{ABC}=\frac{AB\cdot BC}{2}=\frac{6\cdot 4}{2}=12 \text{ (đơn vị diện tích)}. \]

\step{4}{Tìm điểm $D$}
Hình chữ nhật $ABCD$ có $D$ ``đối diện'' với $B$: $D$ phải thẳng hàng dọc với $A$ (cùng hoành độ $-2$) và thẳng hàng ngang với $C$ (cùng tung độ $-1$). Vậy $D(-2;\,-1)$, nằm ở góc phần tư III.
\verify{Hình chữ nhật $ABCD$ có diện tích $6\cdot4=24$, đúng bằng hai lần diện tích tam giác $ABC$ (đường chéo $AC$ chia hình chữ nhật thành hai nửa bằng nhau).}
\ans{a) $A$: II, $B$: I, $C$: IV. \quad b) $AB=6$, $BC=4$, $S=12$. \quad c) $D(-2;\,-1)$.}
\end{exbox}

\begin{exbox}{Ví dụ 2}
Trong hình dưới, tam giác $ABC$ có $\widehat{A}=55^\circ$, $\widehat{C}=60^\circ$. Các cạnh $AB$ và $CB$ được kéo dài qua $B$; cạnh $BC$ được kéo dài qua $C$. Tính số đo $x$ và $y$.
\begin{gmFig}[scale=0.75]
  \coordinate (B) at (0,0); \coordinate (C) at (4,0); \coordinate (A) at (1.787,3.833);
  \coordinate (L) at (-1.6,0); \coordinate (D) at (-0.6,-1.287); \coordinate (E) at (5.6,0);
  \draw[Ink,thick] (L)--(E);
  \draw[Ink,thick] (D)--(A)--(C);
  \pic[draw,tmBlue,thick,angle radius=0.5cm,"$55^\circ$",angle eccentricity=1.7]{angle=B--A--C};
  \pic[draw,tmGreen!70!black,thick,angle radius=0.45cm,"$60^\circ$",angle eccentricity=1.9]{angle=A--C--B};
  \pic[draw,AnsRed,thick,angle radius=0.45cm,"$x$",angle eccentricity=1.6]{angle=L--B--D};
  \pic[draw,tmPurple,thick,angle radius=0.4cm,"$y$",angle eccentricity=1.6]{angle=E--C--A};
  \node[above left] at (A){$A$}; \node[above left] at (B){$B$}; \node[below] at (C){$C$};
\end{gmFig}
\guide
\textbf{Phân tích.} Góc $x$ nằm \emph{bên ngoài} tam giác, nhưng nó \emph{đối đỉnh} với góc $\widehat{ABC}$ của tam giác (hai cạnh của góc $x$ là tia đối của $BC$ và tia đối của $BA$). Vì vậy ta tìm $\widehat{ABC}$ trước.

\step{1}{Tìm góc $B$ của tam giác}
\begin{recallbox}Tổng ba góc của một tam giác bằng $180^\circ$.\end{recallbox}
$\widehat{ABC}=180^\circ-(55^\circ+60^\circ)=180^\circ-115^\circ=65^\circ$.

\step{2}{Tìm $x$}
\begin{recallbox}Hai góc đối đỉnh thì bằng nhau.\end{recallbox}
Góc $x$ và góc $\widehat{ABC}$ đối đỉnh nên $x=65^\circ$.

\step{3}{Tìm $y$}
Góc $y$ kề bù với góc $\widehat{ACB}$ nên $y=180^\circ-60^\circ=120^\circ$.
\verify{Theo tính chất góc ngoài, $y$ phải bằng tổng hai góc trong không kề: $55^\circ+65^\circ=120^\circ$. Khớp!}
\ans{$x=65^\circ$, $y=120^\circ$.}
\end{exbox}

\begin{exbox}{Ví dụ 3}
\begin{enumerate}[label=\alph*),leftmargin=1.6em,itemsep=0pt,topsep=2pt]
  \item Một tam giác cân có một góc bằng $40^\circ$. Tìm số đo hai góc còn lại (xét mọi khả năng).
  \item Một tam giác cân có một góc bằng $100^\circ$. Tìm số đo hai góc còn lại.
  \item Bộ ba độ dài nào sau đây tạo thành một tam giác: $(3;\,4;\,8)$ và $(3;\,5;\,7)$? Một tam giác cân có hai cạnh dài $4$ cm và $9$ cm. Tính chu vi của nó.
\end{enumerate}
\guide
\step{1}{Câu a -- xét hai trường hợp}
\begin{recallbox}Tam giác cân có hai góc ở đáy bằng nhau; tổng ba góc bằng $180^\circ$.\end{recallbox}
\emph{Trường hợp 1: $40^\circ$ là góc ở đỉnh.} Hai góc đáy bằng nhau và có tổng $180^\circ-40^\circ=140^\circ$, nên mỗi góc là $140^\circ:2=70^\circ$. Tam giác có các góc $40^\circ,\,70^\circ,\,70^\circ$.

\emph{Trường hợp 2: $40^\circ$ là góc ở đáy.} Góc đáy còn lại cũng bằng $40^\circ$; góc ở đỉnh là $180^\circ-40^\circ-40^\circ=100^\circ$. Tam giác có các góc $40^\circ,\,40^\circ,\,100^\circ$.

\step{2}{Câu b -- chỉ có một trường hợp}
Góc $100^\circ$ không thể là góc ở đáy: nếu vậy thì hai góc đáy đã có tổng $200^\circ>180^\circ$. Vậy $100^\circ$ là góc ở đỉnh, mỗi góc đáy là $(180^\circ-100^\circ):2=40^\circ$.

\step{3}{Câu c -- dùng bất đẳng thức tam giác}
\begin{recallbox}Ba độ dài tạo thành tam giác khi và chỉ khi cạnh dài nhất nhỏ hơn tổng hai cạnh còn lại.\end{recallbox}
$(3;\,4;\,8)$: $8>3+4=7$, \emph{không} tạo thành tam giác. $\;(3;\,5;\,7)$: $7<3+5=8$, tạo thành tam giác.

Tam giác cân có hai cạnh $4$ và $9$ thì cạnh thứ ba bằng $4$ hoặc bằng $9$:
\begin{itemize}[leftmargin=1.3em,itemsep=0pt,topsep=0pt]
  \item $(4;\,4;\,9)$: $9>4+4=8$, loại;
  \item $(4;\,9;\,9)$: $9<4+9=13$, được. Chu vi $=4+9+9=22$ cm.
\end{itemize}
\ans{a) $70^\circ$ và $70^\circ$, hoặc $40^\circ$ và $100^\circ$. \ b) $40^\circ$ và $40^\circ$. \ c) chỉ $(3;5;7)$; chu vi $22$ cm.}
\end{exbox}

\begin{exbox}{Ví dụ 4}
Tam giác $OAB$ với $O(0;\,0)$, $A(3;\,0)$, $B(3;\,2)$ được quay quanh gốc $O$ \emph{theo chiều kim đồng hồ} và trở thành tam giác $OA'B'$ với $A'(0;\,-3)$, $B'(2;\,-3)$. Hỏi góc quay bằng bao nhiêu độ? Nếu quay ngược chiều kim đồng hồ thì phải quay bao nhiêu độ để được cùng kết quả?
\guide
\begin{gmFig}[scale=0.6]
  \draw[Line,line width=0.3pt] (-1.5,-3.5) grid (4.5,2.5);
  \draw[->,Navy,thick] (-1.7,0)--(4.8,0) node[right,text=Ink]{$x$};
  \draw[->,Navy,thick] (0,-3.7)--(0,2.8) node[above,text=Ink]{$y$};
  \fill[tmBlue!30,draw=tmBlue,thick] (0,0)--(3,0)--(3,2)--cycle;
  \fill[tmOrange!30,draw=tmOrange!90!black,thick] (0,0)--(0,-3)--(2,-3)--cycle;
  \node[below left] at (0,0){$O$};
  \node[below right] at (3,0){$A$}; \node[above right] at (3,2){$B$};
  \node[left] at (0,-3){$A'$}; \node[below right] at (2,-3){$B'$};
  \draw[->,AnsRed,thick] (1.3,-0.1) arc[start angle=-4,end angle=-86,radius=1.3];
  \node[AnsRed,font=\footnotesize] at (1.55,-1.3){$90^\circ$};
\end{gmFig}
\step{1}{Chọn một đoạn thẳng xuất phát từ tâm quay}
Đoạn $OA$ nằm trên tia $Ox$, hướng sang \emph{phải}. Ảnh của nó là $OA'$ với $A'(0;\,-3)$, nằm trên trục $Oy$, hướng \emph{xuống dưới}.

\step{2}{Đo góc quét theo chiều kim đồng hồ}
Kim đồng hồ đi từ hướng ``3 giờ'' (sang phải) đến hướng ``6 giờ'' (xuống dưới) thì quét một góc vuông. Vậy góc quay theo chiều kim đồng hồ là $90^\circ$.

\step{3}{Đổi sang chiều ngược lại}
\begin{recallbox}Quay $\alpha$ theo chiều kim đồng hồ cho cùng kết quả với quay $360^\circ-\alpha$ ngược chiều kim đồng hồ.\end{recallbox}
Ngược chiều kim đồng hồ: $360^\circ-90^\circ=270^\circ$.
\verify{Theo bảng ở mục 4, quay $90^\circ$ cùng chiều kim đồng hồ biến $(x;\,y)$ thành $(y;\,-x)$: $A(3;\,0)\to(0;\,-3)=A'$ và $B(3;\,2)\to(2;\,-3)=B'$. Đúng.}
\ans{$90^\circ$ theo chiều kim đồng hồ, tương đương $270^\circ$ ngược chiều kim đồng hồ.}
\end{exbox}

\begin{exbox}{Ví dụ 5}
Trên một trục số, các vạch chia cách đều nhau. Vạch đầu tiên ứng với số $0$, vạch thứ bảy ứng với số $3$ (tức là cách vạch $0$ sáu khoảng). Điểm $M$ ở vạch cách $0$ ba khoảng, điểm $N$ ở vạch cách $0$ chín khoảng.
\begin{enumerate}[label=\alph*),leftmargin=1.6em,itemsep=0pt,topsep=2pt]
  \item Tìm các số ứng với $M$, $N$ và tính $M:N$.
  \item Nếu đơn vị trên trục là mét, hãy đổi $M$ và $N$ ra xăng-ti-mét.
\end{enumerate}
\guide
\begin{gmFig}[x=0.95cm]
  \draw[->,thick,Navy] (-0.4,0)--(10.5,0);
  \foreach \i in {0,...,10}{\draw[Navy,line width=0.8pt] (\i,0.13)--(\i,-0.13);}
  \node[below=3pt] at (0,0){$0$}; \node[below=3pt] at (6,0){$3$};
  \fill[AnsRed] (3,0) circle (2.6pt) node[above=4pt,text=AnsRed]{$M$};
  \fill[AnsRed] (9,0) circle (2.6pt) node[above=4pt,text=AnsRed]{$N$};
\end{gmFig}
\step{1}{Tìm giá trị của một khoảng}
Từ $0$ đến $3$ có $6$ khoảng bằng nhau, nên mỗi khoảng ứng với $3:6=0{,}5$.

\step{2}{Tìm $M$ và $N$}
$M$ cách $0$ ba khoảng: $M=3\cdot0{,}5=1{,}5$. $\;N$ cách $0$ chín khoảng: $N=9\cdot0{,}5=4{,}5$.

\step{3}{Tính thương}
$M:N=1{,}5:4{,}5=\dfrac{15}{45}=\dfrac13\approx 0{,}33$. (Nhận xét nhanh: $N$ cách $0$ gấp $3$ lần $M$, nên $M:N=\frac13$.)

\step{4}{Đổi đơn vị}
\begin{recallbox}$1\text{ m}=100\text{ cm}$; đổi từ đơn vị lớn sang đơn vị nhỏ thì nhân.\end{recallbox}
$1{,}5\text{ m}=1{,}5\cdot100\text{ cm}=150\text{ cm}$; $\;4{,}5\text{ m}=450\text{ cm}$.
\ans{a) $M=1{,}5$; $N=4{,}5$; $M:N=\frac13$. \quad b) $150$ cm và $450$ cm.}
\end{exbox}

\begin{exbox}{Ví dụ 6}
Một hình chữ nhật dài $14$ cm, rộng $5$ cm. Trên cạnh trên, người ta đánh dấu hai điểm cách hai đỉnh trên mỗi điểm $3$ cm, rồi nối chúng với hai đỉnh dưới để được hình thang tô màu (hình vẽ).
\begin{enumerate}[label=\alph*),leftmargin=1.6em,itemsep=0pt,topsep=2pt]
  \item Tính diện tích hình thang bằng hai cách.
  \item Một tam giác có diện tích $24\text{ cm}^2$ và một cạnh đáy dài $8$ cm. Tính chiều cao ứng với cạnh đáy đó.
\end{enumerate}
\begin{gmFig}[scale=0.42]
  \fill[tmBlue!35] (0,0)--(14,0)--(11,5)--(3,5)--cycle;
  \draw[Ink,thick] (0,0) rectangle (14,5);
  \draw[Ink,thick] (0,0)--(3,5); \draw[Ink,thick] (14,0)--(11,5);
  \node[below] at (7,0){$14$ cm}; \node[right] at (14,2.5){$5$ cm};
  \node[above] at (1.5,5){$3$ cm}; \node[above] at (12.5,5){$3$ cm};
\end{gmFig}
\guide
\step{1}{Cách 1 -- dùng công thức hình thang}
\begin{recallbox}$S_{\text{hình thang}}=\dfrac{(a+b)\cdot h}{2}$, với $a$, $b$ là hai đáy, $h$ là chiều cao.\end{recallbox}
Đáy lớn $a=14$ cm. Đáy nhỏ $b=14-3-3=8$ cm. Chiều cao bằng chiều rộng hình chữ nhật: $h=5$ cm.
\[ S=\frac{(14+8)\cdot 5}{2}=\frac{22\cdot5}{2}=\frac{110}{2}=55\ (\text{cm}^2). \]

\step{2}{Cách 2 -- hình lớn trừ phần thừa}
Diện tích hình chữ nhật: $14\cdot5=70\ (\text{cm}^2)$. Hai phần không tô màu là hai tam giác vuông, mỗi tam giác có hai cạnh góc vuông $3$ cm và $5$ cm, diện tích $\dfrac{3\cdot5}{2}=7{,}5\ (\text{cm}^2)$. Diện tích hình thang: $70-2\cdot7{,}5=70-15=55\ (\text{cm}^2)$.

\step{3}{Câu b -- tính ngược chiều cao}
Từ $S=\dfrac{a\cdot h}{2}$ ta có $h=\dfrac{2\cdot S}{a}=\dfrac{2\cdot24}{8}=\dfrac{48}{8}=6$ (cm).
\verify{$\dfrac{8\cdot 6}{2}=24$. Đúng.}
\ans{a) $55\text{ cm}^2$. \quad b) $6$ cm.}
\end{exbox}

\begin{exbox}{Ví dụ 7}
Một vòng quay may mắn hình tròn tâm $O$ được chia thành bốn ô $A$, $B$, $C$, $D$: ô $A$ là nửa hình tròn, ô $B$ là một phần tư hình tròn, phần còn lại được chia thành ô $C$ và ô $D$ sao cho góc ở tâm của ô $C$ gấp rưỡi góc ở tâm của ô $D$. Mỗi ô chiếm bao nhiêu phần trăm diện tích vòng quay? Khi quay, kim dừng ở ô $C$ trong khoảng bao nhiêu phần trăm số lần?
\begin{gmFig}[scale=0.9]
  \fill[tmBlue!20] (0,0)--(90:1.7) arc[start angle=90,end angle=270,radius=1.7]--cycle;
  \fill[tmGreen!22] (0,0)--(270:1.7) arc[start angle=270,end angle=360,radius=1.7]--cycle;
  \fill[tmOrange!30] (0,0)--(0:1.7) arc[start angle=0,end angle=54,radius=1.7]--cycle;
  \fill[tmPurple!22] (0,0)--(54:1.7) arc[start angle=54,end angle=90,radius=1.7]--cycle;
  \draw[Ink,thick] (0,0) circle (1.7);
  \draw[Ink,thick] (90:1.7)--(270:1.7); \draw[Ink,thick] (0,0)--(0:1.7); \draw[Ink,thick] (0,0)--(54:1.7);
  \draw[Ink] (0,0) rectangle (0.2,-0.2);
  \node at (180:0.9){$A$}; \node at (315:0.95){$B$}; \node at (27:1.05){$C$}; \node at (72:1.15){$D$};
  \node[below left,font=\footnotesize] at (0,0){$O$};
\end{gmFig}
\guide
\step{1}{Góc ở tâm của ô $A$ và ô $B$}
Cả vòng tròn ứng với $360^\circ$. Ô $A$ (nửa hình tròn) ứng với $180^\circ$, ô $B$ (một phần tư) ứng với $90^\circ$.

\step{2}{Góc ở tâm của ô $C$ và ô $D$}
Phần còn lại: $360^\circ-180^\circ-90^\circ=90^\circ$. Ta cần chia $90^\circ$ thành hai phần, phần $C$ gấp rưỡi phần $D$, tức là nếu $D$ gồm $2$ phần bằng nhau thì $C$ gồm $3$ phần như thế. Tổng $5$ phần ứng với $90^\circ$, nên một phần là $90^\circ:5=18^\circ$. Vậy góc của $D$ là $2\cdot18^\circ=36^\circ$, góc của $C$ là $3\cdot18^\circ=54^\circ$.

\step{3}{Đổi ra phần trăm}
\begin{recallbox}Hình quạt có góc ở tâm $n^\circ$ chiếm $\dfrac{n}{360}$ hình tròn.\end{recallbox}
$A$: $\dfrac{180}{360}=50\%$; $\;B$: $\dfrac{90}{360}=25\%$; $\;C$: $\dfrac{54}{360}=0{,}15=15\%$; $\;D$: $\dfrac{36}{360}=0{,}1=10\%$.
\verify{$50\%+25\%+15\%+10\%=100\%$, và $54^\circ=1{,}5\cdot36^\circ$. Đúng.}
\ans{$A$: $50\%$, $B$: $25\%$, $C$: $15\%$, $D$: $10\%$; kim dừng ở ô $C$ khoảng $15\%$ số lần.}
\end{exbox}

\begin{exbox}{Ví dụ 8}
\begin{enumerate}[label=\alph*),leftmargin=1.6em,itemsep=0pt,topsep=2pt]
  \item Cần bao nhiêu khối lập phương cạnh $1$ cm để xếp kín một hình hộp chữ nhật kích thước $5\text{ cm}\times3\text{ cm}\times2\text{ cm}$? Tính diện tích toàn phần của hình hộp đó.
  \item Một hình lăng trụ đứng có đáy là tam giác vuông với hai cạnh góc vuông $3$ cm và $4$ cm, chiều cao lăng trụ $10$ cm. Tính thể tích. Nếu một lăng trụ có cùng đáy như vậy nhưng thể tích $90\text{ cm}^3$ thì chiều cao của nó bằng bao nhiêu?
\end{enumerate}
\guide
\step{1}{Câu a -- đếm theo lớp}
\begin{recallbox}$V_{\text{hộp}}=a\cdot b\cdot c$; $\;S_{\text{toàn phần}}=2\cdot(ab+bc+ca)$.\end{recallbox}
Lớp đáy gồm $5\cdot3=15$ khối; có $2$ lớp nên cần $15\cdot2=30$ khối (thể tích $30\text{ cm}^3$).

Ba cặp mặt đối diện có diện tích $5\cdot3=15$, $\;3\cdot2=6$, $\;5\cdot2=10\ (\text{cm}^2)$, nên
$S_{\text{toàn phần}}=2\cdot(15+6+10)=2\cdot31=62\ (\text{cm}^2)$.

\step{2}{Câu b -- diện tích đáy rồi nhân chiều cao}
\begin{recallbox}$V_{\text{lăng trụ}}=S_{\text{đáy}}\cdot h$.\end{recallbox}
Đáy là tam giác vuông: $S_{\text{đáy}}=\dfrac{3\cdot4}{2}=6\ (\text{cm}^2)$. Thể tích: $V=6\cdot10=60\ (\text{cm}^3)$.

Với thể tích $90\text{ cm}^3$: $h=V:S_{\text{đáy}}=90:6=15$ (cm).
\verify{$6\cdot15=90$. Đúng.}
\ans{a) $30$ khối; $62\text{ cm}^2$. \quad b) $60\text{ cm}^3$; $15$ cm.}
\end{exbox}

\begin{parentbox}
Sau mỗi ví dụ, hãy yêu cầu con \emph{nói lại bằng lời} công thức vừa dùng và \emph{vì sao} nó đúng (ví dụ: ``tam giác bằng nửa hình chữ nhật cùng đáy, cùng cao''). Hiểu được ``vì sao'' giúp con không phải học thuộc máy móc và ít nhầm khi gặp hình lạ.
\end{parentbox}

\clearpage
\dividerbar{Luyện tập cơ bản}

% ======================================================================
% BÀI 1--15
% ======================================================================
\begin{probbox}{1}{\relax}
Trong mặt phẳng toạ độ dưới đây, mỗi vạch chia trên hai trục ứng với $1$ đơn vị. Toạ độ của điểm $P$ có thể là cặp số nào?
\begin{gmFig}[scale=0.42]
  \draw[->,Ink,thick] (-4.8,0)--(5.2,0) node[right]{$x$};
  \draw[->,Ink,thick] (0,-4.8)--(0,5.2) node[above]{$y$};
  \foreach \t in {-4,...,4}{\draw[Ink] (\t,0.15)--(\t,-0.15); \draw[Ink] (0.15,\t)--(-0.15,\t);}
  \foreach \t in {-4,-2,2,4}{\node[below,font=\scriptsize] at (\t,-0.15){$\t$}; \node[right,font=\scriptsize] at (0.15,\t){$\t$};}
  \fill[AnsRed] (-1,3) circle (4pt) node[left=2pt,text=Ink]{$P$};
\end{gmFig}
\choices{$(1;\,3)$}{$(1;\,-3)$}{$(-3;\,1)$}{$(3;\,-1)$}{$(-1;\,3)$}
\end{probbox}
\probsol{
\textbf{Phân tích đề.} Ta cần đọc vị trí của $P$ so với hai trục để xác định \emph{dấu} và \emph{độ lớn} của hoành độ, tung độ, rồi so với năm phương án.
\begin{recallbox}Điểm $M(x_0;\,y_0)$: $x_0$ (hoành độ) cho biết điểm nằm bên phải ($x_0>0$) hay bên trái ($x_0<0$) trục tung; $y_0$ (tung độ) cho biết điểm nằm phía trên ($y_0>0$) hay phía dưới ($y_0<0$) trục hoành. Hoành độ luôn viết trước.\end{recallbox}
\step{1}{Xác định dấu của toạ độ}
$P$ nằm \emph{bên trái} trục tung $Oy$, nên hoành độ của $P$ là số \textbf{âm}.
$P$ nằm \emph{phía trên} trục hoành $Ox$, nên tung độ của $P$ là số \textbf{dương}.
Vậy $P$ ở góc phần tư II, dạng $(-;\,+)$.

\step{2}{Loại bớt phương án}
\begin{itemize}[leftmargin=1.3em,itemsep=0pt,topsep=0pt]
  \item (A) $(1;\,3)$: dạng $(+;+)$, góc phần tư I -- loại.
  \item (B) $(1;\,-3)$: dạng $(+;-)$, góc phần tư IV -- loại.
  \item (D) $(3;\,-1)$: dạng $(+;-)$ -- loại.
  \item Còn lại (C) $(-3;\,1)$ và (E) $(-1;\,3)$, cả hai đều dạng $(-;+)$.
\end{itemize}

\step{3}{So sánh độ lớn}
Từ $P$ kẻ đường vuông góc xuống trục $Ox$: chân đường vuông góc nằm giữa $0$ và $-2$, đúng ở vạch $-1$. Kẻ đường vuông góc sang trục $Oy$: chân nằm giữa $2$ và $4$, đúng ở vạch $3$. Như vậy $P$ ở rất gần trục tung (cách $1$ đơn vị) và ở khá cao (cách trục hoành $3$ đơn vị): $P(-1;\,3)$.

Điểm $(-3;\,1)$ ở phương án (C) thì ngược lại: xa trục tung ($3$ đơn vị) và thấp ($1$ đơn vị), không khớp với hình.
\begin{warnbox}Phương án (C) là ``bẫy'' dành cho bạn đọc tung độ trước, hoành độ sau. Hãy nhớ: \emph{đi ngang trước, đi dọc sau}.\end{warnbox}
\verify{Từ $O$ đi sang trái $1$ vạch, rồi đi lên $3$ vạch: ta đến đúng vị trí của $P$.}
Vậy chọn \textbf{(E)}.
\ans{(E) $(-1;\,3)$}
}

\begin{probbox}{2}{\relax}
Một tam giác có hai góc lần lượt bằng $25^\circ$ và $70^\circ$. Góc thứ ba của tam giác đó bằng
\choices{$85^\circ$}{$105^\circ$}{$65^\circ$}{$95^\circ$}{$75^\circ$}
\end{probbox}
\probsol{
\textbf{Phân tích đề.} Đã biết hai góc của tam giác, cần tìm góc còn lại. Ta dùng định lí tổng ba góc.
\begin{recallbox}Tổng ba góc của một tam giác bằng $180^\circ$.\end{recallbox}
\step{1}{Tính tổng hai góc đã biết}
$25^\circ+70^\circ=95^\circ$.

\step{2}{Lấy $180^\circ$ trừ đi tổng đó}
Góc thứ ba $=180^\circ-95^\circ=85^\circ$.
\begin{gmFig}[scale=0.8]
  \coordinate (A) at (0,0); \coordinate (B) at (4,0); \coordinate (C) at (3.419,1.595);
  \fill[tmBlue!8] (A)--(B)--(C)--cycle;
  \draw[Ink,thick] (A)--(B)--(C)--cycle;
  \pic[draw,tmBlue,thick,angle radius=0.9cm,"$25^\circ$",angle eccentricity=1.35]{angle=B--A--C};
  \pic[draw,tmGreen!70!black,thick,angle radius=0.35cm,"$70^\circ$",angle eccentricity=2.0]{angle=C--B--A};
  \pic[draw,AnsRed,thick,angle radius=0.35cm,"$85^\circ$",angle eccentricity=1.9]{angle=A--C--B};
\end{gmFig}
\begin{warnbox}Phương án (D) $95^\circ$ là tổng của hai góc đã cho, không phải góc thứ ba. Đừng dừng lại ở bước 1!\end{warnbox}
\verify{$25^\circ+70^\circ+85^\circ=180^\circ$. Đúng.}
Vậy chọn \textbf{(A)}.
\ans{(A) $85^\circ$}
}

\begin{probbox}{3}{\relax}
Quay tứ giác tô màu quanh điểm $Z$ \emph{theo chiều kim đồng hồ} (chiều của mũi tên trong hình) thì được tứ giác không tô màu. Góc quay xấp xỉ bằng
\begin{gmFig}[scale=0.85]
  \fill[Slate!70,draw=Ink,thick] (0,0)--(3,0)--(1.87,1.47)--(0.37,1.47)--cycle;
  \fill[white,draw=Ink,thick] (0,0)--(0,3)--(-1.47,1.87)--(-1.47,0.37)--cycle;
  \fill (0,0) circle (2.5pt);
  \node[below=6pt] at (0,0){$Z$};
  \draw[->,thick] (0.75,-0.1) arc[start angle=-8,end angle=-200,radius=0.75];
\end{gmFig}
\choices{$180^\circ$}{$270^\circ$}{$360^\circ$}{$45^\circ$}{$135^\circ$}
\end{probbox}
\probsol{
\textbf{Phân tích đề.} Hình bị quay quanh $Z$ theo chiều kim đồng hồ. Muốn biết góc quay, ta theo dõi \emph{một cạnh xuất phát từ $Z$} và xem nó đã quét qua góc bao nhiêu.
\begin{recallbox}Trong phép quay, mọi điểm quay cùng một góc quanh tâm; độ dài các cạnh không đổi. Một vòng trọn là $360^\circ$; quay $\alpha$ theo chiều kim đồng hồ giống quay $360^\circ-\alpha$ ngược chiều kim đồng hồ.\end{recallbox}
\step{1}{Tìm cạnh tương ứng}
Tứ giác tô màu có cạnh đáy dài nhất xuất phát từ $Z$, nằm ngang và hướng sang \emph{phải} (hướng ``3 giờ'' trên mặt đồng hồ). Ở tứ giác không tô màu, cạnh dài nhất xuất phát từ $Z$ là cạnh \emph{thẳng đứng}, hướng \emph{lên trên} (hướng ``12 giờ''). Vì phép quay giữ nguyên độ dài, cạnh nằm ngang đã biến thành cạnh thẳng đứng này.

\step{2}{Đo góc quét theo chiều kim đồng hồ}
Đi theo chiều kim đồng hồ từ hướng 3 giờ đến hướng 12 giờ, ta phải đi qua hướng 6 giờ (xuống dưới), hướng 9 giờ (sang trái), rồi mới lên 12 giờ: đó là $3$ góc vuông.
\[ 3\cdot 90^\circ=270^\circ. \]
\begin{gmFig}[scale=0.7]
  \draw[->,tmBlue,thick] (0,0)--(1.6,0) node[right,font=\footnotesize]{3 giờ (ban đầu)};
  \draw[->,tmOrange!90!black,thick] (0,0)--(0,1.6) node[above,font=\footnotesize]{12 giờ (ảnh)};
  \draw[Slate,dashed] (0,0)--(0,-1.3) node[below,font=\footnotesize]{6 giờ};
  \draw[Slate,dashed] (0,0)--(-1.3,0) node[left,font=\footnotesize]{9 giờ};
  \draw[->,AnsRed,thick] (0.9,-0.1) arc[start angle=-6,end angle=-264,radius=0.9];
  \node[AnsRed,font=\footnotesize] at (1.3,-1.1){$270^\circ$};
\end{gmFig}

\step{3}{Kiểm tra với cạnh thứ hai}
Cạnh xiên ngắn của tứ giác tô màu đi từ $Z$ lên phía trên, hơi lệch sang phải. Quay thêm $270^\circ$ theo chiều kim đồng hồ (tức $90^\circ$ ngược chiều), nó phải hướng sang trái, hơi lệch lên trên -- đúng là cạnh dưới của tứ giác không tô màu.
\begin{warnbox}Nếu quay \emph{ngược} chiều kim đồng hồ thì chỉ cần $90^\circ$, nhưng đề yêu cầu chiều kim đồng hồ (chiều mũi tên), mà $90^\circ$ không có trong các phương án. Đừng nhầm hai chiều quay.\end{warnbox}
\verify{$270^\circ+90^\circ=360^\circ$: quay $270^\circ$ theo chiều kim đồng hồ và quay $90^\circ$ ngược chiều cho cùng một kết quả, và quay $90^\circ$ ngược chiều rõ ràng đưa cạnh ``sang phải'' thành cạnh ``lên trên''.}
Vậy chọn \textbf{(B)}.
\ans{(B) $270^\circ$}
\begin{parentbox}Lấy một chiếc đồng hồ (hoặc vẽ mặt đồng hồ) và cho con xoay kim phút từ số $3$ theo chiều kim đồng hồ đến số $12$: kim đi qua $9$ ``khoảng $5$ phút'' $=45$ phút, tức $\frac34$ vòng $=\frac34\cdot360^\circ=270^\circ$. Hoạt động này giúp con cảm nhận được sự khác nhau giữa hai chiều quay.\end{parentbox}
}

\begin{probbox}{4}{\relax}
Trong hình vẽ, hai cạnh của tam giác được kéo dài như hình. Giá trị của $x$ là
\begin{gmFig}[scale=0.7]
  \coordinate (A) at (0,0); \coordinate (B) at (4,0); \coordinate (C) at (1.547,4.249);
  \coordinate (L) at (-1.5,0); \coordinate (D) at (-0.5,-1.374);
  \draw[Ink,thick] (L)--(B)--(C)--(D);
  \pic[draw,thick,angle radius=0.5cm,"$50^\circ$",angle eccentricity=1.75]{angle=A--C--B};
  \pic[draw,thick,angle radius=0.5cm,"$60^\circ$",angle eccentricity=1.8]{angle=C--B--A};
  \pic[draw,AnsRed,thick,angle radius=0.45cm,"$x^\circ$",angle eccentricity=1.75]{angle=L--A--D};
\end{gmFig}
\choices{$110$}{$50$}{$10$}{$60$}{$70$}
\end{probbox}
\probsol{
\textbf{Phân tích đề.} Hình gồm một tam giác có hai góc $50^\circ$ và $60^\circ$. Ở đỉnh thứ ba (đỉnh dưới bên trái), cạnh đáy và cạnh bên được kéo dài ra ngoài, tạo thành góc $x^\circ$. Ta tìm góc trong ở đỉnh đó trước, rồi nối nó với $x$.
\begin{gmFig}[scale=0.7]
  \coordinate (A) at (0,0); \coordinate (B) at (4,0); \coordinate (C) at (1.547,4.249);
  \coordinate (L) at (-1.5,0); \coordinate (D) at (-0.5,-1.374);
  \draw[Ink,thick] (L)--(B)--(C)--(D);
  \pic[draw,thick,angle radius=0.5cm,"$50^\circ$",angle eccentricity=1.75]{angle=A--C--B};
  \pic[draw,thick,angle radius=0.5cm,"$60^\circ$",angle eccentricity=1.8]{angle=C--B--A};
  \pic[draw,AnsRed,thick,angle radius=0.45cm,"$x^\circ$",angle eccentricity=1.75]{angle=L--A--D};
  \pic[draw,tmBlue,thick,fill=tmBlue!15,angle radius=0.5cm,"$y$",angle eccentricity=1.6]{angle=B--A--C};
  \node[below right,font=\footnotesize] at (A){$A$};
\end{gmFig}
\step{1}{Tìm góc trong $y$ tại đỉnh $A$}
\begin{recallbox}Tổng ba góc của một tam giác bằng $180^\circ$.\end{recallbox}
Gọi $y$ là góc trong của tam giác tại đỉnh $A$:
\[ y=180^\circ-(50^\circ+60^\circ)=180^\circ-110^\circ=70^\circ. \]

\step{2}{Nhận ra cặp góc đối đỉnh}
Góc $x^\circ$ được tạo bởi tia kéo dài của cạnh đáy (sang trái) và tia kéo dài của cạnh bên (xuống dưới). Đó chính là hai tia đối của hai cạnh tạo nên góc $y$. Vậy $x^\circ$ và $y$ là \textbf{hai góc đối đỉnh}.
\begin{recallbox}Hai góc đối đỉnh thì bằng nhau.\end{recallbox}
Do đó $x^\circ=y=70^\circ$, tức $x=70$.
\begin{warnbox}Phương án (A) $110$ là kết quả nếu nhầm $x$ với góc \emph{kề bù} với $y$ (góc nằm ngay cạnh, phía trên đường đáy kéo dài). Hãy nhìn kĩ: góc $x$ nằm \emph{đối diện} với góc trong qua đỉnh $A$.\end{warnbox}
\verify{Góc kề bù với $y$ là $180^\circ-70^\circ=110^\circ$; góc $x$ lại kề bù với góc $110^\circ$ đó, nên $x^\circ=180^\circ-110^\circ=70^\circ$. Hai cách cho cùng một kết quả.}
Vậy chọn \textbf{(E)}.
\ans{(E) $70$}
}

\begin{probbox}{5}{\relax}
Hình vuông $PQRS$ được chia thành bốn hình vuông nhỏ bằng nhau như hình vẽ. Phần tô màu chiếm bao nhiêu phần trăm diện tích hình vuông $PQRS$?
\begin{gmFig}[scale=0.8]
  \fill[Slate!70] (1,1) rectangle (2,2);
  \draw[Ink,thick] (0,0) rectangle (2,2);
  \draw[Ink,thick] (1,0)--(1,2); \draw[Ink,thick] (0,1)--(2,1);
  \node[above left] at (0,2){$P$}; \node[above right] at (2,2){$Q$};
  \node[below right] at (2,0){$R$}; \node[below left] at (0,0){$S$};
\end{gmFig}
\choices{$33\%$}{$40\%$}{$25\%$}{$50\%$}{$15\%$}
\end{probbox}
\probsol{
\textbf{Phân tích đề.} Hình vuông lớn được chia thành $4$ phần \emph{bằng nhau}, trong đó $1$ phần được tô màu. Ta đổi tỉ số này ra phần trăm.
\begin{recallbox}Nếu một hình được chia thành $n$ phần bằng nhau và tô $m$ phần thì phần tô màu chiếm $\dfrac{m}{n}$ hình, tức $\dfrac{m}{n}\cdot100\%$.\end{recallbox}
\step{1}{Viết phân số}
Phần tô màu là $1$ trong $4$ hình vuông nhỏ bằng nhau: chiếm $\dfrac14$ hình vuông $PQRS$.

\step{2}{Đổi ra phần trăm}
$\dfrac14=\dfrac{25}{100}=25\%$.
\verify{Nếu cạnh $PQRS$ là $2$ cm thì diện tích lớn là $4\text{ cm}^2$, ô tô màu có cạnh $1$ cm, diện tích $1\text{ cm}^2$; $1:4=0{,}25=25\%$.}
\begin{warnbox}Phương án (A) $33\%$ xuất hiện nếu so phần tô màu với phần \emph{không} tô màu ($1$ so với $3$). Phần trăm phải so với \emph{toàn bộ} hình.\end{warnbox}
Vậy chọn \textbf{(C)}.
\ans{(C) $25\%$}
}

\begin{probbox}{6}{\relax}
Một khối lập phương chỉ được sơn đen một phần của \emph{đúng một mặt} như hình bên (phần sơn là tam giác vuông ở một góc của mặt, hai cạnh góc vuông dài bằng nửa cạnh khối lập phương). Năm mặt còn lại không sơn. Sau khi lăn khối lập phương, hình nào dưới đây có thể là chính khối lập phương đó?
\begin{gmFig}
  \gmCube{0}{}{\fill[Ink] (0.5,1)--(1,1)--(1,0.5)--cycle;}
\end{gmFig}
\begin{gmFig}[scale=0.88]
  \gmCube{0}{(A)}{\fill[Ink] (0,0)--(1,0)--(1,1)--cycle;}
  \gmCube{2.9}{(B)}{\fill[Ink] (0.5,0)--(1,0)--(1,0.5)--cycle;}
  \gmCube{5.8}{(C)}{\fill[Ink] (0.5,0)--(1,0)--(0.5,0.5)--cycle;}
  \gmCube{8.7}{(D)}{\fill[Ink] (0.5,1)--(1,1)--(1,0.5)--cycle; \fill[Ink] (1,0.5)--(1,0)--(0.5,0)--cycle;}
  \gmCube{11.6}{(E)}{\fill[Ink] (0,1)--(0.5,1)--(0,0.5)--cycle; \fill[Ink] (0.5,0)--(1,0)--(1,0.5)--cycle;}
\end{gmFig}
\end{probbox}
\probsol{
\textbf{Phân tích đề.} Khối lập phương chỉ có \emph{một} mặt được sơn, và trên mặt ấy chỉ có \emph{một} tam giác đen. Ta mô tả thật chính xác tam giác đó rồi so với từng phương án.
\begin{recallbox}Khi lăn khối lập phương, mỗi mặt được mang theo nguyên vẹn: họa tiết trên mặt có thể đổi chỗ hoặc bị xoay $90^\circ$, $180^\circ$, $270^\circ$, nhưng \emph{không} thay đổi kích thước, \emph{không} bị lật gương và \emph{không} có thêm phần sơn.\end{recallbox}
\step{1}{Mô tả họa tiết ban đầu}
Ba đặc điểm của họa tiết:
\begin{itemize}[leftmargin=1.3em,itemsep=0pt,topsep=0pt]
  \item chỉ có \textbf{một} tam giác đen;
  \item tam giác có \textbf{góc vuông nằm đúng ở một đỉnh} của mặt vuông;
  \item hai cạnh góc vuông dài bằng \textbf{nửa cạnh} của mặt (tam giác chiếm $\frac18$ mặt).
\end{itemize}

\step{2}{Xét từng phương án}
\begin{itemize}[leftmargin=1.3em,itemsep=1pt,topsep=0pt]
  \item (A): tam giác đen chiếm \emph{nửa} mặt (hai cạnh góc vuông bằng cả cạnh mặt) -- quá lớn, loại.
  \item (B): một tam giác đen, góc vuông ở đỉnh dưới bên phải của mặt, hai cạnh góc vuông bằng nửa cạnh -- khớp cả ba đặc điểm.
  \item (C): một tam giác cỡ đúng, nhưng góc vuông nằm ở \emph{trung điểm} cạnh dưới, không phải ở một đỉnh của mặt -- loại.
  \item (D) và (E): có \emph{hai} tam giác đen, nhiều sơn hơn khối ban đầu -- loại.
\end{itemize}

\step{3}{Kiểm tra rằng (B) thật sự đạt được}
Họa tiết (B) chính là họa tiết ban đầu \emph{xoay $90^\circ$ theo chiều kim đồng hồ} trong mặt phẳng của mặt đó: đỉnh trên bên phải chuyển thành đỉnh dưới bên phải. Đặt khối lập phương lên bàn với mặt sơn hướng về phía ta rồi xoay khối một phần tư vòng quanh trục vuông góc với mặt ấy (lăn qua lại vài lần là làm được), ta nhận đúng hình (B).
\begin{tipbox}Để kiểm tra, cắt một hình vuông giấy, tô tam giác ở một góc và xoay tờ giấy: tam giác chỉ có thể ``chạy'' qua bốn góc, không bao giờ nằm ở giữa cạnh như (C).\end{tipbox}
Vậy chọn \textbf{(B)}.
\ans{(B)}
}

\begin{probbox}{7}{\relax}
Trên một trục số, các vạch chia cách đều nhau. Các số $P$, $Q$, $R$, $S$, $T$, $U$ được đánh dấu như hình. Số nào dưới đây gần nhất với giá trị của $R:U$?
\begin{gmFig}[x=1.0cm]
  \draw[->,thick] (-0.3,0)--(6.6,0);
  \foreach \i in {0,...,6}{\draw[thick] (\i,0.15)--(\i,-0.15);}
  \foreach \i/\l in {1/P,2/Q,3/R,4/S,5/T,6/U}{\node[above=3pt] at (\i,0){$\l$};}
  \node[below=4pt] at (0,0){$0$}; \node[below=4pt] at (4,0){$1{,}0$};
\end{gmFig}
\choices{$0{,}25$}{$0{,}50$}{$0{,}75$}{$1{,}25$}{$1{,}50$}
\end{probbox}
\probsol{
\textbf{Phân tích đề.} Hai vạch đã biết giá trị là $0$ và $1{,}0$ (vạch $S$). Từ đó ta tìm giá trị của mỗi khoảng, rồi suy ra $R$ và $U$.
\begin{recallbox}Trên trục số có vạch chia đều: giá trị một khoảng $=$ (hiệu hai số đã biết) $:$ (số \emph{khoảng} giữa chúng).\end{recallbox}
\step{1}{Đếm số khoảng từ $0$ đến $1{,}0$}
Từ vạch $0$ đến vạch $S$ (ứng với $1{,}0$) có các khoảng: $0\to P$, $P\to Q$, $Q\to R$, $R\to S$, tức là $4$ khoảng.

\step{2}{Tính giá trị một khoảng}
$1{,}0:4=0{,}25$.

\step{3}{Tìm $R$ và $U$}
$R$ cách $0$ ba khoảng: $R=3\cdot0{,}25=0{,}75$.
$U$ cách $0$ sáu khoảng: $U=6\cdot0{,}25=1{,}5$.
\begin{gmFig}[x=1.0cm]
  \draw[->,thick] (-0.3,0)--(6.6,0);
  \foreach \i in {0,...,6}{\draw[thick] (\i,0.15)--(\i,-0.15);}
  \foreach \i/\l/\v in {0/0/0,1/P/0{,}25,2/Q/0{,}5,3/R/0{,}75,4/S/1,5/T/1{,}25,6/U/1{,}5}{\node[below=4pt,font=\footnotesize,text=tmBlue] at (\i,0){$\v$};}
  \foreach \i/\l in {1/P,2/Q,3/R,4/S,5/T,6/U}{\node[above=3pt] at (\i,0){$\l$};}
\end{gmFig}

\step{4}{Tính thương}
$R:U=0{,}75:1{,}5=\dfrac{75}{150}=\dfrac12=0{,}5$.

\textbf{Cách nhìn nhanh:} $R$ cách $0$ ba khoảng, $U$ cách $0$ sáu khoảng, nên $U$ gấp đôi $R$ và $R:U=\frac12$ -- không cần biết mỗi khoảng là bao nhiêu!
\begin{warnbox}Phương án (C) $0{,}75$ là giá trị của $R$, chưa chia cho $U$. Ngoài ra, nếu đếm $5$ \emph{vạch} từ $0$ đến $S$ thay vì $4$ \emph{khoảng}, ta sẽ tính sai giá trị mỗi khoảng.\end{warnbox}
\verify{$1{,}5\cdot0{,}5=0{,}75=R$. Đúng.}
Vậy chọn \textbf{(B)}.
\ans{(B) $0{,}50$}
}

\begin{probbox}{8}{\relax}
Một tam giác cân có một góc bằng $50^\circ$. Hai góc còn lại của tam giác này có thể là
\choicesII{$50^\circ$ và $90^\circ$}{$40^\circ$ và $50^\circ$}{$50^\circ$ và $80^\circ$}{$30^\circ$ và $100^\circ$}{$60^\circ$ và $70^\circ$}
\end{probbox}
\probsol{
\textbf{Phân tích đề.} Đề không nói góc $50^\circ$ là góc ở đỉnh hay góc ở đáy, nên ta xét cả hai trường hợp, rồi đối chiếu với các phương án.
\begin{recallbox}Tam giác cân có hai góc ở đáy bằng nhau; tổng ba góc của tam giác bằng $180^\circ$.\end{recallbox}
\step{1}{Trường hợp $50^\circ$ là góc ở đỉnh}
Hai góc đáy bằng nhau, có tổng $180^\circ-50^\circ=130^\circ$, nên mỗi góc là $130^\circ:2=65^\circ$. Hai góc còn lại: $65^\circ$ và $65^\circ$ (không có trong các phương án).

\step{2}{Trường hợp $50^\circ$ là góc ở đáy}
Góc đáy kia cũng bằng $50^\circ$. Góc ở đỉnh: $180^\circ-50^\circ-50^\circ=80^\circ$. Hai góc còn lại: $50^\circ$ và $80^\circ$ -- đây là phương án (C).

\step{3}{Kiểm tra vì sao các phương án khác sai}
Mỗi bộ ba góc phải có tổng $180^\circ$ \emph{và} có hai góc bằng nhau:
\begin{center}
\begin{tabular}{|c|c|c|c|}
\hline
Phương án & Ba góc & Tổng & Có hai góc bằng nhau? \\ \hline
(A) & $50^\circ,\,50^\circ,\,90^\circ$ & $190^\circ$ & có, nhưng tổng sai \\ \hline
(B) & $50^\circ,\,40^\circ,\,50^\circ$ & $140^\circ$ & có, nhưng tổng sai \\ \hline
(C) & $50^\circ,\,50^\circ,\,80^\circ$ & $180^\circ$ & có \\ \hline
(D) & $50^\circ,\,30^\circ,\,100^\circ$ & $180^\circ$ & không \\ \hline
(E) & $50^\circ,\,60^\circ,\,70^\circ$ & $180^\circ$ & không \\ \hline
\end{tabular}
\end{center}
Chỉ (C) thỏa mãn cả hai điều kiện.
\verify{$50^\circ+50^\circ+80^\circ=180^\circ$ và có hai góc bằng nhau ($50^\circ$).}
Vậy chọn \textbf{(C)}.
\ans{(C) $50^\circ$ và $80^\circ$}
}

\begin{probbox}{9}{\relax}
Bạn Minh đo được chiều dài của chiếc bàn học là $3{,}5$ m. Độ dài đó bằng bao nhiêu xăng-ti-mét?
\choices{$350$}{$30{,}5$}{$3{,}05$}{$3{,}5$}{$305$}
\end{probbox}
\probsol{
\textbf{Phân tích đề.} Đổi từ mét (đơn vị lớn) sang xăng-ti-mét (đơn vị nhỏ), nên ta \emph{nhân}.
\begin{recallbox}$1\text{ m}=10\text{ dm}=100\text{ cm}$. Đổi từ đơn vị lớn sang đơn vị nhỏ thì nhân, từ đơn vị nhỏ sang đơn vị lớn thì chia.\end{recallbox}
\step{1}{Viết phép nhân}
$3{,}5\text{ m}=3{,}5\cdot100\text{ cm}$.

\step{2}{Nhân với $100$}
Nhân một số thập phân với $100$ thì dời dấu phẩy sang phải $2$ chữ số: $3{,}5\to 35\to 350$. Vậy $3{,}5\text{ m}=350\text{ cm}$.
\begin{gmFig}[x=1.35cm]
  \foreach \u [count=\i] in {km,hm,dam,m,dm,cm,mm}{
    \node[draw=tmBlue,fill=tmBlue!8,rounded corners=2pt,minimum width=1.0cm,minimum height=0.6cm] (u\i) at (\i,0){\u};}
  \foreach \i [evaluate=\i as \j using int(\i+1)] in {1,...,6}{
    \draw[->,tmOrange!90!black,thick] (u\i.north east) to[bend left=40] node[above,font=\scriptsize]{$\times10$} (u\j.north west);}
  \node[font=\footnotesize,text=AnsRed] at (5,-0.75){m $\to$ dm $\to$ cm: $\times10\times10=\times100$};
\end{gmFig}

\textbf{Cách hiểu khác:} $3{,}5\text{ m}=3\text{ m}+0{,}5\text{ m}=300\text{ cm}+50\text{ cm}=350\text{ cm}$ (nửa mét là $50$ cm).
\begin{warnbox}Phương án (E) $305$ là lỗi đọc ``$3$ mét $5$ xăng-ti-mét''. Nhưng $0{,}5$ m là \emph{nửa mét} $=50$ cm, không phải $5$ cm.\end{warnbox}
\verify{$350:100=3{,}5$. Đúng.}
Vậy chọn \textbf{(A)}.
\ans{(A) $350$}
}

\begin{probbox}{10}{\relax}
Bạn An muốn xếp một khối hộp chữ nhật đặc như hình vẽ từ các khối lập phương nhỏ kích thước $1\text{ cm}\times1\text{ cm}\times1\text{ cm}$. An cần bao nhiêu khối lập phương nhỏ?
\begin{gmFig}[scale=0.7]
  \draw[Ink,thick,fill=white] (0,0) rectangle (4,3);
  \draw[Ink,thick,fill=Tint] (0,3)--(1.5,4)--(5.5,4)--(4,3)--cycle;
  \draw[Ink,thick,fill=Line!50] (4,0)--(5.5,1)--(5.5,4)--(4,3)--cycle;
  \node[below] at (2,0){$4$ cm};
  \node[right] at (5.5,2.5){$3$ cm};
  \node[below right] at (4.75,0.5){$3$ cm};
\end{gmFig}
\choices{$10$}{$12$}{$33$}{$66$}{$36$}
\end{probbox}
\probsol{
\textbf{Phân tích đề.} Hình hộp có ba kích thước: dài $4$ cm, cao $3$ cm, sâu (rộng) $3$ cm. Số khối lập phương $1\text{ cm}^3$ cần dùng chính là thể tích hình hộp tính bằng $\text{cm}^3$.
\begin{recallbox}Thể tích hình hộp chữ nhật: $V=a\cdot b\cdot c$ (dài $\times$ rộng $\times$ cao).\end{recallbox}
\step{1}{Đếm số khối ở một lớp}
Ta xếp từng lớp nằm ngang. Lớp dưới cùng có đáy $4\text{ cm}\times3\text{ cm}$ (dài $4$ cm, sâu $3$ cm), nên gồm $3$ hàng, mỗi hàng $4$ khối: $4\cdot3=12$ khối.
\begin{gmFig}[scale=0.55]
  \foreach \z/\c in {0/tmBlue,1/tmGreen,2/tmOrange}{
    \fill[\c!25,draw=\c!80!black] (0,\z) rectangle (4,\z+1);
    \foreach \x in {1,2,3}{\draw[\c!80!black] (\x,\z)--(\x,\z+1);}
    \fill[\c!12,draw=\c!80!black] (4,\z)--(5.5,\z+1)--(5.5,\z+2)--(4,\z+1)--cycle;
    \draw[\c!80!black] (4.5,\z+0.333)--(4.5,\z+1.333); \draw[\c!80!black] (5,\z+0.667)--(5,\z+1.667);
  }
  \fill[tmOrange!8,draw=tmOrange!80!black] (0,3)--(1.5,4)--(5.5,4)--(4,3)--cycle;
  \foreach \x in {1,2,3}{\draw[tmOrange!80!black] (\x,3)--(\x+1.5,4);}
  \draw[tmOrange!80!black] (0.5,3.333)--(4.5,3.333); \draw[tmOrange!80!black] (1,3.667)--(5,3.667);
  \node[anchor=west,font=\footnotesize] at (6.1,0.9){lớp 1: $12$ khối};
  \node[anchor=west,font=\footnotesize] at (6.1,1.9){lớp 2: $12$ khối};
  \node[anchor=west,font=\footnotesize] at (6.1,2.9){lớp 3: $12$ khối};
\end{gmFig}

\step{2}{Đếm số lớp}
Chiều cao là $3$ cm, nên có $3$ lớp chồng lên nhau.

\step{3}{Nhân lại}
Tổng số khối: $12\cdot3=36$ khối. Tức là $V=4\cdot3\cdot3=36\ (\text{cm}^3)$.
\begin{warnbox}Phương án (A) $10$ là \emph{tổng} $4+3+3$, còn (C) $33$ là ghép chữ số. Thể tích là \emph{tích} ba kích thước, không phải tổng.\end{warnbox}
\verify{Đếm theo cách khác: mỗi ``lát'' thẳng đứng song song với mặt trước có $4\cdot3=12$ khối, có $3$ lát theo chiều sâu: $12\cdot3=36$.}
Vậy chọn \textbf{(E)}.
\ans{(E) $36$}
}

\begin{probbox}{11}{\relax}
Một mảnh vườn hình vuông có chu vi $36$ m. Diện tích của mảnh vườn đó, tính bằng $\text{m}^2$, là
\choices{$24$}{$81$}{$36$}{$1296$}{$324$}
\end{probbox}
\probsol{
\textbf{Phân tích đề.} Biết chu vi, muốn tính diện tích. Cầu nối giữa hai đại lượng là \emph{độ dài cạnh}: từ chu vi tìm cạnh, từ cạnh tìm diện tích.
\begin{recallbox}Hình vuông cạnh $a$: chu vi $P=4\cdot a$, diện tích $S=a\cdot a$.\end{recallbox}
\step{1}{Tìm cạnh hình vuông}
Hình vuông có $4$ cạnh bằng nhau, nên $a=P:4=36:4=9$ (m).

\step{2}{Tính diện tích}
$S=9\cdot9=81\ (\text{m}^2)$.
\begin{gmFig}[scale=0.3]
  \fill[tmGreen!15] (0,0) rectangle (9,9);
  \draw[Line] (0,0) grid (9,9);
  \draw[tmGreen!70!black,thick] (0,0) rectangle (9,9);
  \node[below] at (4.5,0){$9$ m}; \node[right] at (9,4.5){$9$ m};
  \node[font=\footnotesize,align=left,anchor=west] at (10.5,4.5){$9$ hàng, mỗi hàng $9$ ô $1\text{ m}^2$:\\ $9\cdot9=81\text{ m}^2$};
\end{gmFig}
\begin{warnbox}
\begin{itemize}[leftmargin=1.3em,itemsep=0pt,topsep=0pt]
  \item (C) $36$: nhầm chu vi với diện tích.
  \item (D) $1296=36\cdot36$: bình phương chu vi thay vì bình phương cạnh.
  \item (E) $324=36\cdot9$: nhân chu vi với cạnh.
\end{itemize}
\end{warnbox}
\verify{Cạnh $9$ m $\Rightarrow$ chu vi $4\cdot9=36$ m, đúng đề bài.}
Vậy chọn \textbf{(B)}.
\ans{(B) $81$}
\begin{tipbox}Với hình vuông, \emph{luôn} tìm cạnh trước. Bảng nhỏ nên nhớ: cạnh $5$ $\to$ chu vi $20$, diện tích $25$; cạnh $9$ $\to$ chu vi $36$, diện tích $81$; cạnh $10$ $\to$ chu vi $40$, diện tích $100$. Khi cạnh tăng gấp đôi, chu vi tăng gấp đôi nhưng diện tích tăng gấp \emph{bốn}.\end{tipbox}
}

\begin{probbox}{12}{\relax}
Hai tam giác cân, mỗi tam giác có ít nhất một góc bằng $70^\circ$. Ở tam giác thứ nhất, số đo (tính bằng độ) của mỗi góc trong hai góc còn lại là một số chẵn. Ở tam giác thứ hai, số đo của mỗi góc trong hai góc còn lại là một số lẻ. Gọi $S$ là tổng hai góc bằng nhau của tam giác thứ nhất, $T$ là tổng hai góc bằng nhau của tam giác thứ hai. Giá trị của $S+T$ là
\choices{$280^\circ$}{$250^\circ$}{$220^\circ$}{$200^\circ$}{$300^\circ$}
\end{probbox}
\probsol{
\textbf{Phân tích đề.} Câu này trông dài nhưng thực chất là bài ``tam giác cân có một góc $70^\circ$'' quen thuộc. Ta liệt kê \emph{mọi} tam giác cân có góc $70^\circ$, rồi dùng điều kiện chẵn/lẻ để biết tam giác nào là thứ nhất, tam giác nào là thứ hai.
\begin{recallbox}Tam giác cân có hai góc ở đáy bằng nhau; tổng ba góc bằng $180^\circ$. Một góc cho trước có thể là góc ở đỉnh hoặc góc ở đáy -- phải xét cả hai.\end{recallbox}
\step{1}{Liệt kê các tam giác cân có một góc $70^\circ$}
\emph{Trường hợp 1: $70^\circ$ là góc ở đỉnh.} Mỗi góc đáy: $(180^\circ-70^\circ):2=110^\circ:2=55^\circ$. Tam giác: $70^\circ,\,55^\circ,\,55^\circ$.

\emph{Trường hợp 2: $70^\circ$ là góc ở đáy.} Góc đáy kia cũng $70^\circ$, góc ở đỉnh: $180^\circ-2\cdot70^\circ=40^\circ$. Tam giác: $70^\circ,\,70^\circ,\,40^\circ$.

Chỉ có hai tam giác như vậy.
\begin{gmFig}[scale=0.75]
  \coordinate (A) at (0,0); \coordinate (B) at (3,0); \coordinate (C) at (1.5,2.142);
  \fill[tmPurple!10] (A)--(B)--(C)--cycle; \draw[Ink,thick] (A)--(B)--(C)--cycle;
  \pic[draw,thick,angle radius=0.45cm,"$55^\circ$",angle eccentricity=1.8]{angle=B--A--C};
  \pic[draw,thick,angle radius=0.45cm,"$55^\circ$",angle eccentricity=1.8]{angle=C--B--A};
  \pic[draw,AnsRed,thick,angle radius=0.45cm,"$70^\circ$",angle eccentricity=1.7]{angle=A--C--B};
  \node[font=\footnotesize] at (1.5,-0.5){$70^\circ$ ở đỉnh};
  \begin{scope}[shift={(5,0)}]
    \coordinate (A) at (0,0); \coordinate (B) at (2,0); \coordinate (C) at (1,2.747);
    \fill[tmTeal!10] (A)--(B)--(C)--cycle; \draw[Ink,thick] (A)--(B)--(C)--cycle;
    \pic[draw,AnsRed,thick,angle radius=0.4cm,"$70^\circ$",angle eccentricity=1.9]{angle=B--A--C};
    \pic[draw,AnsRed,thick,angle radius=0.4cm,"$70^\circ$",angle eccentricity=1.9]{angle=C--B--A};
    \pic[draw,thick,angle radius=0.6cm,"$40^\circ$",angle eccentricity=1.5]{angle=A--C--B};
    \node[font=\footnotesize] at (1,-0.5){$70^\circ$ ở đáy};
  \end{scope}
\end{gmFig}

\step{2}{Xác định tam giác thứ nhất (hai góc còn lại đều chẵn)}
Bỏ đi một góc $70^\circ$:
\begin{itemize}[leftmargin=1.3em,itemsep=0pt,topsep=0pt]
  \item tam giác $70^\circ,55^\circ,55^\circ$ còn lại $55^\circ$ và $55^\circ$ -- đều \emph{lẻ};
  \item tam giác $70^\circ,70^\circ,40^\circ$ còn lại $70^\circ$ và $40^\circ$ -- đều \emph{chẵn}.
\end{itemize}
Vậy tam giác thứ nhất là $70^\circ,\,70^\circ,\,40^\circ$; tam giác thứ hai là $70^\circ,\,55^\circ,\,55^\circ$.

\step{3}{Tính $S$ và $T$}
Tam giác thứ nhất có hai góc bằng nhau là $70^\circ$ và $70^\circ$: $S=70^\circ+70^\circ=140^\circ$.
Tam giác thứ hai có hai góc bằng nhau là $55^\circ$ và $55^\circ$: $T=55^\circ+55^\circ=110^\circ$.

\step{4}{Tính tổng}
$S+T=140^\circ+110^\circ=250^\circ$.
\begin{warnbox}Đọc kĩ: $S$ và $T$ là tổng của \emph{hai góc bằng nhau}, không phải tổng của hai góc ``còn lại'' (vì vậy với tam giác thứ nhất ta lấy $70^\circ+70^\circ$, không lấy $70^\circ+40^\circ$).\end{warnbox}
\verify{$70+70+40=180$ và $70+55+55=180$. Hai tam giác đều hợp lệ; $140+110=250$.}
Vậy chọn \textbf{(B)}.
\ans{(B) $250^\circ$}
}

\begin{probbox}{13}{\relax}
Trên vòng quay may mắn hình tròn tâm $O$ như hình, $PQ$ là một đường kính, bán kính nằm ngang vuông góc với $PQ$ tại $O$, và hai hình quạt $R$, $S$ có diện tích bằng nhau. Khi quay, kim dừng ở phần tô màu trong khoảng bao nhiêu phần trăm số lần?
\begin{gmFig}[scale=0.95]
  \fill[Slate!45] (0,0)--(90:1.7) arc[start angle=90,end angle=45,radius=1.7]--cycle;
  \draw[Ink,thick] (0,0) circle (1.7);
  \draw[Ink,thick] (0,1.7)--(0,-1.7); \draw[Ink,thick] (0,0)--(1.7,0); \draw[Ink,thick] (0,0)--(45:1.7);
  \draw[Ink] (0,0) rectangle (0.2,-0.2);
  \node[above] at (0,1.7){$P$}; \node[below] at (0,-1.7){$Q$}; \node[left] at (0,0){$O$};
  \node at (67.5:1.15){$R$}; \node at (22.5:1.15){$S$};
\end{gmFig}
\choices{$50\%$}{$22{,}5\%$}{$25\%$}{$45\%$}{$12{,}5\%$}
\end{probbox}
\probsol{
\textbf{Phân tích đề.} Kim quay dừng ở mỗi hướng với khả năng như nhau, nên tỉ lệ số lần dừng ở phần tô màu bằng tỉ lệ diện tích (cũng là tỉ lệ góc ở tâm) của phần tô màu so với cả hình tròn. Ta cần tìm góc ở tâm của hình quạt $R$.
\begin{recallbox}Một vòng quanh tâm là $360^\circ$. Hình quạt có góc ở tâm $n^\circ$ chiếm $\dfrac{n}{360}$ diện tích hình tròn.\end{recallbox}
\step{1}{Góc của phần $R$ và $S$ gộp lại}
Bán kính $OP$ hướng lên, bán kính nằm ngang hướng sang phải, hai bán kính này vuông góc. Vậy $R$ và $S$ gộp lại là một phần tư hình tròn, ứng với góc ở tâm $90^\circ$.

\step{2}{Góc của riêng phần $R$}
$R$ và $S$ có diện tích bằng nhau nên góc ở tâm của chúng bằng nhau: mỗi phần $90^\circ:2=45^\circ$.

\step{3}{Đổi ra phần trăm}
\[ \frac{45}{360}=\frac18=0{,}125=12{,}5\%. \]
\begin{proofbox}Vì sao ``diện tích bằng nhau'' kéo theo ``góc bằng nhau''? Diện tích hình quạt bằng $\frac{n}{360}$ diện tích hình tròn, tỉ lệ thuận với góc $n$. Hai hình quạt cùng hình tròn có diện tích bằng nhau thì $\frac{n_1}{360}=\frac{n_2}{360}$, nên $n_1=n_2$.\end{proofbox}
\begin{warnbox}Phương án (C) $25\%$ là phần của \emph{cả} $R$ và $S$; còn (B) $22{,}5\%$ là nhầm lẫn giữa $45^\circ$ với $45\%$ rồi chia đôi. Chỉ phần $R$ được tô màu, và phải chia góc cho $360$.\end{warnbox}
\verify{Phần tô màu là nửa của một phần tư hình tròn, tức $\frac12\cdot\frac14=\frac18=12{,}5\%$.}
Vậy chọn \textbf{(E)}.
\ans{(E) $12{,}5\%$}
\begin{parentbox}Gấp một tờ giấy tròn làm đôi, rồi làm tư, rồi làm tám: mỗi lần gấp, phần nhỏ nhất bị chia đôi. Mở ra, con sẽ thấy $8$ hình quạt bằng nhau, mỗi hình quạt có góc $45^\circ$ và chiếm $12{,}5\%$ -- đúng bằng phần $R$ trong bài.\end{parentbox}
}

\begin{probbox}{14}{\relax}
Trong hình chữ nhật dưới đây (dài $12$ cm, rộng $5$ cm), hai đoạn đánh dấu giống nhau trên cạnh trên đều dài $2$ cm. Diện tích phần tô màu là
\begin{gmFig}[scale=0.45]
  \fill[Slate!70,draw=Ink,thick] (0,0)--(12,0)--(10,5)--(2,5)--cycle;
  \draw[Ink,thick] (0,0) rectangle (12,5);
  \draw[Ink] (0.95,4.75)--(0.95,5.25); \draw[Ink] (1.05,4.75)--(1.05,5.25);
  \draw[Ink] (10.95,4.75)--(10.95,5.25); \draw[Ink] (11.05,4.75)--(11.05,5.25);
  \node[above] at (11,5.2){$2$ cm};
  \node[right] at (12,2.5){$5$ cm};
  \node[below] at (6,0){$12$ cm};
\end{gmFig}
\choices{$60\text{ cm}^2$}{$20\text{ cm}^2$}{$30\text{ cm}^2$}{$40\text{ cm}^2$}{$50\text{ cm}^2$}
\end{probbox}
\probsol{
\textbf{Phân tích đề.} Phần tô màu là một hình thang: đáy lớn là cạnh dưới của hình chữ nhật, đáy nhỏ nằm trên cạnh trên. Hai đoạn ở hai đầu cạnh trên được đánh dấu giống nhau, nghĩa là chúng \emph{bằng nhau}, mỗi đoạn $2$ cm. Ta giải bằng hai cách để tự kiểm tra.

\step{1}{Cách 1 -- hình chữ nhật trừ hai tam giác}
\begin{recallbox}$S_{\text{hình chữ nhật}}=$ dài $\times$ rộng; $\;S_{\text{tam giác vuông}}=\dfrac{\text{tích hai cạnh góc vuông}}{2}$.\end{recallbox}
Diện tích hình chữ nhật: $12\cdot5=60\ (\text{cm}^2)$.

Hai phần không tô màu là hai tam giác vuông ở hai góc; mỗi tam giác có hai cạnh góc vuông là $2$ cm (trên cạnh trên) và $5$ cm (cạnh bên của hình chữ nhật):
\[ S_{\text{một tam giác}}=\frac{2\cdot5}{2}=5\ (\text{cm}^2). \]
Diện tích phần tô màu: $60-5-5=50\ (\text{cm}^2)$.
\begin{gmFig}[scale=0.42]
  \fill[Slate!45] (0,0)--(12,0)--(10,5)--(2,5)--cycle;
  \fill[tmOrange!35] (0,0)--(2,5)--(0,5)--cycle;
  \fill[tmOrange!35] (12,0)--(10,5)--(12,5)--cycle;
  \draw[Ink,thick] (0,0) rectangle (12,5); \draw[Ink,thick] (0,0)--(2,5); \draw[Ink,thick] (12,0)--(10,5);
  \node[font=\footnotesize] at (0.6,3.6){$5$}; \node[font=\footnotesize] at (11.4,3.6){$5$};
  \node[above,font=\footnotesize] at (1,5){$2$}; \node[above,font=\footnotesize] at (11,5){$2$};
  \node[above,font=\footnotesize] at (6,5){$8$};
  \node[font=\footnotesize] at (6,2.3){$60-5-5=50$};
\end{gmFig}

\step{2}{Cách 2 -- công thức hình thang}
\begin{recallbox}$S_{\text{hình thang}}=\dfrac{(a+b)\cdot h}{2}$.\end{recallbox}
Đáy lớn $a=12$ cm; đáy nhỏ $b=12-2-2=8$ cm; chiều cao $h=5$ cm (khoảng cách giữa cạnh trên và cạnh dưới của hình chữ nhật).
\[ S=\frac{(12+8)\cdot5}{2}=\frac{20\cdot5}{2}=\frac{100}{2}=50\ (\text{cm}^2). \]
\begin{warnbox}Phương án (A) $60\text{ cm}^2$ là diện tích \emph{cả} hình chữ nhật. Ngoài ra, nhớ rằng chiều cao của hình thang là $5$ cm (đoạn vuông góc), không phải độ dài cạnh bên xiên.\end{warnbox}
\verify{Hai cách cho cùng kết quả $50\text{ cm}^2$.}
Vậy chọn \textbf{(E)}.
\ans{(E) $50\text{ cm}^2$}
}

\begin{probbox}{15}{\relax}
Tam giác $PQR$ có diện tích $27\text{ cm}^2$ và cạnh đáy $QR=6$ cm. Chiều cao $h$ của tam giác $PQR$ (ứng với đáy $QR$) bằng
\begin{gmFig}[scale=0.8]
  \coordinate (Q) at (0,0); \coordinate (R) at (3,0); \coordinate (P) at (1.6,3.2);
  \fill[tmBlue!10] (P)--(Q)--(R)--cycle;
  \draw[Ink,thick] (P)--(Q)--(R)--cycle;
  \draw[Ink] (1.6,3.2)--(1.6,0); \draw[Ink] (1.6,0) rectangle (1.8,0.2);
  \node[above] at (P){$P$}; \node[below left] at (Q){$Q$}; \node[below right] at (R){$R$};
  \node[right] at (1.6,1.4){$h$}; \node[below=6pt] at (1.5,0){$6$ cm};
\end{gmFig}
\choices{$9$ cm}{$18$ cm}{$4{,}5$ cm}{$2{,}25$ cm}{$7$ cm}
\end{probbox}
\probsol{
\textbf{Phân tích đề.} Biết diện tích và cạnh đáy, cần tìm chiều cao: đây là bài ``tính ngược'' công thức diện tích tam giác.
\begin{recallbox}$S_{\text{tam giác}}=\dfrac{a\cdot h}{2}$, trong đó $a$ là cạnh đáy, $h$ là chiều cao ứng với đáy đó. Suy ra $h=\dfrac{2\cdot S}{a}$.\end{recallbox}
\step{1}{Thay số vào công thức}
$27=\dfrac{6\cdot h}{2}$.

\step{2}{Bỏ phép chia cho $2$}
Nhân hai vế với $2$: $6\cdot h=54$.

\step{3}{Tìm $h$}
$h=54:6=9$ (cm).
\begin{proofbox}Vì sao lại nhân đôi diện tích? Tam giác bằng nửa hình chữ nhật có cùng đáy $6$ cm và cùng chiều cao $h$. Hình chữ nhật đó có diện tích $2\cdot27=54\text{ cm}^2$, và một cạnh $6$ cm, nên cạnh kia (chính là $h$) bằng $54:6=9$ cm.\end{proofbox}
\begin{gmFig}[scale=0.55]
  \fill[tmBlue!10] (0,0) rectangle (6,4.5);
  \fill[tmBlue!35] (0,0)--(6,0)--(3.2,4.5)--cycle;
  \draw[Ink,thick] (0,0) rectangle (6,4.5); \draw[Ink,thick] (0,0)--(3.2,4.5)--(6,0);
  \draw[Ink,dashed] (3.2,4.5)--(3.2,0);
  \node[below] at (3,0){$6$ cm}; \node[right] at (6,2.25){$h=9$ cm};
  \node[font=\footnotesize,align=left,anchor=west] at (8.3,2.25){hình chữ nhật: $54\text{ cm}^2$\\ tam giác: $27\text{ cm}^2$};
\end{gmFig}
\begin{warnbox}Phương án (C) $4{,}5$ cm là kết quả của $27:6$ -- quên rằng diện tích tam giác chỉ bằng \emph{một nửa} tích đáy và chiều cao.\end{warnbox}
\verify{$\dfrac{6\cdot9}{2}=\dfrac{54}{2}=27\ (\text{cm}^2)$. Đúng.}
Vậy chọn \textbf{(A)}.
\ans{(A) $9$ cm}
}

% ---------------- part b ----------------
% ===== Topic 2, part b: problems 16-30 (VI) =====
% helper: 2x2 grid with one letter; #1,#2 = centre of the cell, #3 = letter, #4 = node transform
\providecommand{\tbGrid}[4]{\begin{tikzpicture}[baseline=0.2cm,x=0.5cm,y=0.5cm]
  \draw[Ink,line width=0.7pt] (0,0) grid[step=1] (2,2);
  \node[font=\sffamily\bfseries,text=tmBlue!80!black,inner sep=0pt,#4] at (#1,#2) {#3};\end{tikzpicture}}

% ---------------------------------------------------------------- 16
\begin{probbox}{16}{\relax}
Một mảnh vườn hình chữ nhật có chiều dài hơn \emph{hai lần} chiều rộng là $6$ m.
Chu vi của mảnh vườn là $120$ m. Hỏi chiều rộng của mảnh vườn bằng bao nhiêu mét?
\choices{$8$ m}{$18$ m}{$27$ m}{$38$ m}{$22$ m}
\end{probbox}
\probsol{
\textbf{Phân tích đề.} Ta biết hai điều: (1) chiều dài $=$ hai lần chiều rộng, cộng thêm $6$ m;
(2) chu vi $=120$ m. Cần tìm chiều rộng.

\step{1}{Gọi chiều rộng là $r$ (m) và biểu diễn chiều dài theo $r$}
``Hai lần chiều rộng'' là $2r$; ``hơn hai lần chiều rộng $6$ m'' nên chiều dài là $2r+6$ (m).
Ta vẽ sơ đồ đoạn thẳng:
\begin{barfigC}
\rowlab{-0.15}{0}{Dài}
\blk{0}{2}{0}{0.6}{tmBlue}{$r$}\blk{2}{4}{0}{0.6}{tmBlue}{$r$}\blk{4}{5.1}{0}{0.6}{tmOrange}{$6$}
\rowlab{-0.15}{-0.9}{Rộng}
\blk{0}{2}{-0.9}{0.6}{tmBlue}{$r$}
\vbrace{5.3}{0.3}{-1.2}{nửa chu vi $=60$ m}
\end{barfigC}

\begin{recallbox}Chu vi hình chữ nhật $=(\text{dài}+\text{rộng})\times 2$. Do đó
$\text{dài}+\text{rộng}=$ chu vi $:2$ (gọi là \emph{nửa chu vi}).\end{recallbox}

\step{2}{Tính nửa chu vi}
$120:2=60$ (m). Vậy chiều dài $+$ chiều rộng $=60$ m.

\step{3}{Lập và giải phương trình}
Nhìn sơ đồ: dài $+$ rộng gồm $3$ đoạn $r$ và một đoạn $6$ m:
\[ (2r+6)+r=60 \;\Rightarrow\; 3r+6=60 \;\Rightarrow\; 3r=60-6=54 \;\Rightarrow\; r=54:3=18. \]
Chiều rộng là $18$ m, chiều dài là $2\cdot 18+6=42$ (m).

\begin{warnbox}Nhiều bạn quên chia đôi chu vi và viết $3r+6=120$, được $r=38$ -- đó chính là
phương án (D) ``bẫy''. Hãy nhớ: chu vi gồm \emph{hai} chiều dài và \emph{hai} chiều rộng.\end{warnbox}

\verify{Chiều dài $42$ m, chiều rộng $18$ m: $42=2\cdot 18+6$ (đúng) và chu vi
$(42+18)\cdot 2=60\cdot 2=120$ m (đúng).}
Chiều rộng $18$ m ứng với phương án \textbf{(B)}.
\ans{(B) $18$ m}
}

% ---------------------------------------------------------------- 17
\begin{probbox}{17}{\relax}
Chữ P được viết trong một lưới ô vuông $2\times 2$ như hình bên: \tbGrid{1.5}{1.5}{P}{}\,.
Người ta thực hiện liên tiếp một số phép quay quanh tâm của lưới và một số phép lật (đối xứng)
qua hai đường thẳng đi qua tâm lưới (đường nằm ngang và đường thẳng đứng), và nhận được kết quả:
\tbGrid{0.5}{0.5}{P}{cm={0,-1,-1,0,(0,0)}}\,.

Nếu thực hiện \emph{đúng dãy phép biến hình đó} với lưới \tbGrid{0.5}{1.5}{A}{} thì ta được hình nào?
\par\smallskip\noindent
\begin{tabularx}{\linewidth}{@{}XXXXX@{}}
(A) \tbGrid{0.5}{1.5}{A}{} & (B) \tbGrid{0.5}{1.5}{A}{rotate=90} & (C) \tbGrid{1.5}{0.5}{A}{rotate=-90} &
(D) \tbGrid{1.5}{0.5}{A}{} & (E) \tbGrid{1.5}{0.5}{A}{rotate=180}
\end{tabularx}
\end{probbox}
\probsol{
\textbf{Phân tích đề.} Ta không biết dãy phép biến hình cụ thể gồm những bước nào, nhưng biết
\emph{kết quả cuối cùng} của nó trên chữ P. Ý tưởng: tưởng tượng lưới là một \emph{tấm thẻ} vuông
bằng giấy kính trong, có vẽ chữ trên đó. Mỗi phép quay hay lật chỉ là xoay hoặc lật tấm thẻ rồi đặt lại
đúng chỗ cũ. Cả dãy phép biến hình cũng chỉ là \emph{một cách đặt lại tấm thẻ}. Ta đọc cách đặt đó từ
hình chữ P, rồi áp dụng cho chữ A.

\step{1}{Quan sát chữ P đã đi đâu}
Chữ P ở ô \emph{trên -- phải} đã chuyển sang ô \emph{dưới -- trái}. Đồng thời, nét thẳng (nét ``thân'')
của chữ P, vốn thẳng đứng, nay nằm ngang ở phía trên; phần ``bụng'' tròn vốn ở bên phải thân,
nay nằm ở phía dưới.

\step{2}{Tìm xem tấm thẻ đã được đặt lại như thế nào}
Có hai cách đặt lại tấm thẻ đưa ô trên -- phải về ô dưới -- trái:
\begin{itemize}[leftmargin=1.4em,itemsep=1pt]
\item quay $180^\circ$: khi đó chữ P sẽ bị lộn ngược (thân vẫn \emph{thẳng đứng}, trông giống chữ ``d'');
\item lật qua đường chéo nối góc trên -- trái với góc dưới -- phải (đường nét đứt trong hình dưới):
khi đó thân thẳng đứng của P sẽ thành \emph{nằm ngang}.
\end{itemize}
Hình đề bài cho thân nằm ngang, nên dãy phép biến hình có tác dụng như
\textbf{phép lật qua đường chéo trên-trái -- dưới-phải}.
\begin{figC}\begin{scope}[x=0.8cm,y=0.8cm]
\draw[Ink,line width=0.8pt] (0,0) grid[step=1] (2,2);
\draw[tmOrange,dashed,line width=1pt] (-0.3,2.3)--(2.3,-0.3);
\node[font=\sffamily\bfseries\large,text=tmBlue] at (1.5,1.5){P};
\node[font=\sffamily\bfseries\large,text=tmGreen!80!black,cm={0,-1,-1,0,(0,0)}] at (0.5,0.5){P};
\draw[->,tmPurple,thick] (1.55,1.15) to[bend left=35] (0.85,0.45);
\node[font=\footnotesize,text=tmOrange!80!black,anchor=west] at (2.3,-0.3){trục lật};
\begin{scope}[xshift=4.2cm]
\draw[Ink,line width=0.8pt] (0,0) grid[step=1] (2,2);
\draw[tmOrange,dashed,line width=1pt] (-0.3,2.3)--(2.3,-0.3);
\node[font=\sffamily\bfseries\large,text=tmBlue] at (0.5,1.5){A};
\node[font=\footnotesize,text=Slate] at (1,-0.45){A nằm ngay trên trục};
\end{scope}
\end{scope}\end{figC}
\begin{recallbox}Một dãy phép quay quanh tâm và phép lật qua đường ngang, đường dọc có thể cho kết quả giống
một phép lật qua đường chéo. Chẳng hạn: \emph{quay $90^\circ$ theo chiều kim đồng hồ rồi lật qua đường
thẳng đứng}. Thử với P: quay $90^\circ$ đưa P từ ô trên -- phải xuống ô dưới -- phải, lật qua đường
dọc đưa P sang ô dưới -- trái, và thân P nằm ngang -- đúng như hình đề bài.\end{recallbox}

\step{3}{Áp dụng cho chữ A}
Ô trên -- trái nằm \emph{ngay trên} trục lật (đường chéo đi qua nó), nên chữ A \emph{vẫn ở ô trên -- trái}.
Phép lật qua đường chéo này biến hướng ``lên trên'' thành hướng ``sang trái''. Đỉnh nhọn của chữ A vốn
chỉ lên trên, nay chỉ sang trái. Chữ A đối xứng trái -- phải nên không cần lo chuyện bị ``soi gương''.

\begin{warnbox}Đừng chỉ nhìn vị trí ô! Chữ P đi từ ô trên -- phải sang ô dưới -- trái, nhưng chữ A
\emph{không} vì thế mà cũng ``đi chéo xuống''. Phải tìm phép biến hình tác động lên \emph{cả tấm thẻ},
rồi mới áp dụng.\end{warnbox}

\verify{Dùng dãy cụ thể ``quay $90^\circ$ theo chiều kim đồng hồ rồi lật qua đường dọc'': quay đưa
A từ ô trên -- trái sang ô trên -- phải, đỉnh A chỉ sang phải; lật qua đường dọc đưa A về ô trên -- trái,
đỉnh A chỉ sang trái. Đúng là hình (B).}
\ans{(B) chữ A ở ô trên -- trái, đỉnh chỉ sang trái}
}

% ---------------------------------------------------------------- 18
\begin{probbox}{18}{\relax}
Trong hình vuông cạnh $8$ cm, người ta vẽ hai đường chéo. Diện tích phần tô màu bằng bao nhiêu?
\begin{figC}\begin{scope}[x=0.38cm,y=0.38cm]
\fill[tmBlue!35] (0,0)--(8,0)--(4,4)--cycle;
\draw[Ink,line width=1pt] (0,0) rectangle (8,8);
\draw[Ink,line width=0.9pt] (0,0)--(8,8) (0,8)--(8,0);
\node[left] at (0,4){$8$ cm};
\end{scope}\end{figC}
\choices{$4$ cm$^2$}{$8$ cm$^2$}{$16$ cm$^2$}{$56$ cm$^2$}{$64$ cm$^2$}
\end{probbox}
\probsol{
\textbf{Phân tích đề.} Hai đường chéo chia hình vuông thành $4$ tam giác. Phần tô màu là một trong
bốn tam giác đó. Cần tìm diện tích của nó.

\step{1}{Tính diện tích hình vuông}
\begin{recallbox}Diện tích hình vuông cạnh $a$ là $a\times a$.\end{recallbox}
$S_{\text{vuông}}=8\times 8=64$ (cm$^2$).

\step{2}{Chứng tỏ bốn tam giác có diện tích bằng nhau}
\begin{proofbox}Hai đường chéo của hình vuông bằng nhau và cắt nhau tại trung điểm $O$ của mỗi đường, nên
$O$ cách đều bốn đỉnh. Quay hình vuông $90^\circ$ quanh $O$ thì mỗi tam giác chồng khít lên tam giác
bên cạnh. Vậy bốn tam giác bằng nhau, mỗi tam giác chiếm $\tfrac14$ hình vuông.\end{proofbox}

\step{3}{Tính diện tích phần tô màu}
$S=64:4=16$ (cm$^2$).

\verify{Tính trực tiếp: tam giác tô màu có đáy là cạnh dưới dài $8$ cm, chiều cao là khoảng cách từ
tâm $O$ xuống cạnh đáy, bằng nửa cạnh, tức $4$ cm. Diện tích $=\dfrac{8\times 4}{2}=16$ cm$^2$. Khớp.}
Phương án \textbf{(C)}.
\ans{(C) $16$ cm$^2$}
}

% ---------------------------------------------------------------- 19
\begin{probbox}{19}{\relax}
Lưới ô vuông dưới đây gồm các ô vuông nhỏ bằng nhau. Phần tô màu chiếm bao nhiêu phần trăm diện tích
cả lưới?
\begin{figC}\begin{scope}[x=0.36cm,y=0.36cm]
\foreach \r in {0,...,7}{\foreach \c in {0,...,7}{\fill[tmBlue!40] (\c,\r) rectangle (\c+1,\r+1);}}
\foreach \c/\r in {2/7,5/7,2/6,5/6,0/5,1/5,6/5,7/5,0/2,1/2,6/2,7/2,2/1,5/1,2/0,5/0}{\fill[white] (\c,\r) rectangle (\c+1,\r+1);}
\draw[Ink,line width=0.7pt] (0,0) grid[step=1] (8,8);
\end{scope}\end{figC}
\choices{$84\%$}{$80\%$}{$90\%$}{$75\%$}{$66\%$}
\end{probbox}
\probsol{
\textbf{Phân tích đề.} Muốn tính phần trăm, ta cần biết: có tất cả bao nhiêu ô, và có bao nhiêu ô tô màu.
Đếm ô \emph{trắng} sẽ nhanh hơn vì chúng ít hơn.

\step{1}{Đếm tổng số ô}
Lưới có $8$ hàng, mỗi hàng $8$ ô: $8\times 8=64$ ô.

\step{2}{Đếm số ô trắng}
Các ô trắng xếp thành $8$ ``cặp'' hai ô liền nhau, nằm đối xứng ở bốn phía:
\begin{itemize}[leftmargin=1.4em,itemsep=1pt]
\item hai hàng trên cùng: cột thứ $3$ và cột thứ $6$ có ô trắng ở cả hai hàng: $4$ ô;
\item hai hàng dưới cùng: tương tự, $4$ ô;
\item hàng thứ $3$ (từ trên xuống): $2$ ô trắng sát mép trái và $2$ ô sát mép phải: $4$ ô;
\item hàng thứ $6$: tương tự, $4$ ô.
\end{itemize}
Tổng số ô trắng: $4+4+4+4=16$ ô.

\step{3}{Tính số ô tô màu và đổi ra phần trăm}
Số ô tô màu: $64-16=48$ ô.
\begin{recallbox}Tỉ số phần trăm của $a$ so với $b$ là $\dfrac{a}{b}\times 100\%$.\end{recallbox}
\[\frac{48}{64}=\frac{3}{4}=\frac{75}{100}=75\%.\]

\begin{tipbox}Có thể không cần đếm hết: ô trắng chiếm $\dfrac{16}{64}=\dfrac14=25\%$, nên ô tô màu chiếm
$100\%-25\%=75\%$.\end{tipbox}

\verify{$75\%$ của $64$ là $64\times\dfrac34=48$ ô, và $48+16=64$. Đúng.}
Phương án \textbf{(D)}.
\ans{(D) $75\%$}
}

% ---------------------------------------------------------------- 20
\begin{probbox}{20}{\relax}
Cho $JKLM$ là hình vuông và $PQRS$ là hình chữ nhật, đặt chồng lên nhau như hình vẽ. Biết $JK$ song song
với $PQ$, $JK=8$ và $PS=2$. Tính tổng diện tích các phần tô màu.
\begin{figC}\begin{scope}[x=0.3cm,y=0.3cm]
\fill[tmBlue!35] (5,-3) rectangle (13,5);
\draw[Ink,line width=0.9pt] (5,-3) rectangle (13,5);
\fill[white] (0,0) rectangle (18,2);
\draw[Ink,line width=0.9pt] (0,0) rectangle (18,2);
\node[above left] at (0,2){$P$};\node[above right] at (18,2){$Q$};
\node[below right] at (18,0){$R$};\node[below left] at (0,0){$S$};
\node[above] at (5,5){$J$};\node[above] at (13,5){$K$};
\node[below] at (13,-3){$L$};\node[below] at (5,-3){$M$};
\end{scope}\end{figC}
\choices{$32$}{$16$}{$56$}{$48$}{$62$}
\end{probbox}
\probsol{
\textbf{Phân tích đề.} Phần tô màu là phần của hình vuông \emph{không} bị dải hình chữ nhật che.
Vậy: diện tích tô màu $=$ diện tích hình vuông $-$ diện tích phần chung của hai hình.

\step{1}{Tính diện tích hình vuông $JKLM$}
\begin{recallbox}Diện tích hình vuông $=$ cạnh $\times$ cạnh; diện tích hình chữ nhật $=$ dài $\times$ rộng.\end{recallbox}
$S_{JKLM}=8\times 8=64$.

\step{2}{Xác định phần chung}
Vì $JK\parallel PQ$ nên các cạnh của dải $PQRS$ song song với cạnh $JK$ của hình vuông, tức là dải
nằm ``ngang'' cắt qua hình vuông. Dải dài hơn hình vuông nên cắt ngang trọn bề rộng hình vuông.
Phần chung là một hình chữ nhật có:
\begin{itemize}[leftmargin=1.4em,itemsep=1pt]
\item chiều dài bằng cạnh hình vuông: $8$;
\item chiều rộng bằng bề rộng của dải: $PS=2$.
\end{itemize}
Diện tích phần chung: $8\times 2=16$.

\step{3}{Tính diện tích tô màu}
$64-16=48$.

\begin{warnbox}Không cần biết dải nằm cao hay thấp, và hai mảnh tô màu to nhỏ ra sao: tổng của chúng
luôn là $64-16=48$. Nhiều bạn chỉ tính $8\times 2=16$ (phương án B) -- đó là phần \emph{bị che}, không phải
phần tô màu.\end{warnbox}

\verify{Giả sử dải cách cạnh trên $JK$ một khoảng $3$; mảnh trên có diện tích $8\times 3=24$, mảnh dưới có
chiều cao $8-3-2=3$, diện tích $24$. Tổng $48$. Đúng.}
Phương án \textbf{(D)}.
\ans{(D) $48$}
}

% ---------------------------------------------------------------- 21
\begin{probbox}{21}{\relax}
Ta dựng được tam giác có ba cạnh dài $4$, $5$, $8$, nhưng \emph{không} dựng được tam giác có ba cạnh
$4$, $5$, $10$. Chỉ dùng các độ dài $2$, $3$, $5$, $7$, $11$ (mỗi độ dài có thể được dùng nhiều lần), có thể dựng
được bao nhiêu tam giác khác nhau có \emph{đúng hai cạnh bằng nhau}?
\choices{$8$}{$5$}{$20$}{$10$}{$14$}
\end{probbox}
\probsol{
\textbf{Phân tích đề.} Tam giác có đúng hai cạnh bằng nhau là tam giác cân nhưng không đều: hai cạnh bên
bằng $a$, cạnh đáy bằng $b$ với $b\neq a$. Ta phải đếm các cặp $(a,b)$ lấy từ $\{2,3,5,7,11\}$ sao cho ba
đoạn $a,a,b$ thật sự dựng được tam giác.

\begin{recallbox}\textbf{Bất đẳng thức tam giác.} Ba đoạn thẳng là ba cạnh của một tam giác khi và chỉ khi
tổng hai đoạn bất kì lớn hơn đoạn còn lại. Thực hành: chỉ cần kiểm tra
\emph{tổng hai đoạn ngắn hơn lớn hơn đoạn dài nhất}. Ví dụ: $4+5=9>8$ nên dựng được; $4+5=9<10$ nên không
dựng được.\end{recallbox}

\step{1}{Rút gọn điều kiện cho tam giác cân $a,a,b$}
Các tổng cần kiểm tra là $a+a>b$ và $a+b>a$. Điều thứ hai luôn đúng (vì $b>0$). Vậy chỉ cần
\[ 2a>b \quad\text{(hai lần cạnh bên lớn hơn cạnh đáy).}\]

\step{2}{Lập bảng theo cạnh bên $a$}
\begin{center}
\begin{tabular}{|c|c|c|c|}
\hline
Cạnh bên $a$ & $2a$ & Cạnh đáy $b\ne a$ thỏa $b<2a$ & Số tam giác\\ \hline
$2$ & $4$ & $3$ & $1$\\
$3$ & $6$ & $2,\ 5$ & $2$\\
$5$ & $10$ & $2,\ 3,\ 7$ & $3$\\
$7$ & $14$ & $2,\ 3,\ 5,\ 11$ & $4$\\
$11$ & $22$ & $2,\ 3,\ 5,\ 7$ & $4$\\ \hline
\end{tabular}
\end{center}
Chẳng hạn với $a=3$: đáy $b=7$ không được vì $3+3=6<7$; $b=11$ cũng không được.
Với $a=5$: $b=11$ không được vì $5+5=10<11$.

\step{3}{Cộng lại}
$1+2+3+4+4=14$ tam giác.

\begin{warnbox}Tam giác $2,2,5$ hay $3,3,7$ ``trông có vẻ'' được nhưng thật ra hai cạnh bên quá ngắn,
không chạm nhau được. Ngoài ra, tam giác $5,5,5$ có \emph{ba} cạnh bằng nhau nên không được đếm.\end{warnbox}

\verify{Đếm ngược: có $5\times 4=20$ cặp $(a,b)$ với $b\neq a$. Các cặp bị loại (vì $b\ge 2a$) là
$(2;5),(2;7),(2;11),(3;7),(3;11),(5;11)$: $6$ cặp. $20-6=14$. Khớp.}
Phương án \textbf{(E)}.
\ans{(E) $14$}
}

% ---------------------------------------------------------------- 22
\begin{probbox}{22}{\relax}
Hình bình hành $PQRS$ có diện tích bằng $40$. Gọi $T$ và $V$ lần lượt là trung điểm của các cạnh $PS$
và $RS$. Tính diện tích tứ giác $PRVT$.
\begin{figC}\begin{scope}[x=0.55cm,y=0.55cm]
\fill[tmBlue!25] (0,0)--(8,4)--(7,2)--(3,0)--cycle;
\draw[Ink,line width=1pt] (0,0)--(2,4)--(8,4)--(6,0)--cycle;
\draw[Ink,dashed,line width=0.8pt] (0,0)--(8,4) (3,0)--(7,2);
\foreach \x in {1.5,4.5}{\draw[Ink] (\x,-0.15)--(\x,0.15);}
\foreach \x/\y in {6.5/1,7.5/3}{\draw[Ink] (\x-0.15,\y+0.07)--(\x+0.15,\y-0.07);\draw[Ink] (\x-0.15,\y+0.17)--(\x+0.15,\y+0.03);}
\node[below left] at (0,0){$P$};\node[above left] at (2,4){$Q$};
\node[above right] at (8,4){$R$};\node[below right] at (6,0){$S$};
\node[below] at (3,0){$T$};\node[right] at (7,2){$V$};
\end{scope}\end{figC}
\choices{$10$}{$12$}{$15$}{$16$}{$18$}
\end{probbox}
\probsol{
\textbf{Phân tích đề.} Tứ giác $PRVT$ là phần còn lại khi lấy tam giác $PRS$ bỏ đi tam giác nhỏ $TVS$
ở góc $S$. Vậy ta tính $S_{PRS}$ và $S_{TVS}$.

\step{1}{Diện tích tam giác $PRS$}
\begin{recallbox}Đường chéo chia hình bình hành thành hai tam giác bằng nhau, nên mỗi tam giác có diện tích
bằng nửa diện tích hình bình hành.\end{recallbox}
$S_{PRS}=40:2=20$.

\step{2}{Diện tích tam giác $TRS$ (nối thêm $T$ với $R$)}
\begin{recallbox}Hai tam giác có \emph{chung chiều cao} thì tỉ số diện tích bằng tỉ số hai đáy. (Vì diện tích
tam giác $=\tfrac12\times$ đáy $\times$ chiều cao.)\end{recallbox}
Tam giác $TRS$ và tam giác $PRS$ có chung đỉnh $R$, hai đáy $TS$ và $PS$ nằm trên cùng một đường thẳng,
nên có chung chiều cao hạ từ $R$. Vì $T$ là trung điểm $PS$ nên $TS=\tfrac12 PS$. Do đó
$S_{TRS}=\tfrac12\cdot 20=10$.

\step{3}{Diện tích tam giác $TVS$}
Tam giác $TVS$ và tam giác $TRS$ có chung đỉnh $T$, đáy $VS$ và $RS$ cùng nằm trên đường thẳng $RS$.
Vì $V$ là trung điểm $RS$ nên $VS=\tfrac12 RS$, suy ra $S_{TVS}=\tfrac12\cdot 10=5$.

\step{4}{Diện tích tứ giác $PRVT$}
$S_{PRVT}=S_{PRS}-S_{TVS}=20-5=15$.

\begin{tipbox}Tam giác ``nửa -- nửa'' ở góc có diện tích bằng $\tfrac12\times\tfrac12=\tfrac14$ tam giác lớn.
Đây là một kết quả rất hay dùng!\end{tipbox}

\verify{Đặt hình lên lưới: $P(0;0)$, $S(6;0)$, $R(8;4)$, $Q(2;4)$ (diện tích $6\times4=24$, không phải $40$,
nhưng tỉ lệ thì như nhau). Khi đó $T(3;0)$, $V(7;2)$; $S_{PRS}=12$, $S_{TVS}=\tfrac12\cdot 3\cdot 2=3$,
nên $S_{PRVT}=9=\tfrac{3}{8}\cdot 24$. Với hình đề bài: $\tfrac38\cdot 40=15$. Khớp.}
Phương án \textbf{(C)}.
\ans{(C) $15$}
}

% ---------------------------------------------------------------- 23
\begin{probbox}{23}{\stars{1}}
Trên lưới ô vuông $10\times 10$ (mỗi ô có cạnh $1$), cho tam giác $FGH$ như hình vẽ. Điểm $P$ có thể là
một trong $41$ giao điểm của các đường kẻ lưới nằm \emph{bên trong} (không nằm trên cạnh) tam giác $FGH$.
Với mỗi vị trí của $P$ ta được đúng ba tam giác: $FPG$, $GPH$, $HPF$. Trong tất cả $123$ tam giác đó, có
bao nhiêu tam giác có diện tích \emph{đúng bằng một nửa} diện tích tam giác $FGH$?
\begin{figC}\begin{scope}[x=0.36cm,y=0.36cm]
\draw[Line,line width=0.5pt] (0,0) grid[step=1] (10,10);
\draw[Ink,line width=1.3pt] (0,0)--(10,0)--(5,10)--cycle;
\node[above] at (5,10){$F$};\node[below left] at (0,0){$G$};\node[below right] at (10,0){$H$};
\end{scope}\end{figC}
\choices{$5$}{$7$}{$3$}{$11$}{$9$}
\end{probbox}
\probsol{
\textbf{Phân tích đề.} Đặt hệ trục: $G(0;0)$, $H(10;0)$, $F(5;10)$. Tam giác $FGH$ có đáy $GH=10$ và chiều
cao $10$, nên
\[ S_{FGH}=\frac{10\times 10}{2}=50,\qquad \text{một nửa là } 25. \]
Ta đếm, lần lượt cho từng loại tam giác $GPH$, $FPG$, $HPF$, những vị trí của $P$ cho diện tích $25$.

\begin{recallbox}Hai tam giác có \emph{chung đáy} thì tỉ số diện tích bằng tỉ số hai chiều cao.
Các điểm cách một đường thẳng cùng một khoảng (về cùng một phía) nằm trên một đường thẳng song song
với đường thẳng đó.\end{recallbox}

\step{1}{Tam giác $GPH$}
Tam giác $GPH$ có chung đáy $GH$ với tam giác $FGH$. Muốn diện tích bằng một nửa thì chiều cao hạ từ $P$
phải bằng một nửa chiều cao $10$, tức là $5$. Vậy $P$ nằm trên đường kẻ ngang thứ $5$ (tính từ $GH$).
Trên đường này, cạnh $FG$ đi qua điểm có hoành độ $2{,}5$ và cạnh $FH$ đi qua điểm có hoành độ $7{,}5$
(trung điểm của hai cạnh). Các giao điểm lưới nằm \emph{bên trong} có hoành độ $3,4,5,6,7$:
\[ (3;5),\ (4;5),\ (5;5),\ (6;5),\ (7;5)\quad\Rightarrow\quad 5 \text{ tam giác.} \]

\step{2}{Tam giác $FPG$}
Lần này đáy chung là $FG$. Gọi $M(5;0)$ là trung điểm của $GH$ và $N(7{,}5;5)$ là trung điểm của $FH$.
\begin{proofbox}Tam giác $FMG$ và $FHG$ có chung chiều cao hạ từ $F$, còn đáy $GM=\tfrac12 GH$, nên
$S_{FMG}=\tfrac12 S_{FHG}=25$. Như vậy điểm $M$ cách $FG$ một khoảng đúng bằng nửa khoảng cách từ $H$
đến $FG$. Tương tự với $N$ (trung điểm của $FH$). Do đó $S_{FPG}=25$ khi và chỉ khi $P$ cách $FG$ đúng bằng
$M$, tức là $P$ nằm trên đường thẳng $MN$ (song song với $FG$).\end{proofbox}
Đi từ $G$ đến $F$ ta sang phải $5$, lên trên $10$, tức là ``sang phải $1$, lên $2$''. Đường $MN$ song song với
$FG$ nên từ $M(5;0)$ ta cũng đi ``sang phải $1$, lên $2$'': gặp $(6;2)$, $(7;4)$, rồi tới $N(7{,}5;5)$ (không
phải giao điểm lưới). Hai điểm $(6;2)$, $(7;4)$ nằm bên trong tam giác $\Rightarrow$ $2$ tam giác.

\step{3}{Tam giác $HPF$}
Hình vẽ đối xứng qua đường thẳng đứng $x=5$ nên kết quả như bước 2: $P=(4;2)$ hoặc $(3;4)$
$\Rightarrow$ $2$ tam giác.

\step{4}{Cộng lại}
$5+2+2=9$ tam giác.

\verify{Kiểm tra $P(6;2)$ bằng cách ``bao hình chữ nhật'': hình chữ nhật từ $x=0$ đến $6$, $y=0$ đến $10$ có
diện tích $60$. Bỏ đi ba tam giác vuông ở ngoài tam giác $FPG$: $\tfrac{5\cdot 10}{2}=25$, $\tfrac{1\cdot 8}{2}=4$,
$\tfrac{6\cdot 2}{2}=6$. Còn lại $60-35=25$. Đúng bằng nửa của $50$.}
\begin{warnbox}Đừng quên hai loại tam giác $FPG$ và $HPF$! Nếu chỉ xét tam giác có đáy nằm ngang $GH$
thì sẽ chọn nhầm phương án (A) $5$.\end{warnbox}
Phương án \textbf{(E)}.
\ans{(E) $9$}
}

% ---------------------------------------------------------------- 24
\begin{probbox}{24}{\stars{1}}
Các điểm $T,U,V,W,X,Y$ nằm trên các cạnh của hình vuông $PQRS$ như hình vẽ, với
\[PT=TU=UQ=QV=VW=WR=XS=SY.\]
Phần tô màu chiếm bao nhiêu phần diện tích hình vuông $PQRS$?
\begin{figC}\begin{scope}[x=1.25cm,y=1.25cm]
\fill[tmBlue!35] (1,3)--(2,3)--(1.5,2.5)--(2.5,1.5)--(3,2)--(3,1)--(1,1)--cycle;
\draw[Ink,line width=1pt] (0,0) rectangle (3,3);
\draw[Ink,line width=0.9pt] (1,3)--(1,0) (0,1)--(2,3) (1,3)--(3,1) (1,0)--(3,2) (0,1)--(3,1);
\foreach \x in {0.5,1.5,2.5}{\draw[Ink] (\x-0.03,2.92)--(\x-0.03,3.08) (\x+0.03,2.92)--(\x+0.03,3.08);}
\foreach \y in {0.5,1.5,2.5}{\draw[Ink] (2.92,\y-0.03)--(3.08,\y-0.03) (2.92,\y+0.03)--(3.08,\y+0.03);}
\draw[Ink] (0.47,-0.08)--(0.47,0.08) (0.53,-0.08)--(0.53,0.08);
\draw[Ink] (-0.08,0.47)--(0.08,0.47) (-0.08,0.53)--(0.08,0.53);
\node[above left] at (0,3){$P$};\node[above right] at (3,3){$Q$};\node[below right] at (3,0){$R$};
\node[below left] at (0,0){$S$};\node[above] at (1,3){$T$};\node[above] at (2,3){$U$};
\node[right] at (3,2){$V$};\node[right] at (3,1){$W$};\node[below] at (1,0){$X$};\node[left] at (0,1){$Y$};
\end{scope}\end{figC}
\par\smallskip\noindent
\begin{tabularx}{\linewidth}{@{}XXXXX@{}}(A)\ $\dfrac{5}{18}$&(B)\ $\dfrac13$&(C)\ $\dfrac29$&(D)\ $\dfrac14$&(E)\ $\dfrac16$\end{tabularx}
\end{probbox}
\probsol{
\textbf{Phân tích đề.} Mỗi cạnh hình vuông được chia thành $3$ đoạn bằng nhau. Ta chọn mỗi đoạn dài $1$ đơn
vị, khi đó hình vuông có cạnh $3$ và diện tích $3\times 3=9$. Ý hay: kẻ thêm các đường chia hình vuông thành
lưới $3\times 3$ ô vuông đơn vị (nét chấm), mọi đường trong hình sẽ hiện ra rất ``ngay ngắn''.
\begin{figC}\begin{scope}[x=1.25cm,y=1.25cm]
\fill[tmBlue!12] (1,3)--(3,3)--(3,1)--(1,1)--cycle;
\fill[tmBlue!40] (1,3)--(3,1)--(1,1)--cycle;
\fill[tmOrange!50] (1,3)--(2,3)--(1.5,2.5)--cycle;
\fill[tmGreen!50] (3,2)--(3,1)--(2.5,1.5)--cycle;
\draw[Slate,densely dotted] (0,0) grid[step=1] (3,3);
\draw[Ink,line width=1pt] (0,0) rectangle (3,3);
\draw[Ink,line width=0.9pt] (1,3)--(1,0) (0,1)--(2,3) (1,3)--(3,1) (1,0)--(3,2) (0,1)--(3,1);
\draw[tmOrange!80!black,dashed] (1,2)--(2,3) (1,3)--(2,2);
\draw[tmGreen!70!black,dashed] (2,2)--(3,1) (2,1)--(3,2);
\node[above left] at (0,3){$P$};\node[above right] at (3,3){$Q$};\node[below right] at (3,0){$R$};
\node[below left] at (0,0){$S$};\node[above] at (1,3){$T$};\node[above] at (2,3){$U$};
\node[right] at (3,2){$V$};\node[right] at (3,1){$W$};\node[below] at (1,0){$X$};\node[left] at (0,1){$Y$};
\node[below left] at (1,1){$Z$};\node[font=\footnotesize] at (1.6,1.45){$2$};
\end{scope}\end{figC}

\step{1}{Nhận ra ba mảnh của phần tô màu}
Gọi $Z$ là giao điểm của $TX$ và $YW$. Phần tô màu gồm ba mảnh không chồng lên nhau:
tam giác lớn $TZW$ (xanh đậm), tam giác nhỏ sát cạnh $TU$ (cam) và tam giác nhỏ sát cạnh $VW$ (xanh lá).

\step{2}{Tam giác $TZW$}
$TQWZ$ là hình vuông $2\times 2$ (gồm $4$ ô đơn vị) và $TW$ là đường chéo của nó.
\begin{recallbox}Đường chéo chia hình vuông (hình chữ nhật) thành hai tam giác có diện tích bằng nhau.\end{recallbox}
Vậy $S_{TZW}=\tfrac12\cdot 4=2$.

\step{3}{Hai tam giác nhỏ}
Xét ô đơn vị có cạnh trên là $TU$. Đường $YU$ đi ``sang phải $1$, lên $1$'' qua các góc ô, nên chứa một
đường chéo của ô này; đường $TW$ đi ``sang phải $1$, xuống $1$'' từ $T$, nên chứa đường chéo còn lại.
Hai đường chéo của một hình vuông chia nó thành $4$ tam giác bằng nhau, và tam giác cam là một trong số đó:
diện tích $\tfrac14\cdot 1=\tfrac14$. Hoàn toàn tương tự, tam giác xanh lá nằm trong ô đơn vị có cạnh phải
là $VW$, được cắt bởi hai đường chéo $XV$ và $TW$: diện tích cũng là $\tfrac14$.

\step{4}{Tính phân số}
Diện tích phần tô màu: $2+\tfrac14+\tfrac14=\tfrac52$. Phân số cần tìm:
\[ \frac{5}{2}:9=\frac{5}{2}\cdot\frac{1}{9}=\frac{5}{18}. \]

\begin{warnbox}Đừng ``ước lượng bằng mắt''! Phần tô màu trông gần bằng $\tfrac13$, nhưng tính chính xác
cho $\tfrac{5}{18}$ (nhỏ hơn $\tfrac{6}{18}=\tfrac13$ một chút).\end{warnbox}
\verify{Đặt $S(0;0)$, cạnh $3$; tính diện tích đa giác tô màu bằng cách cộng các tam giác cho cùng kết quả
$\tfrac52$, và $\tfrac{5}{2}=\tfrac{5}{18}\cdot 9$. Đúng.}
Phương án \textbf{(A)}.
\ans{(A) $\dfrac{5}{18}$}
}

% ---------------------------------------------------------------- 25
\begin{probbox}{25}{\relax}
Người ta xếp các hình chữ nhật kích thước $4\times 2$ theo quy tắc: kể từ hình thứ hai, đỉnh trên -- trái
của mỗi hình được đặt đúng vào trung điểm của cạnh dưới của hình ngay phía trên nó (như hình vẽ). Khi xếp
tất cả $10$ hình chữ nhật theo cách này, chu vi của cả hình là bao nhiêu?
\begin{figC}\begin{scope}[x=0.42cm,y=0.42cm]
\foreach \i in {0,1,2,3}{\draw[Ink,line width=0.9pt] (2*\i,-2*\i) rectangle (2*\i+4,-2*\i+2);}
\foreach \i in {0,1,2}{\draw[Ink] (2*\i+1,-2*\i-0.25)--(2*\i+1,-2*\i+0.25) (2*\i+3,-2*\i-0.25)--(2*\i+3,-2*\i+0.25);}
\node[above] at (2,2){$4$};\node[left] at (0,1){$2$};
\foreach \k in {0,1,2}{\fill[Ink] (9+0.6*\k,-6.8-0.6*\k) circle (2.5pt);}
\end{scope}\end{figC}
\choices{$48$}{$64$}{$90$}{$84$}{$100$}
\end{probbox}
\probsol{
\textbf{Phân tích đề.} Mỗi hình mới lùi sang phải $2$ (nửa chiều dài) và xuống dưới $2$ (đúng bằng chiều
rộng). Cả hình có dạng ``cầu thang''. Ta tính chu vi theo hai cách.

\step{1}{Cách 1: đếm từng loại cạnh}
\emph{Cạnh đứng.} Mỗi hình chữ nhật có cạnh trái và cạnh phải dài $2$, và cả hai đều nằm trên đường viền
(không bị hình khác che). Tổng: $10\times(2+2)=40$.

\emph{Cạnh ngang phía trên.} Hình đầu tiên lộ cả cạnh trên: $4$. Mỗi hình sau, nửa trái cạnh trên áp vào
cạnh dưới của hình trước, chỉ lộ nửa phải: $2$. Tổng: $4+9\times 2=22$.

\emph{Cạnh ngang phía dưới.} Mỗi hình trong $9$ hình đầu chỉ lộ nửa trái cạnh dưới: $2$; hình cuối lộ cả
cạnh dưới: $4$. Tổng: $9\times 2+4=22$.

Chu vi: $40+22+22=84$.

\step{2}{Cách 2: ``kéo thẳng'' cầu thang}
\begin{proofbox}Đi dọc đường viền phía trên -- phải từ góc trên -- trái của hình đầu tiên đến góc dưới --
phải của hình cuối, ta chỉ đi sang phải hoặc đi xuống. Các đoạn ngang cộng lại đúng bằng bề ngang cả hình,
các đoạn đứng cộng lại đúng bằng chiều cao cả hình. Đường viền phía dưới -- trái cũng vậy. Do đó chu vi
cầu thang \emph{bằng chu vi hình chữ nhật nhỏ nhất bao quanh nó}.\end{proofbox}
Bề ngang: $4+9\times 2=22$; chiều cao: $10\times 2=20$. Chu vi: $(22+20)\times 2=84$.

\begin{warnbox}Không lấy $10\times$ chu vi một hình ($10\times 12=120$): những đoạn hai hình áp vào nhau
nằm \emph{bên trong}, không thuộc đường viền.\end{warnbox}
\verify{Thử với $2$ hình: cách 1 cho $8+(4+2)+(2+4)=20$; hình bao có kích thước $6\times 4$, chu vi $20$.
Với $n$ hình, chu vi là $2\big((4+2(n-1))+2n\big)=8n+4$; $n=10$ cho $84$. Khớp.}
Phương án \textbf{(D)}.
\ans{(D) $84$}
}

% ---------------------------------------------------------------- 26
\begin{probbox}{26}{\relax}
Một tứ giác có bốn đỉnh là các điểm $A(7;6)$, $B(-5;1)$, $C(-2;-3)$ và $D(10;2)$ trên mặt phẳng tọa độ.
Diện tích tứ giác (tính theo đơn vị diện tích) là
\choices{$60$}{$63$}{$67$}{$70$}{$72$}
\end{probbox}
\probsol{
\textbf{Phân tích đề.} Tứ giác có các cạnh ``xiên'', không có sẵn công thức. Kĩ thuật hay dùng:
\emph{bao tứ giác bằng một hình chữ nhật có cạnh song song với các trục}, rồi trừ đi các tam giác vuông ở
bốn góc.
\begin{figC}\begin{scope}[x=0.4cm,y=0.4cm]
\fill[tmOrange!15] (-5,-3) rectangle (10,6);
\fill[tmBlue!30] (7,6)--(-5,1)--(-2,-3)--(10,2)--cycle;
\draw[Line] (-5,-3) grid[step=1] (10,6);
\draw[->,Slate] (-6,0)--(11,0) node[right]{$x$};\draw[->,Slate] (0,-4)--(0,7) node[above]{$y$};
\draw[tmOrange!80!black,line width=0.9pt] (-5,-3) rectangle (10,6);
\draw[tmBlue!80!black,line width=1.2pt] (7,6)--(-5,1)--(-2,-3)--(10,2)--cycle;
\node[above] at (7,6){$A$};\node[left] at (-5,1){$B$};\node[below] at (-2,-3){$C$};\node[right] at (10,2){$D$};
\node[font=\footnotesize] at (-2.5,4.6){$30$};\node[font=\footnotesize] at (-4.2,-2){$6$};
\node[font=\footnotesize] at (6.5,-1.6){$30$};\node[font=\footnotesize] at (9.1,4.9){$6$};
\end{scope}\end{figC}

\step{1}{Hình chữ nhật bao quanh}
Hoành độ nhỏ nhất là $-5$, lớn nhất là $10$; tung độ nhỏ nhất là $-3$, lớn nhất là $6$.
\begin{recallbox}Khoảng cách giữa hai điểm trên một trục số bằng số lớn trừ số nhỏ: từ $-5$ đến $10$ là
$10-(-5)=15$.\end{recallbox}
Hình chữ nhật có chiều dài $10-(-5)=15$, chiều rộng $6-(-3)=9$, diện tích $15\times 9=135$.

\step{2}{Bốn tam giác vuông ở góc}
\begin{recallbox}Diện tích tam giác vuông $=\tfrac12\times$ tích hai cạnh góc vuông.\end{recallbox}
\begin{itemize}[leftmargin=1.4em,itemsep=1pt]
\item Góc trên -- trái (giữa $B$ và $A$): hai cạnh $7-(-5)=12$ và $6-1=5$: diện tích $\tfrac{12\cdot5}{2}=30$.
\item Góc dưới -- trái (giữa $B$ và $C$): hai cạnh $-2-(-5)=3$ và $1-(-3)=4$: diện tích $\tfrac{3\cdot4}{2}=6$.
\item Góc dưới -- phải (giữa $C$ và $D$): hai cạnh $10-(-2)=12$ và $2-(-3)=5$: diện tích $30$.
\item Góc trên -- phải (giữa $D$ và $A$): hai cạnh $10-7=3$ và $6-2=4$: diện tích $6$.
\end{itemize}
Tổng: $30+6+30+6=72$.

\step{3}{Diện tích tứ giác}
$135-72=63$.

\begin{warnbox}Với tọa độ âm, rất dễ tính nhầm độ dài, chẳng hạn viết $7-5=2$ thay vì $7-(-5)=12$.
Hãy đếm ô trên hình để kiểm tra.\end{warnbox}
\verify{Thực ra $ABCD$ là hình bình hành ($A\to B$ và $D\to C$ đều là ``sang trái $12$, xuống $5$''). Tam giác
$ABD$ có diện tích bằng nửa tứ giác; bao nó bằng hình chữ nhật $15\times 5$ (từ $y=1$ đến $6$):
$75-30-\tfrac{15\cdot 1}{2}-6=31{,}5$. Nhân đôi: $63$. Khớp.}
Phương án \textbf{(B)}.
\ans{(B) $63$}
}

% ---------------------------------------------------------------- 27
\begin{probbox}{27}{\relax}
Hình dưới đây gồm $8$ hình bình hành nhỏ giống hệt nhau ghép lại. Kể cả $8$ hình bình hành nhỏ đó, trong
hình có tất cả bao nhiêu hình bình hành?
\begin{figC}\begin{scope}[x=0.6cm,y=0.6cm]
\foreach \k in {0,...,4}{\draw[Ink,line width=0.9pt] (2.2*\k,0)--(2.2*\k+2,2);}
\foreach \j in {0,1,2}{\draw[Ink,line width=0.9pt] (\j,\j)--(\j+8.8,\j);}
\end{scope}\end{figC}
\choices{$29$}{$30$}{$26$}{$27$}{$28$}
\end{probbox}
\probsol{
\textbf{Phân tích đề.} Mỗi hình bình hành trong hình có hai cạnh nằm trên hai đường \emph{ngang} và hai cạnh
nằm trên hai đường \emph{xiên} của hình. Ta đếm có hệ thống để không sót, không trùng.

\step{1}{Đếm theo chiều cao và chiều ngang (tính bằng số ô)}
\emph{Hình cao $1$ hàng} (có $2$ hàng để chọn). Trong một hàng có $4$ ô liền nhau:
\begin{itemize}[leftmargin=1.4em,itemsep=1pt]
\item rộng $1$ ô: $4$ hình; rộng $2$ ô: $3$ hình; rộng $3$ ô: $2$ hình; rộng $4$ ô: $1$ hình;
\item mỗi hàng có $4+3+2+1=10$ hình, hai hàng có $20$ hình.
\end{itemize}
\emph{Hình cao $2$ hàng}: tương tự có $4+3+2+1=10$ hình.

Tổng: $20+10=30$.

\step{2}{Cách đếm nhanh: chọn đường}
\begin{proofbox}Một hình bình hành được xác định hoàn toàn khi ta chọn $2$ trong $3$ đường ngang (cạnh trên,
cạnh dưới) và $2$ trong $5$ đường xiên (cạnh trái, cạnh phải). Ngược lại, mỗi cách chọn như vậy cho đúng một
hình bình hành.\end{proofbox}
\begin{recallbox}Số cách chọn $2$ đường từ $k$ đường là $\dfrac{k\times(k-1)}{2}$ (mỗi cặp được đếm hai lần
nếu ta chọn ``đường thứ nhất, đường thứ hai'').\end{recallbox}
Chọn $2$ trong $3$ đường ngang: $\tfrac{3\cdot 2}{2}=3$ cách. Chọn $2$ trong $5$ đường xiên:
$\tfrac{5\cdot4}{2}=10$ cách. Số hình bình hành: $3\times 10=30$.

\begin{warnbox}Các bạn hay quên những hình cao $2$ hàng, hoặc quên chính hình lớn nhất (cả hình).\end{warnbox}
\verify{Hai cách đếm cho cùng kết quả $30$.}
Phương án \textbf{(B)}.
\ans{(B) $30$}
}

% ---------------------------------------------------------------- 28
\begin{probbox}{28}{\relax}
Một hình lăng trụ đứng tam giác có thể tích $120$ cm$^3$. Mặt đáy là tam giác vuông có hai cạnh góc vuông
dài $3$ cm và $4$ cm (như hình vẽ). Chiều cao của hình lăng trụ là bao nhiêu xăng-ti-mét?
\begin{figC}\begin{scope}[x=0.75cm,y=0.75cm]
\coordinate (A) at (0,2.8);\coordinate (B) at (1.9,3.8);\coordinate (C) at (4.3,2.2);
\coordinate (A1) at (0,0);\coordinate (B1) at (1.9,1.0);\coordinate (C1) at (4.3,-0.6);
\fill[tmBlue!15] (A)--(B)--(C)--cycle;
\draw[Ink,line width=1pt] (A)--(B)--(C)--cycle (A)--(A1)--(C1)--(C);
\draw[Ink,dashed,line width=0.8pt] (B)--(B1) (A1)--(B1)--(C1);
\draw[Ink] ($(B)!0.3cm!(A)$)--($(B)!0.3cm!(A)+(B)!0.3cm!(C)-(B)$)--($(B)!0.3cm!(C)$);
\draw[Ink] ($(B1)!0.3cm!(A1)$)--($(B1)!0.3cm!(A1)+(B1)!0.3cm!(C1)-(B1)$)--($(B1)!0.3cm!(C1)$);
\node[above left] at ($(A)!0.5!(B)$){$3$ cm};\node[above right] at ($(B)!0.5!(C)$){$4$ cm};
\end{scope}\end{figC}
\choices{$12$}{$20$}{$10$}{$16$}{$8$}
\end{probbox}
\probsol{
\textbf{Phân tích đề.} Biết thể tích và biết đáy (tam giác vuông). Cần tìm chiều cao lăng trụ, tức là độ dài
cạnh bên (nối đáy trên với đáy dưới).

\begin{recallbox}Thể tích hình lăng trụ đứng $=$ diện tích đáy $\times$ chiều cao:
$V=S_{\text{đáy}}\cdot h$. Do đó $h=V:S_{\text{đáy}}$.\end{recallbox}

\step{1}{Diện tích đáy}
Đáy là tam giác vuông, lấy một cạnh góc vuông làm đáy thì cạnh góc vuông kia là chiều cao:
\[ S_{\text{đáy}}=\frac{3\times 4}{2}=6 \text{ (cm}^2). \]

\step{2}{Chiều cao lăng trụ}
$h=120:6=20$ (cm).

\begin{warnbox}Đừng nhầm ``chiều cao của lăng trụ'' (khoảng cách giữa hai đáy) với chiều cao của tam giác đáy.
Cũng đừng quên chia $2$ khi tính diện tích tam giác: nếu lấy $3\times 4=12$ thì sẽ ra $10$ -- phương án (C)
``bẫy''.\end{warnbox}
\verify{$6\times 20=120$ cm$^3$. Đúng.}
Phương án \textbf{(B)}.
\ans{(B) $20$ cm}
}

% ---------------------------------------------------------------- 29
\begin{probbox}{29}{\stars{1}}
Người ta dùng các khối lập phương nhỏ cạnh $1$ để xếp thành một hình hộp chữ nhật kích thước
$8\times 8\times n$ ($n$ là số nguyên dương). Gọi $A$ là diện tích toàn phần của hình hộp, và $B$ là tổng diện
tích toàn phần của tất cả các khối lập phương nhỏ (tính riêng từng khối rồi cộng lại). Tính tổng tất cả các giá
trị của $n$ sao cho $\dfrac{B}{A}$ là một số nguyên.
\choices{$86$}{$90$}{$70$}{$78$}{$96$}
\end{probbox}
\probsol{
\textbf{Phân tích đề.} Ta viết $A$ và $B$ theo $n$, lập tỉ số $\dfrac BA$, rồi tìm những $n$ làm tỉ số
đó nguyên.

\step{1}{Tính $A$}
\begin{recallbox}Hình hộp chữ nhật kích thước $a\times b\times c$ có $6$ mặt, gồm ba cặp mặt bằng nhau, nên
diện tích toàn phần là $2(ab+bc+ca)$.\end{recallbox}
$A=2(8\cdot 8+8\cdot n+8\cdot n)=2(64+16n)=128+32n$.

\step{2}{Tính $B$}
Số khối lập phương nhỏ: $8\cdot 8\cdot n=64n$. Mỗi khối có $6$ mặt, mỗi mặt diện tích $1$, nên
$B=64n\cdot 6=384n$.

\step{3}{Rút gọn tỉ số}
\[ \frac{B}{A}=\frac{384n}{128+32n}=\frac{32\cdot 12n}{32\cdot(4+n)}=\frac{12n}{n+4}. \]
Ta ``tách phần nguyên'': vì $12n=12(n+4)-48$ nên
\[ \frac{12n}{n+4}=\frac{12(n+4)-48}{n+4}=12-\frac{48}{n+4}. \]

\step{4}{Điều kiện để tỉ số là số nguyên}
$12-\dfrac{48}{n+4}$ là số nguyên khi và chỉ khi $\dfrac{48}{n+4}$ là số nguyên, tức là $n+4$ là
\emph{ước} của $48$.
\begin{recallbox}Các ước dương của $48$: $1,2,3,4,6,8,12,16,24,48$.\end{recallbox}
Vì $n\ge 1$ nên $n+4\ge 5$, ta chỉ lấy các ước $6,8,12,16,24,48$, ứng với
\[ n=2,\ 4,\ 8,\ 12,\ 20,\ 44. \]

\step{5}{Tính tổng}
$2+4+8+12+20+44=90$.

\begin{center}\begin{tabular}{|c|c|c|c|c|c|c|}\hline
$n$ & $2$ & $4$ & $8$ & $12$ & $20$ & $44$\\\hline
$B/A=\dfrac{12n}{n+4}$ & $\tfrac{24}{6}=4$ & $\tfrac{48}{8}=6$ & $\tfrac{96}{12}=8$ & $\tfrac{144}{16}=9$ & $\tfrac{240}{24}=10$ & $\tfrac{528}{48}=11$\\\hline
\end{tabular}\end{center}
\begin{warnbox}Đừng quên $n+4$ có thể bằng $48$ (tức $n=44$), một giá trị khá lớn. Và nhớ loại các ước
$1,2,3,4$ vì khi đó $n\le 0$.\end{warnbox}
\verify{Bảng trên cho thấy cả $6$ giá trị đều cho tỉ số nguyên. Thử một giá trị khác, $n=1$:
$\tfrac{12}{5}$ không nguyên. Đúng.}
Phương án \textbf{(B)}.
\ans{(B) $90$}
}

% ---------------------------------------------------------------- 30
\begin{probbox}{30}{\relax}
Hình chữ nhật $PQRS$ được ghép từ sáu hình vuông bằng nhau. Các điểm $U,V,W,X,Y,Z$ là trung điểm của
cạnh các hình vuông như hình vẽ. Tam giác nào dưới đây có diện tích lớn nhất?
\begin{figC}\begin{scope}[x=0.6cm,y=0.6cm]
\draw[Ink,line width=1pt] (0,0) rectangle (6,4);
\draw[Ink,line width=0.9pt] (2,0)--(2,4) (4,0)--(4,4) (0,2)--(6,2);
\foreach \p in {(1,4),(3,4),(4,3),(5,2),(4,1),(6,1)}{\fill[Ink] \p circle (2.2pt);}
\node[above] at (1,4){$V$};\node[above] at (3,4){$W$};\node[right] at (4,3){$X$};
\node[below] at (5,2){$Y$};\node[left] at (4,1){$U$};\node[right] at (6,1){$Z$};
\node[above left] at (0,4){$Q$};\node[above right] at (6,4){$R$};
\node[below left] at (0,0){$P$};\node[below right] at (6,0){$S$};
\end{scope}\end{figC}
\choices{$PVU$}{$PXZ$}{$PVX$}{$PYS$}{$PQW$}
\end{probbox}
\probsol{
\textbf{Phân tích đề.} Ta cần tính (hoặc so sánh) diện tích của $5$ tam giác. Để tránh phân số, chọn cạnh
mỗi hình vuông nhỏ là $2$; khi đó các trung điểm có tọa độ nguyên. Đặt $P(0;0)$, $S(6;0)$, $R(6;4)$, $Q(0;4)$;
suy ra $V(1;4)$, $W(3;4)$, $X(4;3)$, $Y(5;2)$, $U(4;1)$, $Z(6;1)$.

\begin{recallbox}Diện tích tam giác $=\tfrac12\times$ đáy $\times$ chiều cao. Với tam giác ``xiên'', bao nó
bằng hình chữ nhật có cạnh theo lưới rồi trừ các tam giác vuông ở góc (như Bài 26).\end{recallbox}

\step{1}{Hai tam giác có cạnh nằm trên cạnh hình chữ nhật}
\begin{itemize}[leftmargin=1.4em,itemsep=1pt]
\item $PYS$: đáy $PS=6$, chiều cao là khoảng cách từ $Y$ đến $PS$, bằng $2$: $S=\tfrac{6\cdot 2}{2}=6$.
\item $PQW$: đáy $PQ=4$, chiều cao là khoảng cách từ $W$ đến $PQ$, bằng $3$: $S=\tfrac{4\cdot 3}{2}=6$.
\end{itemize}

\step{2}{Ba tam giác ``xiên''}
\begin{itemize}[leftmargin=1.4em,itemsep=1pt]
\item $PVU$: hình chữ nhật bao $4\times 4=16$; trừ $\tfrac{1\cdot 4}{2}=2$, $\tfrac{3\cdot 3}{2}=4{,}5$,
$\tfrac{4\cdot 1}{2}=2$. Còn $16-8{,}5=7{,}5$.
\item $PXZ$: hình chữ nhật bao $6\times 3=18$; trừ $\tfrac{4\cdot3}{2}=6$, $\tfrac{2\cdot2}{2}=2$,
$\tfrac{6\cdot1}{2}=3$. Còn $18-11=7$.
\item $PVX$: hình chữ nhật bao $4\times 4=16$; trừ $\tfrac{1\cdot4}{2}=2$, $\tfrac{3\cdot1}{2}=1{,}5$,
$\tfrac{4\cdot3}{2}=6$. Còn $16-9{,}5=6{,}5$.
\end{itemize}
\begin{figC}\begin{scope}[x=0.6cm,y=0.6cm]
\fill[tmOrange!15] (0,0) rectangle (4,4);
\fill[tmBlue!35] (0,0)--(1,4)--(4,1)--cycle;
\draw[Line] (0,0) grid[step=1] (6,4);
\draw[Ink,line width=1pt] (0,0) rectangle (6,4);
\draw[tmBlue!80!black,line width=1.1pt] (0,0)--(1,4)--(4,1)--cycle;
\node[font=\footnotesize] at (0.35,3.4){$2$};\node[font=\footnotesize] at (3,3){$4{,}5$};\node[font=\footnotesize] at (3,0.35){$2$};
\node[below left] at (0,0){$P$};\node[above] at (1,4){$V$};\node[right] at (4,1){$U$};
\end{scope}\end{figC}

\step{3}{So sánh}
$7{,}5>7>6{,}5>6=6$. Tam giác lớn nhất là $PVU$.

\begin{warnbox}Đừng đoán bằng mắt: tam giác $PXZ$ trông ``to'' vì kéo dài sang tận cạnh $RS$, nhưng diện tích
chỉ là $7$, nhỏ hơn $PVU$.\end{warnbox}
\verify{Nếu lấy cạnh hình vuông là $1$ thì mọi diện tích đều chia cho $4$ ($1{,}875$; $1{,}75$; $1{,}625$;
$1{,}5$; $1{,}5$), thứ tự không đổi.}
Phương án \textbf{(A)}.
\ans{(A) tam giác $PVU$}
}

% ---------------- part c ----------------
% ======================================================================
% Chủ đề 2 -- Hình học và Đo lường -- phần c (bài 31--40)
% ======================================================================
\clearpage
\dividerbar{Mở rộng -- Nâng cao}

\begin{probbox}{31}{\relax}
Trong mặt phẳng toạ độ $Oxy$, cho điểm $A(-3;\,2)$.
\begin{enumerate}[label=\alph*),leftmargin=1.6em,itemsep=0pt,topsep=2pt]
  \item Gọi $B$ là điểm đối xứng với $A$ qua trục $Ox$, $C$ là điểm đối xứng với $A$ qua trục $Oy$, $D$ là điểm đối xứng với $A$ qua gốc $O$. Tìm toạ độ của $B$, $C$, $D$ và cho biết mỗi điểm nằm ở góc phần tư nào.
  \item Bốn điểm $A$, $C$, $D$, $B$ là bốn đỉnh của một hình chữ nhật. Tính chu vi và diện tích hình chữ nhật đó.
\end{enumerate}
\end{probbox}
\probsol{
\textbf{Phân tích đề.} ``Đối xứng qua một trục'' giống như soi gương qua trục đó: điểm ảnh nằm ở phía bên kia trục, cách trục một khoảng bằng điểm ban đầu. Ta xem toạ độ nào đổi dấu, toạ độ nào giữ nguyên.
\begin{recallbox}Đối xứng qua trục $Ox$: $(x;\,y)\to(x;\,-y)$. Đối xứng qua trục $Oy$: $(x;\,y)\to(-x;\,y)$. Đối xứng qua gốc $O$ (cũng là quay $180^\circ$ quanh $O$): $(x;\,y)\to(-x;\,-y)$.\end{recallbox}
\step{1}{Điểm $B$ -- đối xứng qua $Ox$}
Soi gương qua trục hoành: điểm giữ nguyên vị trí trái -- phải (hoành độ không đổi), chỉ đổi từ phía trên xuống phía dưới (tung độ đổi dấu). $B(-3;\,-2)$, dạng $(-;-)$: góc phần tư III.

\step{2}{Điểm $C$ -- đối xứng qua $Oy$}
Soi gương qua trục tung: tung độ giữ nguyên, hoành độ đổi dấu. $C(3;\,2)$, dạng $(+;+)$: góc phần tư I.

\step{3}{Điểm $D$ -- đối xứng qua gốc $O$}
Cả hai toạ độ đổi dấu: $D(3;\,-2)$, dạng $(+;-)$: góc phần tư IV.
\begin{gmFig}[scale=0.55]
  \draw[Line,line width=0.3pt] (-4.5,-3.5) grid (4.5,3.5);
  \draw[->,Navy,thick] (-4.8,0)--(4.9,0) node[right,text=Ink]{$x$};
  \draw[->,Navy,thick] (0,-3.8)--(0,3.9) node[above,text=Ink]{$y$};
  \foreach \t in {-4,...,-1,1,2,3,4}{\node[below,font=\scriptsize] at (\t,0){$\t$};}
  \foreach \t in {-3,...,-1,1,2,3}{\node[left,font=\scriptsize] at (0,\t){$\t$};}
  \fill[tmBlue!15] (-3,2) rectangle (3,-2);
  \draw[tmBlue,thick] (-3,2) rectangle (3,-2);
  \fill[AnsRed] (-3,2) circle (3pt) node[above left]{$A$};
  \fill[tmPurple] (-3,-2) circle (3pt) node[below left]{$B$};
  \fill[tmGreen!70!black] (3,2) circle (3pt) node[above right]{$C$};
  \fill[tmOrange] (3,-2) circle (3pt) node[below right]{$D$};
  \node[above,font=\footnotesize,text=tmBlue] at (1.5,2){$6$};
  \node[right,font=\footnotesize,text=tmBlue] at (3,1){$4$};
\end{gmFig}

\step{4}{Tính kích thước hình chữ nhật}
$A(-3;\,2)$ và $C(3;\,2)$ cùng tung độ: $AC=3-(-3)=6$.
$A(-3;\,2)$ và $B(-3;\,-2)$ cùng hoành độ: $AB=2-(-2)=4$.

\step{5}{Tính chu vi và diện tích}
\begin{recallbox}Hình chữ nhật dài $a$, rộng $b$: $P=2\cdot(a+b)$, $S=a\cdot b$.\end{recallbox}
$P=2\cdot(6+4)=20$ (đơn vị độ dài); $\;S=6\cdot4=24$ (đơn vị diện tích).
\begin{tipbox}Hình chữ nhật này nhận cả hai trục toạ độ làm trục đối xứng. Mỗi góc phần tư chứa đúng một phần tư hình, là hình chữ nhật $3\times2$ có diện tích $6$; và $4\cdot6=24$.\end{tipbox}
\verify{Mỗi góc phần tư chứa một điểm: $A$ (II), $C$ (I), $B$ (III), $D$ (IV), đúng như hình.}
\ans{a) $B(-3;\,-2)$: III; $C(3;\,2)$: I; $D(3;\,-2)$: IV. \quad b) Chu vi $20$, diện tích $24$.}
}

\begin{probbox}{32}{\relax}
Trong hình vẽ, tam giác $ABC$ có $\widehat{BAC}=40^\circ$; cạnh $BC$ được kéo dài về cả hai phía, cạnh $AB$ được kéo dài qua $B$. Góc ngoài tại đỉnh $C$ bằng $115^\circ$. Giá trị của $x$ là
\begin{gmFig}[scale=0.55]
  \coordinate (B) at (0,0); \coordinate (C) at (4,0); \coordinate (A) at (1.460,5.448);
  \coordinate (L) at (-1.6,0); \coordinate (D) at (-0.42,-1.567); \coordinate (E) at (5.8,0);
  \draw[Ink,thick] (L)--(E); \draw[Ink,thick] (D)--(A)--(C);
  \pic[draw,thick,angle radius=0.7cm,"$40^\circ$",angle eccentricity=1.45]{angle=B--A--C};
  \pic[draw,thick,angle radius=0.45cm,"$115^\circ$",angle eccentricity=2.1]{angle=E--C--A};
  \pic[draw,AnsRed,thick,angle radius=0.45cm,"$x^\circ$",angle eccentricity=1.8]{angle=L--B--D};
  \node[above] at (A){$A$}; \node[above left] at (B){$B$}; \node[below] at (C){$C$};
\end{gmFig}
\choices{$65$}{$75$}{$105$}{$115$}{$140$}
\end{probbox}
\probsol{
\textbf{Phân tích đề.} Góc $x^\circ$ ở đỉnh $B$ được tạo bởi hai tia kéo dài (tia đối của $BC$ và tia đối của $BA$), nên nó \emph{đối đỉnh} với góc trong $\widehat{ABC}$. Vậy chỉ cần tìm $\widehat{ABC}$. Ta giải bằng hai cách.

\step{1}{Cách 1 -- qua góc trong tại $C$}
\begin{recallbox}Hai góc kề bù có tổng $180^\circ$; tổng ba góc trong tam giác bằng $180^\circ$.\end{recallbox}
Góc trong tại $C$ kề bù với góc ngoài $115^\circ$: $\widehat{ACB}=180^\circ-115^\circ=65^\circ$.
Góc trong tại $B$: $\widehat{ABC}=180^\circ-40^\circ-65^\circ=75^\circ$.

\step{2}{Cách 2 -- dùng tính chất góc ngoài}
\begin{recallbox}Góc ngoài của tam giác bằng tổng hai góc trong không kề với nó.\end{recallbox}
Góc ngoài tại $C$ bằng $\widehat{A}+\widehat{B}$: $115^\circ=40^\circ+\widehat{ABC}$, suy ra $\widehat{ABC}=115^\circ-40^\circ=75^\circ$.

\step{3}{Suy ra $x$}
\begin{recallbox}Hai góc đối đỉnh thì bằng nhau.\end{recallbox}
$x^\circ=\widehat{ABC}=75^\circ$, nên $x=75$.
\begin{warnbox}
\begin{itemize}[leftmargin=1.3em,itemsep=0pt,topsep=0pt]
  \item (A) $65$ là góc trong tại $C$, không phải tại $B$.
  \item (C) $105$ là góc kề bù với $\widehat{ABC}$ -- nhầm ``kề bù'' với ``đối đỉnh''.
\end{itemize}
\end{warnbox}
\verify{$40^\circ+75^\circ+65^\circ=180^\circ$ và $40^\circ+75^\circ=115^\circ$ (góc ngoài tại $C$). Cả hai điều kiện đều đúng.}
Vậy chọn \textbf{(B)}.
\ans{(B) $75$}
}

\begin{probbox}{33}{\relax}
\begin{enumerate}[label=\alph*),leftmargin=1.6em,itemsep=0pt,topsep=0pt]
  \item Một tam giác cân có góc ở đỉnh gấp $4$ lần mỗi góc ở đáy. Tính số đo các góc của tam giác.
  \item Một tam giác cân có một góc bằng $36^\circ$. Tìm tất cả các khả năng cho số đo hai góc còn lại. Với mỗi khả năng, tính hiệu giữa góc lớn nhất và góc nhỏ nhất của tam giác.
  \item Có tam giác cân nào có một góc ở đáy bằng $92^\circ$ không? Vì sao?
\end{enumerate}
\end{probbox}
\probsol{
\textbf{Phân tích đề.} Cả ba câu đều dựa trên hai sự thật: tam giác cân có hai góc ở đáy bằng nhau, và tổng ba góc bằng $180^\circ$.
\begin{recallbox}Tam giác cân: góc ở đỉnh $d$, mỗi góc ở đáy $g$, thì $d+2g=180^\circ$.\end{recallbox}
\step{1}{Câu a -- dùng ``phần bằng nhau''}
Coi mỗi góc đáy là $1$ phần. Góc ở đỉnh gấp $4$ lần nên là $4$ phần. Tổng ba góc gồm $1+1+4=6$ phần, ứng với $180^\circ$.
\begin{barfigC}
  \blk{0}{1}{0}{0.6}{tmBlue}{$g$}
  \blk{1}{2}{0}{0.6}{tmBlue}{$g$}
  \blk{2}{3}{0}{0.6}{tmOrange}{}
  \blk{3}{4}{0}{0.6}{tmOrange}{}
  \blk{4}{5}{0}{0.6}{tmOrange}{}
  \blk{5}{6}{0}{0.6}{tmOrange}{}
  \brcup{2}{6}{0.5}{góc ở đỉnh $=4g$}
  \brcdn{0}{6}{-0.5}{$180^\circ$}
\end{barfigC}
Một phần: $180^\circ:6=30^\circ$. Vậy mỗi góc đáy là $30^\circ$, góc ở đỉnh là $4\cdot30^\circ=120^\circ$.

\step{2}{Câu b -- xét hai trường hợp}
\emph{Trường hợp 1: $36^\circ$ là góc ở đỉnh.} Mỗi góc đáy: $(180^\circ-36^\circ):2=144^\circ:2=72^\circ$. Tam giác $36^\circ,\,72^\circ,\,72^\circ$. Hiệu góc lớn nhất và nhỏ nhất: $72^\circ-36^\circ=36^\circ$.

\emph{Trường hợp 2: $36^\circ$ là góc ở đáy.} Góc đáy kia $36^\circ$, góc ở đỉnh $180^\circ-72^\circ=108^\circ$. Tam giác $36^\circ,\,36^\circ,\,108^\circ$. Hiệu: $108^\circ-36^\circ=72^\circ$.

\step{3}{Câu c -- lập luận}
\begin{proofbox}Nếu một góc ở đáy bằng $92^\circ$ thì góc ở đáy còn lại cũng bằng $92^\circ$. Khi đó riêng hai góc đáy đã có tổng $184^\circ$, lớn hơn $180^\circ$ -- vô lí, vì tổng \emph{ba} góc chỉ bằng $180^\circ$. Vậy không có tam giác cân nào như thế. (Tổng quát: góc ở đáy của tam giác cân luôn nhỏ hơn $90^\circ$.)\end{proofbox}
\verify{$30+30+120=180$; $\;36+72+72=180$; $\;36+36+108=180$.}
\ans{a) $30^\circ$, $30^\circ$, $120^\circ$. \ b) $72^\circ$ và $72^\circ$ (hiệu $36^\circ$), hoặc $36^\circ$ và $108^\circ$ (hiệu $72^\circ$). \ c) Không có.}
}

\begin{probbox}{34}{\stars{1}}
Trong mặt phẳng toạ độ, cho điểm $A(4;\,1)$.
\begin{enumerate}[label=\alph*),leftmargin=1.6em,itemsep=0pt,topsep=2pt]
  \item Quay $A$ quanh gốc $O$ một góc $90^\circ$ ngược chiều kim đồng hồ được điểm $B$; quay $B$ thêm $90^\circ$ ngược chiều kim đồng hồ được $C$; quay $C$ thêm $90^\circ$ ngược chiều kim đồng hồ được $E$. Tìm toạ độ $B$, $C$, $E$.
  \item Quay $A$ quanh $O$ một góc $270^\circ$ \emph{theo chiều kim đồng hồ} thì được điểm nào trong ba điểm trên?
  \item Tứ giác $ABCE$ là một hình vuông. Tính diện tích của nó.
\end{enumerate}
\end{probbox}
\probsol{
\textbf{Phân tích đề.} Ta dùng quy tắc quay $90^\circ$ trong mặt phẳng toạ độ, sau đó tính diện tích hình vuông ``nghiêng'' bằng cách đặt nó vào trong một hình vuông lớn có cạnh nằm theo các đường kẻ ô.
\begin{recallbox}Quay $90^\circ$ ngược chiều kim đồng hồ quanh $O$: $(x;\,y)\to(-y;\,x)$. Quay $\alpha$ theo chiều kim đồng hồ giống quay $360^\circ-\alpha$ ngược chiều.\end{recallbox}
\begin{proofbox}Vì sao $(x;\,y)\to(-y;\,x)$? Điểm $A(4;\,1)$ đạt được bằng cách đi ``sang phải $4$, lên $1$''. Quay cả đường đi này $90^\circ$ ngược chiều kim đồng hồ: ``sang phải'' biến thành ``lên trên'', ``lên trên'' biến thành ``sang trái''. Vậy ảnh đạt được bằng ``lên $4$, sang trái $1$'', tức điểm $(-1;\,4)$.\end{proofbox}
\step{1}{Câu a -- áp dụng quy tắc ba lần}
$A(4;\,1)\to B(-1;\,4)\to C(-4;\,-1)\to E(1;\,-4)$.

(Kiểm tra: $C$ là ảnh của $A$ qua phép quay $180^\circ$, đúng là $(-4;\,-1)$, cả hai toạ độ đổi dấu.)

\step{2}{Câu b}
Quay $270^\circ$ theo chiều kim đồng hồ giống quay $360^\circ-270^\circ=90^\circ$ ngược chiều kim đồng hồ, nên được đúng điểm $B(-1;\,4)$.

\step{3}{Câu c -- đặt hình vuông vào khung}
\begin{gmFig}[scale=0.42]
  \draw[Line,line width=0.3pt] (-4,-4) grid (4,4);
  \draw[->,Navy] (-4.6,0)--(4.7,0) node[right,text=Ink]{$x$};
  \draw[->,Navy] (0,-4.6)--(0,4.7) node[above,text=Ink]{$y$};
  \draw[tmOrange!90!black,thick,dashed] (-4,-4) rectangle (4,4);
  \fill[tmOrange!25] (4,1)--(4,4)--(-1,4)--cycle;
  \fill[tmOrange!25] (-1,4)--(-4,4)--(-4,-1)--cycle;
  \fill[tmOrange!25] (-4,-1)--(-4,-4)--(1,-4)--cycle;
  \fill[tmOrange!25] (1,-4)--(4,-4)--(4,1)--cycle;
  \fill[tmBlue!25,draw=tmBlue,thick] (4,1)--(-1,4)--(-4,-1)--(1,-4)--cycle;
  \node[right] at (4,1){$A$}; \node[above] at (-1,4){$B$}; \node[left] at (-4,-1){$C$}; \node[below] at (1,-4){$E$};
  \node[font=\scriptsize] at (2.9,3.2){$5\times3$};
\end{gmFig}
Hình vuông $ABCE$ nằm gọn trong hình vuông lớn có các đỉnh $(\pm4;\,\pm4)$, cạnh $4-(-4)=8$, diện tích $8\cdot8=64$.

Phần thừa gồm $4$ tam giác vuông bằng nhau ở bốn góc. Ví dụ, tam giác ở góc trên bên phải có đỉnh $A(4;\,1)$, $(4;\,4)$, $B(-1;\,4)$: cạnh nằm ngang dài $4-(-1)=5$, cạnh thẳng đứng dài $4-1=3$, diện tích $\dfrac{5\cdot3}{2}=7{,}5$.

Diện tích hình vuông $ABCE$: $64-4\cdot7{,}5=64-30=34$ (đơn vị diện tích).
\begin{tipbox}``Đóng khung rồi trừ bớt'' là cách rất mạnh để tính diện tích một đa giác có các đỉnh nằm ở nút lưới: vẽ hình chữ nhật nhỏ nhất bao quanh, rồi trừ các tam giác vuông (hình thang, hình chữ nhật) ở phía ngoài.\end{tipbox}
\verify{Bốn tam giác vuông bằng nhau (mỗi tam giác $5\times3$) cùng với hình vuông $ABCE$ lấp đầy khung $8\times8$: $4\cdot7{,}5+34=64$. Ngoài ra, cạnh khung $8=5+3$ đúng bằng tổng hai cạnh góc vuông.}
\ans{a) $B(-1;\,4)$, $C(-4;\,-1)$, $E(1;\,-4)$. \ b) Điểm $B$. \ c) $34$ đơn vị diện tích.}
}

\begin{probbox}{35}{\relax}
Điền số thích hợp vào chỗ trống:
\par\smallskip
a) $2{,}4\text{ km}=\ldots\text{ m}$; \qquad b) $75\text{ mm}=\ldots\text{ cm}$; \qquad c) $3\text{ m}^2=\ldots\text{ cm}^2$; \qquad d) $2500\text{ cm}^3=\ldots$ lít.
\par\smallskip
e) Một tấm bảng hình chữ nhật dài $1{,}2$ m, rộng $50$ cm. Tính diện tích tấm bảng theo $\text{cm}^2$ và theo $\text{m}^2$.
\end{probbox}
\probsol{
\textbf{Phân tích đề.} Mỗi câu là một phép đổi đơn vị. Điều cần nhớ: độ dài gấp $10$ lần qua mỗi bậc, diện tích gấp $100$ lần, thể tích gấp $1000$ lần.
\begin{recallbox}$1\text{ km}=1000\text{ m}$; $\;1\text{ cm}=10\text{ mm}$; $\;1\text{ m}^2=10\,000\text{ cm}^2$; $\;1\text{ lít}=1\text{ dm}^3=1000\text{ cm}^3$. Lớn sang nhỏ: nhân; nhỏ sang lớn: chia.\end{recallbox}
\step{1}{Câu a}
km lớn hơn m, nên nhân: $2{,}4\cdot1000=2400$. Vậy $2{,}4\text{ km}=2400\text{ m}$.

\step{2}{Câu b}
mm nhỏ hơn cm, nên chia: $75:10=7{,}5$. Vậy $75\text{ mm}=7{,}5\text{ cm}$.

\step{3}{Câu c}
$3\text{ m}^2=3\cdot10\,000\text{ cm}^2=30\,000\text{ cm}^2$.
\begin{warnbox}Đổi diện tích mà chỉ nhân $100$ (như đổi độ dài m sang cm) sẽ được $300$ -- sai. Hình vuông $1\text{ m}\times1\text{ m}$ là $100\text{ cm}\times100\text{ cm}$, chứa $10\,000$ ô vuông $1\text{ cm}^2$.\end{warnbox}

\step{4}{Câu d}
$2500\text{ cm}^3=2500:1000\text{ dm}^3=2{,}5\text{ dm}^3=2{,}5$ lít.

\step{5}{Câu e -- đưa về cùng đơn vị trước}
Hai kích thước khác đơn vị, nên phải đổi trước khi nhân: $1{,}2\text{ m}=120\text{ cm}$.
Diện tích: $120\cdot50=6000\ (\text{cm}^2)$.
Đổi ra $\text{m}^2$: $6000:10\,000=0{,}6\ (\text{m}^2)$.
\verify{Tính trực tiếp bằng mét: $50\text{ cm}=0{,}5\text{ m}$, nên $1{,}2\cdot0{,}5=0{,}6\ (\text{m}^2)$. Khớp.}
\ans{a) $2400$; \ b) $7{,}5$; \ c) $30\,000$; \ d) $2{,}5$; \ e) $6000\text{ cm}^2=0{,}6\text{ m}^2$.}
}

\begin{probbox}{36}{\stars{1}}
Bạn An xếp được khối hộp chữ nhật $4\text{ cm}\times3\text{ cm}\times3\text{ cm}$ từ $36$ khối lập phương nhỏ cạnh $1$ cm (như Bài 10). An sơn tất cả các mặt bên ngoài của khối hộp, rồi tách rời các khối nhỏ ra. Có bao nhiêu khối nhỏ được sơn đúng $3$ mặt? Đúng $2$ mặt? Đúng $1$ mặt? Không được sơn mặt nào?
\begin{gmFig}[scale=0.55]
  \fill[tmBlue!20] (0,0) rectangle (4,3);
  \fill[tmBlue!10] (0,3)--(1.5,4)--(5.5,4)--(4,3)--cycle;
  \fill[tmBlue!30] (4,0)--(5.5,1)--(5.5,4)--(4,3)--cycle;
  \draw[Ink] (0,0) grid (4,3);
  \foreach \t in {1,2}{\draw[Ink] (4+0.5*\t,0.333*\t)--(4+0.5*\t,3+0.333*\t); \draw[Ink] (0.5*\t,3+0.333*\t)--(4+0.5*\t,3+0.333*\t);}
  \foreach \t in {1,2}{\draw[Ink] (4,\t)--(5.5,\t+1);}
  \foreach \t in {1,2,3}{\draw[Ink] (\t,3)--(\t+1.5,4);}
  \draw[Ink,thick] (0,0) rectangle (4,3); \draw[Ink,thick] (0,3)--(1.5,4)--(5.5,4)--(5.5,1)--(4,0); \draw[Ink,thick] (4,3)--(5.5,4);
\end{gmFig}
\end{probbox}
\probsol{
\textbf{Phân tích đề.} Số mặt được sơn của một khối nhỏ phụ thuộc vào \emph{vị trí} của nó: ở đỉnh, nằm trên cạnh, nằm giữa một mặt, hay nằm hẳn bên trong. Ta đếm từng loại.
\step{1}{Khối ở các đỉnh -- sơn $3$ mặt}
Mỗi khối ở đỉnh của hình hộp lộ ra $3$ mặt (một mặt trước/sau, một mặt trên/dưới, một mặt trái/phải). Hình hộp có $8$ đỉnh, nên có $\mathbf{8}$ khối sơn $3$ mặt.

\step{2}{Khối nằm trên cạnh (trừ hai đầu) -- sơn $2$ mặt}
Hình hộp có $12$ cạnh: $4$ cạnh dài $4$ cm, $4$ cạnh dài $3$ cm (chiều cao) và $4$ cạnh dài $3$ cm (chiều sâu). Trên một cạnh dài $n$ khối, bỏ hai khối ở hai đầu (đã đếm ở bước 1), còn $n-2$ khối.
\begin{itemize}[leftmargin=1.3em,itemsep=0pt,topsep=0pt]
  \item $4$ cạnh dài $4$: mỗi cạnh $4-2=2$ khối, tổng $4\cdot2=8$;
  \item $8$ cạnh dài $3$: mỗi cạnh $3-2=1$ khối, tổng $8\cdot1=8$.
\end{itemize}
Có $8+8=\mathbf{16}$ khối sơn $2$ mặt.

\step{3}{Khối nằm giữa một mặt -- sơn $1$ mặt}
Trên một mặt $m\times n$, bỏ đi viền ngoài, còn $(m-2)\cdot(n-2)$ khối.
\begin{itemize}[leftmargin=1.3em,itemsep=0pt,topsep=0pt]
  \item $2$ mặt $4\times3$ (trước, sau): mỗi mặt $2\cdot1=2$ khối, tổng $4$;
  \item $2$ mặt $4\times3$ (trên, dưới): mỗi mặt $2\cdot1=2$ khối, tổng $4$;
  \item $2$ mặt $3\times3$ (trái, phải): mỗi mặt $1\cdot1=1$ khối, tổng $2$.
\end{itemize}
Có $4+4+2=\mathbf{10}$ khối sơn $1$ mặt.

\step{4}{Khối nằm hẳn bên trong -- không sơn}
Bỏ đi một lớp vỏ ở mỗi phía, phần lõi bên trong có kích thước $(4-2)\times(3-2)\times(3-2)=2\times1\times1$, gồm $\mathbf{2}$ khối.
\begin{gmFig}[scale=0.8]
  \foreach \x/\y/\c/\l in {0/0/AnsRed/3,1.4/0/tmOrange/2,2.8/0/tmGreen/1,4.2/0/Slate/0}{
    \fill[\c!30,draw=\c!80!black] (\x,\y) rectangle ++(0.8,0.8);
    \node[font=\footnotesize,below] at (\x+0.4,\y){$\l$ mặt};}
  \node[font=\footnotesize] at (0.4,1.1){$8$}; \node[font=\footnotesize] at (1.8,1.1){$16$};
  \node[font=\footnotesize] at (3.2,1.1){$10$}; \node[font=\footnotesize] at (4.6,1.1){$2$};
\end{gmFig}
\begin{parentbox}Nếu có bộ khối lập phương đồ chơi (hoặc khối rubik), hãy cùng con xếp hình hộp $4\times3\times3$, dán giấy màu lên các mặt ngoài rồi tách ra để đếm. Con sẽ tự thấy vì sao khối ở đỉnh có $3$ mặt màu.\end{parentbox}
\verify{Tổng bốn loại phải bằng số khối: $8+16+10+2=36=4\cdot3\cdot3$. Đúng.}
\ans{$3$ mặt: $8$ khối; $2$ mặt: $16$ khối; $1$ mặt: $10$ khối; không sơn: $2$ khối.}
}

\begin{probbox}{37}{\relax}
Một hình chữ nhật có chu vi bằng chu vi của một hình vuông có diện tích $81\text{ cm}^2$. Chiều dài hình chữ nhật gấp đôi chiều rộng. Diện tích hình chữ nhật là
\choices{$81\text{ cm}^2$}{$72\text{ cm}^2$}{$54\text{ cm}^2$}{$36\text{ cm}^2$}{$64\text{ cm}^2$}
\end{probbox}
\probsol{
\textbf{Phân tích đề.} Chuỗi suy luận: diện tích hình vuông $\to$ cạnh hình vuông $\to$ chu vi (chung cho cả hai hình) $\to$ nửa chu vi của hình chữ nhật $\to$ chiều dài, chiều rộng $\to$ diện tích.
\begin{recallbox}Hình vuông cạnh $a$: $S=a\cdot a$, $P=4a$. Hình chữ nhật: $P=2\cdot(\text{dài}+\text{rộng})$, $S=\text{dài}\cdot\text{rộng}$.\end{recallbox}
\step{1}{Cạnh hình vuông}
Cần số $a$ với $a\cdot a=81$: $a=9$ cm (vì $9\cdot9=81$).

\step{2}{Chu vi}
Chu vi hình vuông $=4\cdot9=36$ cm, cũng là chu vi hình chữ nhật.

\step{3}{Tổng chiều dài và chiều rộng}
Nửa chu vi hình chữ nhật: $36:2=18$ cm $=$ dài $+$ rộng.

\step{4}{Chia theo tỉ lệ}
Chiều dài gấp đôi chiều rộng: rộng là $1$ phần, dài là $2$ phần, tổng $3$ phần ứng với $18$ cm.
\begin{barfigC}
  \rowlab{-0.2}{0}{rộng}
  \blk{0}{2}{0}{0.55}{tmBlue}{$6$}
  \rowlab{-0.2}{-0.8}{dài}
  \blk{0}{2}{-0.8}{0.55}{tmOrange}{$6$}
  \blk{2}{4}{-0.8}{0.55}{tmOrange}{$6$}
  \vbrace{4.2}{0.3}{-1.1}{tổng $18$ cm}
\end{barfigC}
Một phần $=18:3=6$ cm. Rộng $6$ cm, dài $12$ cm.

\step{5}{Diện tích}
$S=12\cdot6=72\ (\text{cm}^2)$.
\begin{tipbox}Hai hình có cùng chu vi $36$ cm, nhưng hình vuông có diện tích $81\text{ cm}^2$ lớn hơn hình chữ nhật ($72\text{ cm}^2$). Trong các hình chữ nhật cùng chu vi, hình vuông có diện tích lớn nhất.\end{tipbox}
\verify{Chu vi hình chữ nhật: $2\cdot(12+6)=36$ cm $=$ chu vi hình vuông; $12=2\cdot6$. Đúng.}
Vậy chọn \textbf{(B)}.
\ans{(B) $72\text{ cm}^2$}
}

\begin{probbox}{38}{\stars{1}}
Một vòng quay may mắn hình tròn được chia thành bốn ô có góc ở tâm lần lượt là $36^\circ$, $x^\circ$, $2x^\circ$ và $3x^\circ$ (hình vẽ).
\begin{enumerate}[label=\alph*),leftmargin=1.6em,itemsep=0pt,topsep=2pt]
  \item Tìm $x$.
  \item Mỗi ô chiếm bao nhiêu phần trăm diện tích vòng quay?
  \item Khi quay, kim \emph{không} dừng ở ô lớn nhất trong khoảng bao nhiêu phần trăm số lần?
\end{enumerate}
\begin{gmFig}[scale=0.9]
  \fill[tmPurple!20] (0,0)--(90:1.7) arc[start angle=90,end angle=126,radius=1.7]--cycle;
  \fill[tmBlue!20] (0,0)--(126:1.7) arc[start angle=126,end angle=180,radius=1.7]--cycle;
  \fill[tmGreen!20] (0,0)--(180:1.7) arc[start angle=180,end angle=288,radius=1.7]--cycle;
  \fill[tmOrange!25] (0,0)--(288:1.7) arc[start angle=288,end angle=450,radius=1.7]--cycle;
  \draw[Ink,thick] (0,0) circle (1.7);
  \foreach \a in {90,126,180,288}{\draw[Ink,thick] (0,0)--(\a:1.7);}
  \node[font=\footnotesize] at (108:1.2){$36^\circ$};
  \node[font=\footnotesize] at (153:1.1){$x^\circ$};
  \node[font=\footnotesize] at (234:1.0){$2x^\circ$};
  \node[font=\footnotesize] at (9:0.95){$3x^\circ$};
\end{gmFig}
\end{probbox}
\probsol{
\textbf{Phân tích đề.} Bốn góc ở tâm ghép lại thành một vòng trọn. Từ đó lập được một đẳng thức để tìm $x$, rồi đổi từng góc ra phần trăm.
\begin{recallbox}Các góc ở tâm quanh một điểm có tổng $360^\circ$. Ô có góc ở tâm $n^\circ$ chiếm $\dfrac{n}{360}$ hình tròn.\end{recallbox}
\step{1}{Câu a -- lập đẳng thức}
$36+x+2x+3x=360$, tức $36+6x=360$.
Bớt $36$ ở cả hai vế: $6x=324$. Chia cho $6$: $x=54$.

\step{2}{Câu b -- các góc và phần trăm}
Bốn góc: $36^\circ$, $54^\circ$, $108^\circ$, $162^\circ$.
\[ \frac{36}{360}=10\%,\qquad \frac{54}{360}=15\%,\qquad \frac{108}{360}=30\%,\qquad \frac{162}{360}=45\%. \]
(Mẹo tính: $\frac{36}{360}=\frac1{10}$; các góc còn lại gấp $1{,}5$; $3$; $4{,}5$ lần $36^\circ$, nên phần trăm lần lượt là $15\%$, $30\%$, $45\%$.)

\step{3}{Câu c -- dùng phần bù}
Ô lớn nhất là ô $162^\circ$, chiếm $45\%$. Kim \emph{không} dừng ở ô đó trong khoảng $100\%-45\%=55\%$ số lần.
\begin{warnbox}Kết quả ``$55\%$'' chỉ là \emph{khoảng} số lần sau rất nhiều lượt quay. Trong vài lượt quay, số lần thực tế có thể khác; tỉ lệ diện tích cho biết kim dừng ở đâu \emph{thường xuyên đến mức nào} về lâu dài.\end{warnbox}
\verify{$36+54+108+162=360$ và $10\%+15\%+30\%+45\%=100\%$.}
\ans{a) $x=54$. \ b) $10\%$, $15\%$, $30\%$, $45\%$. \ c) Khoảng $55\%$.}
}

\begin{probbox}{39}{\stars{1}}
Một hình chữ nhật dài $15$ cm, rộng $6$ cm. Trên cạnh trên, lấy hai đoạn ở hai đầu có cùng độ dài $a$ cm rồi cắt bỏ hai tam giác vuông ở hai góc trên, còn lại hình thang tô màu như hình. Biết diện tích hình thang là $72\text{ cm}^2$. Tìm $a$ và độ dài đáy nhỏ của hình thang.
\begin{gmFig}[scale=0.42]
  \fill[tmBlue!35] (0,0)--(15,0)--(12,6)--(3,6)--cycle;
  \draw[Ink,thick] (0,0) rectangle (15,6);
  \draw[Ink,thick] (0,0)--(3,6); \draw[Ink,thick] (15,0)--(12,6);
  \node[above] at (1.5,6){$a$}; \node[above] at (13.5,6){$a$};
  \node[below] at (7.5,0){$15$ cm}; \node[right] at (15,3){$6$ cm};
\end{gmFig}
\end{probbox}
\probsol{
\textbf{Phân tích đề.} Bài 14 cho kích thước và hỏi diện tích; bài này làm ngược lại: cho diện tích, hỏi kích thước. Ta viết diện tích hình thang theo $a$ bằng cách ``hình chữ nhật trừ hai tam giác'', rồi tìm $a$.
\begin{recallbox}$S_{\text{hình chữ nhật}}=$ dài $\times$ rộng; $\;S_{\text{tam giác vuông}}=\dfrac{\text{tích hai cạnh góc vuông}}{2}$.\end{recallbox}
\step{1}{Diện tích hình chữ nhật}
$15\cdot6=90\ (\text{cm}^2)$.

\step{2}{Tổng diện tích hai tam giác bị cắt}
Diện tích hai tam giác $=90-72=18\ (\text{cm}^2)$, nên mỗi tam giác có diện tích $18:2=9\ (\text{cm}^2)$.

\step{3}{Tìm $a$}
Mỗi tam giác vuông có hai cạnh góc vuông là $a$ và $6$:
\[ \frac{a\cdot6}{2}=9\quad\Rightarrow\quad 3\cdot a=9\quad\Rightarrow\quad a=3. \]

\step{4}{Đáy nhỏ của hình thang}
Đáy nhỏ $=15-a-a=15-3-3=9$ cm.
\verify{Theo công thức hình thang: $\dfrac{(15+9)\cdot6}{2}=\dfrac{24\cdot6}{2}=72\ (\text{cm}^2)$. Đúng.}
\begin{tipbox}Cũng có thể giải bằng công thức hình thang: đáy nhỏ $b$ thoả mãn $\dfrac{(15+b)\cdot 6}{2}=72$, tức $(15+b)\cdot3=72$, $15+b=24$, $b=9$; rồi $a=(15-9):2=3$.\end{tipbox}
\ans{$a=3$ cm; đáy nhỏ $9$ cm.}
}

\begin{probbox}{40}{\stars{1}}
Tam giác $ABC$ vuông tại $A$ có $AB=6$ cm, $AC=8$ cm và $BC=10$ cm.
\begin{enumerate}[label=\alph*),leftmargin=1.6em,itemsep=0pt,topsep=2pt]
  \item Tính diện tích tam giác $ABC$.
  \item Tính độ dài đường cao $AH$ kẻ từ $A$ xuống cạnh $BC$.
  \item Gọi $M$ là trung điểm của $BC$. Tính diện tích tam giác $ABM$.
\end{enumerate}
\begin{gmFig}[scale=0.55]
  \coordinate (B) at (0,0); \coordinate (C) at (10,0); \coordinate (A) at (3.6,4.8);
  \coordinate (H) at (3.6,0); \coordinate (M) at (5,0);
  \fill[tmBlue!10] (A)--(B)--(C)--cycle;
  \draw[Ink,thick] (A)--(B)--(C)--cycle;
  \draw[AnsRed,thick,dashed] (A)--(H); \draw[Ink] (H) rectangle ++(0.3,0.3);
  \draw[tmGreen!70!black,thick] (A)--(M);
  \draw[Ink] ($(A)+(-0.24,-0.32)$)--($(A)+(0.08,-0.56)$)--($(A)+(0.32,-0.24)$);
  \node[above] at (A){$A$}; \node[below left] at (B){$B$}; \node[below right] at (C){$C$};
  \node[below] at (H){$H$}; \node[below] at (M){$M$};
  \node[left] at (1.8,2.4){$6$}; \node[right] at (6.8,2.6){$8$};
\end{gmFig}
\end{probbox}
\probsol{
\textbf{Phân tích đề.} Một tam giác có thể chọn \emph{bất kì} cạnh nào làm đáy, miễn là dùng đúng chiều cao ứng với đáy đó. Diện tích luôn như nhau. Ta tính diện tích theo cách dễ (đáy $AB$, chiều cao $AC$), rồi dùng nó để tìm chiều cao ứng với đáy $BC$.
\begin{recallbox}$S=\dfrac{\text{đáy}\cdot\text{chiều cao}}{2}$, suy ra chiều cao $=\dfrac{2\cdot S}{\text{đáy}}$. Tam giác vuông: hai cạnh góc vuông là đáy và chiều cao của nhau.\end{recallbox}
\step{1}{Câu a -- chọn đáy $AB$}
Vì góc $A$ vuông nên $AC$ vuông góc với $AB$: lấy $AB$ làm đáy thì $AC$ là chiều cao.
\[ S_{ABC}=\frac{6\cdot8}{2}=24\ (\text{cm}^2). \]

\step{2}{Câu b -- chọn đáy $BC$}
Lấy $BC=10$ cm làm đáy, chiều cao tương ứng là $AH$. Diện tích vẫn là $24\text{ cm}^2$:
\[ \frac{10\cdot AH}{2}=24\quad\Rightarrow\quad 10\cdot AH=48\quad\Rightarrow\quad AH=4{,}8\text{ cm}. \]

\step{3}{Câu c -- trung điểm chia đôi diện tích}
\begin{recallbox}Hai tam giác có đáy bằng nhau và chung chiều cao thì có diện tích bằng nhau.\end{recallbox}
Hai tam giác $ABM$ và $ACM$ có đáy $BM=MC=5$ cm (vì $M$ là trung điểm) và chung chiều cao $AH$ (cùng kẻ từ $A$ xuống đường thẳng $BC$). Vậy chúng có diện tích bằng nhau, mỗi tam giác bằng nửa tam giác $ABC$:
\[ S_{ABM}=\frac{5\cdot4{,}8}{2}=12\ (\text{cm}^2)=\frac{24}{2}. \]
\begin{warnbox}Đừng lấy ``đáy $BC$ nhân cạnh $AB$'' -- cạnh $AB$ không vuông góc với $BC$, nên không phải chiều cao ứng với đáy $BC$.\end{warnbox}
\verify{$AH=4{,}8$ nhỏ hơn cả $AB=6$ và $AC=8$ -- hợp lí, vì đường vuông góc là đường ngắn nhất từ $A$ đến đường thẳng $BC$. Và $\frac{10\cdot4{,}8}{2}=24$.}
\ans{a) $24\text{ cm}^2$. \ b) $AH=4{,}8$ cm. \ c) $12\text{ cm}^2$.}
}

% ---------------- part d ----------------
% ===== Topic 2, part d: problems 41-50 (VI) =====
\providecommand{\tbGrid}[4]{\begin{tikzpicture}[baseline=0.2cm,x=0.5cm,y=0.5cm]
  \draw[Ink,line width=0.7pt] (0,0) grid[step=1] (2,2);
  \node[font=\sffamily\bfseries,text=tmBlue!80!black,inner sep=0pt,#4] at (#1,#2) {#3};\end{tikzpicture}}
% helper inside a tikzpicture: \tbG{x0}{y0}{cx}{cy}{letter}{transform} draws a 2x2 grid (cell 0.5) at (x0,y0)
\providecommand{\tbG}[6]{\draw[Ink,line width=0.7pt] (#1,#2) grid[step=0.5] ++(1,1);
  \node[font=\sffamily\bfseries,text=tmBlue!80!black,inner sep=0pt,#6] at ({#1+#3},{#2+#4}) {#5};}

% ---------------------------------------------------------------- 41
\begin{probbox}{41}{\stars{1}}
\textit{(Mở rộng Bài 16.)} Ta chỉ xét các hình chữ nhật có độ dài hai cạnh là các \emph{số nguyên dương}
(đơn vị cm); hình vuông cũng được coi là hình chữ nhật.
\begin{enumerate}[label=\alph*),leftmargin=1.6em,itemsep=1pt]
\item Tìm tất cả các hình chữ nhật mà số đo diện tích (cm$^2$) \emph{bằng} số đo chu vi (cm).
\item Tìm tất cả các hình chữ nhật mà số đo diện tích \emph{gấp đôi} số đo chu vi.
\end{enumerate}
\end{probbox}
\probsol{
\textbf{Phân tích đề.} Gọi hai cạnh là $a$ và $b$ (số nguyên dương). Không mất tính tổng quát, giả sử
$a\le b$ ($a$ là cạnh ngắn). Diện tích là $ab$, chu vi là $2(a+b)$. Đây là bài toán ``tìm nghiệm nguyên'':
một phương trình có hai ẩn. Mẹo: dùng $a\le b$ để \emph{chặn} $a$ trong một khoảng nhỏ, rồi thử từng giá trị.

\textbf{a)} Điều kiện: $ab=2a+2b$.

\step{1}{Chặn trên cho $a$}
Vì $a\le b$ nên $2a+2b\le 2b+2b=4b$. Suy ra $ab\le 4b$. Chia hai vế cho $b>0$: $a\le 4$.

\step{2}{Chặn dưới cho $a$}
Vì $2a>0$ nên $ab=2a+2b>2b$. Chia hai vế cho $b>0$: $a>2$. Vậy $a=3$ hoặc $a=4$.

\step{3}{Thử từng giá trị}
\begin{itemize}[leftmargin=1.4em,itemsep=1pt]
\item $a=3$: $3b=6+2b\Rightarrow b=6$. Hình chữ nhật $3\times 6$.
\item $a=4$: $4b=8+2b\Rightarrow 2b=8\Rightarrow b=4$. Hình vuông $4\times 4$.
\end{itemize}
\begin{figC}\begin{scope}[x=0.45cm,y=0.45cm]
\fill[tmBlue!15] (0,0) rectangle (6,3);\draw[Line] (0,0) grid[step=1] (6,3);\draw[tmBlue!80!black,line width=1pt] (0,0) rectangle (6,3);
\node[below,font=\footnotesize] at (3,0){$S=18$, $C=18$};
\begin{scope}[xshift=4cm]
\fill[tmGreen!15] (0,0) rectangle (4,4);\draw[Line] (0,0) grid[step=1] (4,4);\draw[tmGreen!70!black,line width=1pt] (0,0) rectangle (4,4);
\node[below,font=\footnotesize] at (2,0){$S=16$, $C=16$};
\end{scope}
\end{scope}\end{figC}

\textbf{b)} Điều kiện: $ab=2\cdot 2(a+b)=4a+4b$.

\step{4}{Chặn $a$ như trên}
$ab=4a+4b\le 8b\Rightarrow a\le 8$; và $ab=4a+4b>4b\Rightarrow a>4$. Vậy $a\in\{5;6;7;8\}$.

\step{5}{Thử từng giá trị}
Từ $ab=4a+4b$ ta có $ab-4b=4a$, tức $b(a-4)=4a$, nên $b=\dfrac{4a}{a-4}$.
\begin{center}\begin{tabular}{|c|c|c|c|c|}\hline
$a$ & $5$ & $6$ & $7$ & $8$\\\hline
$b=\dfrac{4a}{a-4}$ & $\dfrac{20}{1}=20$ & $\dfrac{24}{2}=12$ & $\dfrac{28}{3}$ (không nguyên) & $\dfrac{32}{4}=8$\\\hline
\end{tabular}\end{center}
Có ba hình: $5\times 20$, $6\times 12$, $8\times 8$.

\begin{proofbox}Vì sao không sót? Mọi hình chữ nhật thỏa mãn đều có cạnh ngắn $a$ nằm trong khoảng đã chặn,
và với mỗi $a$, cạnh $b$ được xác định duy nhất từ phương trình. Ta đã thử hết các $a$ đó.\end{proofbox}
\begin{tipbox}Cách khác cho câu a): $ab-2a-2b+4=4$, tức là $(a-2)(b-2)=4=1\cdot 4=2\cdot 2$. Suy ra
$a-2=1,\ b-2=4$ hoặc $a-2=b-2=2$. Lại được $3\times 6$ và $4\times 4$.\end{tipbox}
\verify{$3\times 6$: $S=18$, $C=2\cdot 9=18$. $4\times4$: $16=16$. $5\times 20$: $S=100$, $C=50$;
$6\times 12$: $S=72$, $C=36$; $8\times 8$: $S=64$, $C=32$ -- đều gấp đôi.}
\ans{a) $3\times 6$ và $4\times 4$. \quad b) $5\times 20$, $6\times 12$, $8\times 8$.}
}

% ---------------------------------------------------------------- 42
\begin{probbox}{42}{\stars{1}}
\textit{(Mở rộng Bài 17.)} Chữ P được viết ở ô trên -- phải của lưới $2\times 2$: \tbGrid{1.5}{1.5}{P}{}\,.
Ta được dùng ba phép biến hình, bao nhiêu lần và theo thứ tự nào cũng được:
\begin{itemize}[leftmargin=1.4em,itemsep=1pt]
\item $Q$: quay lưới $90^\circ$ theo chiều kim đồng hồ quanh tâm;
\item $D$: lật lưới qua đường thẳng \emph{đứng} đi qua tâm;
\item $N$: lật lưới qua đường thẳng \emph{nằm ngang} đi qua tâm.
\end{itemize}
\begin{enumerate}[label=\alph*),leftmargin=1.6em,itemsep=1pt]
\item Vẽ kết quả khi thực hiện $D$ rồi $N$. Có thể nhận được kết quả đó chỉ bằng các phép $Q$ không?
\item Chứng tỏ rằng thực hiện $D$ rồi $Q$ cho cùng kết quả với một phép lật qua một \emph{đường chéo} của lưới.
Đó là đường chéo nào?
\item Có tất cả bao nhiêu hình khác nhau có thể nhận được (kể cả hình ban đầu)? Hãy giải thích vì sao không
thể nhiều hơn.
\end{enumerate}
\end{probbox}
\probsol{
\textbf{Phân tích đề.} Như trong Bài 17, ta coi lưới là một \emph{tấm thẻ} vuông trong suốt, có vẽ chữ P. Mỗi
phép biến hình là một cách nhấc tấm thẻ lên rồi đặt lại vừa khít chỗ cũ. Để theo dõi, ta ghi nhớ hai điều:
chữ P nằm ở ô nào, và nó bị quay/lật ra sao.

\textbf{a)}
\step{1}{Thực hiện $D$}
Lật qua đường đứng: ô trên -- phải sang ô trên -- trái, chữ P bị ``soi gương'' (bụng quay sang trái).
\step{2}{Thực hiện tiếp $N$}
Lật qua đường ngang: ô trên -- trái xuống ô dưới -- trái, chữ bị lật ngược trên -- dưới. Kết quả: chữ P ở ô
dưới -- trái, \emph{lộn đầu} (giống chữ ``d''), như hình thứ ba bên dưới.
\begin{figC}
\tbG{0}{0}{0.75}{0.75}{P}{}\node at (1.75,0.5){$\xrightarrow{\;D\;}$};
\tbG{2.5}{0}{0.25}{0.75}{P}{xscale=-1}\node at (4.25,0.5){$\xrightarrow{\;N\;}$};
\tbG{5}{0}{0.25}{0.25}{P}{rotate=180}
\node[anchor=west] at (6.2,0.5){\small$=$ quay $180^\circ$};
\end{figC}
\step{3}{So sánh với các phép $Q$}
Quay $Q$ hai lần (tức quay $180^\circ$): ô trên -- phải đi tới ô đối diện qua tâm là ô dưới -- trái, và chữ bị
lộn đầu. Trùng với kết quả trên. Vậy \textbf{$D$ rồi $N$ cho cùng kết quả với $Q$ rồi $Q$.}
\begin{recallbox}Hai lần lật qua hai đường vuông góc với nhau (cùng đi qua tâm) cho kết quả bằng một phép quay
$180^\circ$ quanh tâm.\end{recallbox}

\textbf{b)}
\step{4}{Thực hiện $D$ rồi $Q$}
$D$: chữ P sang ô trên -- trái, bị soi gương. $Q$ (quay $90^\circ$ theo chiều kim đồng hồ): ô trên -- trái sang ô
trên -- phải. Kết quả: P \emph{trở về ô trên -- phải} nhưng thân nằm ngang, bụng ở phía trên:
\begin{figC}
\tbG{0}{0}{0.75}{0.75}{P}{}\node at (1.75,0.5){$\xrightarrow{\;D\;}$};
\tbG{2.5}{0}{0.25}{0.75}{P}{xscale=-1}\node at (4.25,0.5){$\xrightarrow{\;Q\;}$};
\tbG{5}{0}{0.75}{0.75}{P}{cm={0,1,1,0,(0,0)}}\draw[tmOrange,dashed,line width=0.9pt] (4.85,-0.15)--(6.15,1.15);
\draw[tmOrange,dashed,line width=0.9pt] (4.05,-0.5)--(5.15,0.6);
\end{figC}
\step{5}{Nhận ra phép lật qua đường chéo}
Ô trên -- phải không đổi chỗ, nghĩa là nó nằm trên trục lật. Đường chéo đi qua ô trên -- phải (và ô dưới -- trái)
là đường chéo nối \textbf{góc dưới -- trái với góc trên -- phải} (nét đứt). Lật qua đường này: hướng ``lên''
thành hướng ``sang phải'' và ngược lại, nên thân P (thẳng đứng) thành nằm ngang, bụng P (bên phải) thành ở
phía trên -- đúng như hình. Có thể kiểm tra thêm: một chữ đặt ở ô trên -- trái sẽ chuyển sang ô dưới -- phải
trong cả hai cách.

\textbf{c)}
\step{6}{Không thể có quá $8$ hình}
\begin{proofbox}Đánh dấu góc trên -- phải của tấm thẻ bằng một ngôi sao. Mỗi lần đặt lại, ngôi sao chỉ có thể
nằm ở $1$ trong $4$ góc của khung. Ngoài ra tấm thẻ chỉ có thể đang \emph{ngửa} hoặc \emph{úp} (mỗi lần lật
$D$ hay $N$ đổi ngửa thành úp và ngược lại; quay $Q$ không đổi). Khi đã biết vị trí ngôi sao và thẻ ngửa hay
úp thì vị trí của tấm thẻ, và do đó cả hình, được xác định hoàn toàn. Vậy có nhiều nhất $4\times 2=8$
hình.\end{proofbox}
\step{7}{Có đủ $8$ hình}
Bốn hình ``thẻ ngửa'': quay $0^\circ$, $90^\circ$, $180^\circ$, $270^\circ$ (dùng $Q$ từ $0$ đến $3$ lần).
Bốn hình ``thẻ úp'': mỗi hình trên thực hiện thêm $D$. Chữ P không có trục đối xứng nào, nên $8$ hình này
khác nhau từng đôi một:
\begin{figC}
\tbG{0}{0}{0.75}{0.75}{P}{}\tbG{1.5}{0}{0.75}{0.25}{P}{rotate=-90}
\tbG{3}{0}{0.25}{0.25}{P}{rotate=180}\tbG{4.5}{0}{0.25}{0.75}{P}{rotate=90}
\tbG{0}{-1.5}{0.25}{0.75}{P}{xscale=-1}\tbG{1.5}{-1.5}{0.25}{0.25}{P}{cm={0,-1,-1,0,(0,0)}}
\tbG{3}{-1.5}{0.75}{0.25}{P}{yscale=-1}\tbG{4.5}{-1.5}{0.75}{0.75}{P}{cm={0,1,1,0,(0,0)}}
\node[anchor=east,font=\footnotesize] at (-0.2,0.5){thẻ ngửa};\node[anchor=east,font=\footnotesize] at (-0.2,-1){thẻ úp};
\end{figC}
\begin{warnbox}Đừng đếm ``$4$ ô $\times$ $8$ tư thế $=32$''. Vị trí ô và tư thế của chữ \emph{gắn liền} với
nhau, vì chúng đi theo cùng một tấm thẻ.\end{warnbox}
\ans{a) P ở ô dưới -- trái, lộn đầu; bằng $Q$ rồi $Q$ (quay $180^\circ$). \ b) Đường chéo nối góc dưới -- trái
với góc trên -- phải. \ c) $8$ hình.}
}

% ---------------------------------------------------------------- 43
\begin{probbox}{43}{\stars{1}}
\textit{(Mở rộng Bài 18.)} Cho hình chữ nhật $ABCD$ với $AB=12$ cm, $BC=8$ cm. Điểm $P$ nằm bên trong hình chữ nhật.
\begin{figC}\begin{scope}[x=0.36cm,y=0.36cm]
\fill[tmBlue!25] (0,0)--(12,0)--(7,3)--cycle;\fill[tmOrange!25] (12,8)--(0,8)--(7,3)--cycle;
\draw[Ink,line width=1pt] (0,0) rectangle (12,8);
\draw[Ink] (0,0)--(7,3)--(12,0) (0,8)--(7,3)--(12,8);
\fill[Ink] (7,3) circle (2.5pt);\node[right=2pt] at (7,3){$P$};
\node[below left] at (0,0){$A$};\node[below right] at (12,0){$B$};\node[above right] at (12,8){$C$};\node[above left] at (0,8){$D$};
\end{scope}\end{figC}
\begin{enumerate}[label=\alph*),leftmargin=1.6em,itemsep=1pt]
\item Chứng minh rằng dù $P$ ở đâu, ta luôn có $S_{PAB}+S_{PCD}=48$ cm$^2$.
\item Biết $S_{PAB}=18$ cm$^2$ và $S_{PBC}=20$ cm$^2$. Tính $S_{PCD}$, $S_{PDA}$ và khoảng cách từ $P$ đến
bốn cạnh của hình chữ nhật.
\item Có điểm $P$ nào bên trong hình chữ nhật mà $S_{PAB}=20$, $S_{PBC}=30$ và $S_{PCD}=30$ (cm$^2$) không?
\end{enumerate}
\end{probbox}
\probsol{
\textbf{Phân tích đề.} Cả bốn tam giác đều có đáy là một cạnh hình chữ nhật và đỉnh là $P$. Chiều cao của
mỗi tam giác chính là \emph{khoảng cách từ $P$ đến cạnh đó}. Ta gọi khoảng cách từ $P$ đến $AB$, $BC$, $CD$,
$DA$ lần lượt là $h_1,h_2,h_3,h_4$.
\begin{figC}\begin{scope}[x=0.36cm,y=0.36cm]
\draw[Ink,line width=1pt] (0,0) rectangle (12,8);
\draw[tmBlue,thick,<->] (7,0.1)--(7,2.9);\draw[tmOrange!80!black,thick,<->] (7,3.1)--(7,7.9);
\draw[tmGreen!70!black,thick,<->] (7.1,3)--(11.9,3);\draw[tmPurple,thick,<->] (0.1,3)--(6.9,3);
\fill[Ink] (7,3) circle (2.5pt);
\node[right,font=\footnotesize] at (7,1.5){$h_1$};\node[right,font=\footnotesize] at (7,5.5){$h_3$};
\node[above,font=\footnotesize] at (9.5,3){$h_2$};\node[above,font=\footnotesize] at (3.5,3){$h_4$};
\node[below left] at (0,0){$A$};\node[below right] at (12,0){$B$};\node[above right] at (12,8){$C$};\node[above left] at (0,8){$D$};
\end{scope}\end{figC}

\textbf{a)}
\step{1}{Tổng hai chiều cao}
Kẻ qua $P$ đường thẳng vuông góc với $AB$; vì $AB\parallel CD$ nên nó cũng vuông góc với $CD$. Đoạn nằm giữa
$AB$ và $CD$ có độ dài bằng $BC=8$, và $P$ chia đoạn này thành hai phần $h_1$ và $h_3$. Vậy $h_1+h_3=8$.
Tương tự $h_2+h_4=AB=12$.
\step{2}{Cộng hai diện tích}
\begin{recallbox}Diện tích tam giác $=\tfrac12\times$ đáy $\times$ chiều cao.\end{recallbox}
\[ S_{PAB}+S_{PCD}=\frac{12\cdot h_1}{2}+\frac{12\cdot h_3}{2}=6(h_1+h_3)=6\cdot 8=48 \text{ (cm}^2). \]
Kết quả không phụ thuộc vào vị trí của $P$. (Đó đúng bằng nửa diện tích hình chữ nhật $12\cdot 8=96$.)
Tương tự, $S_{PBC}+S_{PDA}=\tfrac{8}{2}(h_2+h_4)=4\cdot 12=48$ cm$^2$.

\textbf{b)}
\step{3}{Dùng kết quả câu a}
$S_{PCD}=48-18=30$ cm$^2$, \quad $S_{PDA}=48-20=28$ cm$^2$.
\step{4}{Tính bốn khoảng cách}
Từ $S=\tfrac12\cdot$ đáy $\cdot h$ suy ra $h=2S:$ đáy.
\[ h_1=\frac{2\cdot 18}{12}=3,\quad h_2=\frac{2\cdot 20}{8}=5,\quad h_3=\frac{2\cdot 30}{12}=5,\quad h_4=\frac{2\cdot 28}{8}=7 \text{ (cm)}. \]

\textbf{c)}
\step{5}{Kiểm tra điều kiện bắt buộc}
Theo câu a), luôn có $S_{PAB}+S_{PCD}=48$. Nhưng $20+30=50\ne 48$. Vậy \textbf{không có} điểm $P$ nào như thế.

\begin{tipbox}Bài 18 là trường hợp đặc biệt: $P$ là tâm hình vuông nên bốn tam giác bằng nhau. Ở đây $P$ bất kì
nhưng ``tổng hai tam giác đối diện'' vẫn luôn bằng nửa hình.\end{tipbox}
\verify{$h_1+h_3=3+5=8=BC$ và $h_2+h_4=5+7=12=AB$. Đặt $A(0;0)$ thì $P(7;3)$ -- đúng là điểm trong hình vẽ.}
\ans{b) $S_{PCD}=30$ cm$^2$, $S_{PDA}=28$ cm$^2$; khoảng cách đến $AB,BC,CD,DA$ là $3;5;5;7$ cm. \ c) Không có.}
}

% ---------------------------------------------------------------- 44
\begin{probbox}{44}{\stars{1}}
\textit{(Mở rộng Bài 22.)} Hình bình hành $ABCD$ có diện tích $48$ cm$^2$.
\begin{enumerate}[label=\alph*),leftmargin=1.6em,itemsep=1pt]
\item Gọi $M$, $N$ lần lượt là trung điểm của $BC$ và $CD$. Tính diện tích tam giác $AMN$.
\item Giữ $N$ là trung điểm của $CD$, nhưng lấy $M$ trên cạnh $BC$ sao cho $BM=3\,MC$. Tính diện tích tam giác $AMN$.
\end{enumerate}
\begin{figC}\begin{scope}[x=0.55cm,y=0.55cm]
\fill[tmBlue!25] (0,0)--(7,2)--(5,4)--cycle;
\draw[Ink,line width=1pt] (0,0)--(6,0)--(8,4)--(2,4)--cycle;
\draw[tmBlue!80!black,line width=1pt] (0,0)--(7,2)--(5,4)--cycle;
\node[below left] at (0,0){$A$};\node[below right] at (6,0){$B$};\node[above right] at (8,4){$C$};\node[above left] at (2,4){$D$};
\node[right] at (7,2){$M$};\node[above] at (5,4){$N$};
\end{scope}\end{figC}
\end{probbox}
\probsol{
\textbf{Phân tích đề.} Tam giác $AMN$ nằm ``chéo'' bên trong hình bình hành. Ta lấy cả hình bình hành bỏ đi ba
tam giác ở các góc: $ABM$, $MCN$ và $NDA$.
\[ S_{AMN}=S_{ABCD}-S_{ABM}-S_{MCN}-S_{NDA}. \]
\begin{recallbox}(1) Mỗi đường chéo chia hình bình hành thành hai tam giác có diện tích bằng nhau ($=24$ cm$^2$).
(2) Hai tam giác chung đỉnh, có hai đáy nằm trên cùng một đường thẳng thì tỉ số diện tích bằng tỉ số hai đáy.
\end{recallbox}

\textbf{a)}
\step{1}{Tam giác $ABM$}
So với tam giác $ABC$ ($=24$): chung đỉnh $A$, đáy $BM=\tfrac12 BC$. Vậy $S_{ABM}=\tfrac12\cdot 24=12$.
\step{2}{Tam giác $NDA$}
So với tam giác $ACD$ ($=24$): chung đỉnh $A$, đáy $DN=\tfrac12 DC$. Vậy $S_{NDA}=12$.
\step{3}{Tam giác $MCN$ (làm hai bước, như Bài 22)}
Tam giác $BCD$ có diện tích $24$. Tam giác $MCD$ chung đỉnh $D$ với nó, đáy $CM=\tfrac12 CB$: $S_{MCD}=12$.
Tam giác $MCN$ chung đỉnh $M$ với tam giác $MCD$, đáy $CN=\tfrac12 CD$: $S_{MCN}=\tfrac12\cdot 12=6$.
\step{4}{Kết quả}
$S_{AMN}=48-12-6-12=18$ cm$^2$.

\textbf{b)}
\step{5}{Làm lại với $MC=\tfrac14 BC$}
Vì $BM=3\,MC$ nên $BC=4\,MC$, tức $MC=\tfrac14 BC$ và $BM=\tfrac34 BC$.
\begin{itemize}[leftmargin=1.4em,itemsep=1pt]
\item $S_{ABM}=\tfrac34\cdot 24=18$;
\item $S_{NDA}=12$ (không đổi);
\item $S_{MCD}=\tfrac14\cdot 24=6$, rồi $S_{MCN}=\tfrac12\cdot 6=3$.
\end{itemize}
$S_{AMN}=48-18-3-12=15$ cm$^2$.

\begin{tipbox}Quy tắc ``góc'': nếu trên hai cạnh của một tam giác, tính từ cùng một đỉnh, ta lấy các phần
bằng $\tfrac{p}{q}$ và $\tfrac{r}{s}$ của cạnh, thì tam giác nhỏ ở góc đó có diện tích bằng
$\tfrac{p}{q}\cdot\tfrac{r}{s}$ tam giác lớn. Ở câu b): $\tfrac14\cdot\tfrac12=\tfrac18$ của $24$ là $3$.\end{tipbox}
\verify{Đặt $A(0;0)$, $B(8;0)$, $D(2;6)$, $C(10;6)$ (diện tích $8\cdot 6=48$). Câu a): $M(9;3)$, $N(6;6)$; bao tam giác
$AMN$ bằng hình chữ nhật $9\times 6=54$, trừ $\tfrac{9\cdot3}{2}=13{,}5$, $\tfrac{3\cdot 3}{2}=4{,}5$,
$\tfrac{6\cdot6}{2}=18$: còn $18$. Đúng.}
\ans{a) $18$ cm$^2$. \quad b) $15$ cm$^2$.}
}

% ---------------------------------------------------------------- 45
\begin{probbox}{45}{\stars{1}}
\textit{(Mở rộng Bài 21.)} Xét các tam giác có độ dài ba cạnh là các số nguyên dương và chu vi bằng $15$
(hai tam giác có cùng bộ ba độ dài cạnh chỉ tính là một).
\begin{enumerate}[label=\alph*),leftmargin=1.6em,itemsep=1pt]
\item Gọi $c$ là cạnh dài nhất. Chứng minh rằng $5\le c\le 7$.
\item Có tất cả bao nhiêu tam giác như vậy? Hãy liệt kê.
\item Trong đó có bao nhiêu tam giác cân? Bao nhiêu tam giác có \emph{đúng} hai cạnh bằng nhau?
\end{enumerate}
\end{probbox}
\probsol{
\textbf{Phân tích đề.} Gọi ba cạnh là $a\le b\le c$ (sắp xếp từ nhỏ đến lớn để không đếm trùng). Ta có
$a+b+c=15$, và theo bất đẳng thức tam giác chỉ cần kiểm tra $a+b>c$.
\begin{recallbox}Ba đoạn $a\le b\le c$ là ba cạnh của một tam giác khi và chỉ khi $a+b>c$ (tổng hai cạnh nhỏ
lớn hơn cạnh lớn nhất).\end{recallbox}

\textbf{a)}
\step{1}{Chặn trên}
$a+b=15-c$. Điều kiện $a+b>c$ trở thành $15-c>c$, tức $2c<15$, $c<7{,}5$. Vì $c$ nguyên nên $c\le 7$.
\step{2}{Chặn dưới}
$c$ là cạnh lớn nhất nên $a\le c$, $b\le c$, suy ra $15=a+b+c\le 3c$, tức $c\ge 5$.

\textbf{b)}
\step{3}{Liệt kê theo $c$}
Với mỗi $c$, ta có $a+b=15-c$ và $a\le b\le c$.
\begin{center}\begin{tabular}{|c|c|l|c|}\hline
$c$ & $a+b$ & Các cặp $(a;b)$ với $a\le b\le c$ & Số tam giác\\\hline
$7$ & $8$ & $(1;7),\ (2;6),\ (3;5),\ (4;4)$ & $4$\\
$6$ & $9$ & $(3;6),\ (4;5)$ \quad (không lấy $(2;7)$ vì $7>6$) & $2$\\
$5$ & $10$ & $(5;5)$ & $1$\\\hline
\end{tabular}\end{center}
Tổng cộng $4+2+1=7$ tam giác:
\[ (1;7;7),\ (2;6;7),\ (3;5;7),\ (4;4;7),\ (3;6;6),\ (4;5;6),\ (5;5;5). \]
\begin{warnbox}Với $c=7$, cặp $(1;7)$ cho tam giác $1;7;7$ -- rất ``gầy'' nhưng vẫn hợp lệ vì $1+7=8>7$.
Ngược lại, không có tam giác nào có cạnh $8$ vì khi đó hai cạnh còn lại có tổng $7<8$.\end{warnbox}

\textbf{c)}
\step{4}{Chọn ra các tam giác cân}
Tam giác cân có ít nhất hai cạnh bằng nhau: $(1;7;7)$, $(4;4;7)$, $(3;6;6)$, $(5;5;5)$ -- có $4$ tam giác.
Trong đó $(5;5;5)$ là tam giác đều (ba cạnh bằng nhau), nên có đúng hai cạnh bằng nhau là $3$ tam giác.
\verify{Mỗi bộ đều có tổng $15$ và thỏa $a+b>c$: $8>7$, $8>7$, $8>7$, $8>7$, $9>6$, $9>6$, $10>5$. Đúng.}
\ans{a) $5\le c\le 7$. \ b) $7$ tam giác. \ c) $4$ tam giác cân, trong đó $3$ tam giác có đúng hai cạnh bằng nhau.}
}

% ---------------------------------------------------------------- 46
\begin{probbox}{46}{\stars{2}}
\textit{(Mở rộng Bài 23.)} Trên lưới ô vuông (mỗi ô cạnh $1$), cho tam giác $ABC$ với $A(0;0)$, $B(12;0)$,
$C(3;9)$. Có đúng $43$ giao điểm của các đường kẻ lưới nằm \emph{bên trong} (không nằm trên cạnh) tam giác.
Với mỗi giao điểm $P$ như vậy, ta được ba tam giác $PAB$, $PBC$, $PCA$.
\begin{figC}\begin{scope}[x=0.36cm,y=0.36cm]
\draw[Line,line width=0.5pt] (0,0) grid[step=1] (12,9);
\draw[Ink,line width=1.3pt] (0,0)--(12,0)--(3,9)--cycle;
\node[below left] at (0,0){$A$};\node[below right] at (12,0){$B$};\node[above] at (3,9){$C$};
\end{scope}\end{figC}
\begin{enumerate}[label=\alph*),leftmargin=1.6em,itemsep=1pt]
\item Tính diện tích tam giác $ABC$.
\item Tìm điểm $P$ sao cho ba tam giác $PAB$, $PBC$, $PCA$ có diện tích bằng nhau.
\item Trong tất cả $3\cdot 43=129$ tam giác, có bao nhiêu tam giác có diện tích đúng bằng $\tfrac13$ diện tích
tam giác $ABC$? Có bao nhiêu vị trí của $P$ cho ít nhất một tam giác như thế?
\end{enumerate}
\end{probbox}
\probsol{
\textbf{Phân tích đề.} Đây là ``phiên bản một phần ba'' của Bài 23. Ý chính vẫn là: tam giác có đáy cố định thì
diện tích chỉ phụ thuộc vào chiều cao, và các điểm có cùng chiều cao nằm trên một đường thẳng song song với đáy.

\textbf{a)}
\step{1}{Diện tích tam giác $ABC$}
Đáy $AB=12$ nằm trên trục hoành, chiều cao là tung độ của $C$, bằng $9$:
$S_{ABC}=\tfrac{12\cdot 9}{2}=54$. Một phần ba là $18$.

\textbf{c)} (làm trước, câu b) suy ra ngay)
\step{2}{Tam giác $PAB$ có diện tích $18$}
Đáy $AB=12$ nên chiều cao phải là $2\cdot 18:12=3$, tức $P$ nằm trên đường ngang $y=3$.
Trên đường này: cạnh $AC$ (đi ``sang phải $1$, lên $3$'') qua điểm $(1;3)$; cạnh $BC$ (đi ``sang trái $1$,
lên $1$'' từ $B$) qua điểm $(9;3)$. Các điểm bên trong: $x=2,3,\dots,8$, tức $7$ điểm.

\step{3}{Tam giác $PBC$ có diện tích $18$}
\begin{proofbox}Lấy $E$ trên $AB$ với $EB=\tfrac13 AB=4$, tức $E(8;0)$. Tam giác $EBC$ và $ABC$ chung đỉnh $C$,
đáy $EB=\tfrac13 AB$ nên $S_{EBC}=\tfrac13\cdot 54=18$. Do đó $S_{PBC}=18$ khi và chỉ khi $P$ cách $BC$ đúng
bằng $E$, tức $P$ nằm trên đường thẳng qua $E$ song song với $BC$ (về phía trong tam giác).\end{proofbox}
$BC$ có hướng ``sang trái $1$, lên $1$''. Từ $E(8;0)$ đi theo hướng đó: $(7;1)$, $(6;2)$, $(5;3)$, $(4;4)$, $(3;5)$,
rồi $(2;6)$ -- điểm này nằm trên $AC$ (vì trên $AC$, tung độ gấp $3$ lần hoành độ), không tính. Có $5$ điểm.

\step{4}{Tam giác $PCA$ có diện tích $18$}
Tương tự, lấy $F$ trên $AB$ với $AF=4$, tức $F(4;0)$; $P$ phải nằm trên đường thẳng qua $F$ song song với $AC$
(hướng ``sang phải $1$, lên $3$''): $(5;3)$, rồi $(6;6)$ -- nằm trên $BC$ (vì $6+6=12$), không tính. Có $1$ điểm.
\begin{figC}\begin{scope}[x=0.36cm,y=0.36cm]
\draw[Line,line width=0.5pt] (0,0) grid[step=1] (12,9);
\draw[Ink,line width=1.3pt] (0,0)--(12,0)--(3,9)--cycle;
\draw[tmBlue,line width=1pt] (1,3)--(9,3);
\draw[tmOrange!90!black,line width=1pt] (8,0)--(2,6);
\draw[tmGreen!70!black,line width=1pt] (4,0)--(6,6);
\foreach \p in {(2,3),(3,3),(4,3),(6,3),(7,3),(8,3)}{\fill[tmBlue] \p circle (3pt);}
\foreach \p in {(7,1),(6,2),(4,4),(3,5)}{\fill[tmOrange!90!black] \p circle (3pt);}
\fill[AnsRed] (5,3) circle (4pt);
\node[below left] at (0,0){$A$};\node[below right] at (12,0){$B$};\node[above] at (3,9){$C$};
\node[below] at (8,0){$E$};\node[below] at (4,0){$F$};
\end{scope}\end{figC}

\step{5}{Đếm}
Số tam giác: $7+5+1=13$. Điểm $(5;3)$ nằm trên cả ba đường (điểm đỏ), nên được đếm ba lần. Số vị trí khác
nhau của $P$: $13-2=11$.

\textbf{b)}
\step{6}{Ba tam giác bằng nhau}
Ba tam giác có tổng diện tích bằng $54$, nên bằng nhau khi mỗi tam giác bằng $18$, tức $P$ phải nằm trên
\emph{cả ba} đường thẳng vừa tìm. Chỉ có một điểm như vậy: $P(5;3)$.

\begin{tipbox}Điểm $(5;3)$ chính là \emph{trọng tâm} của tam giác: hoành độ $\tfrac{0+12+3}{3}=5$, tung độ
$\tfrac{0+0+9}{3}=3$. Các em sẽ học kĩ về trọng tâm trong chương tam giác.\end{tipbox}
\verify{Với $P(5;3)$: $S_{PAB}=\tfrac{12\cdot 3}{2}=18$. $S_{PCA}$: bao bằng hình chữ nhật $5\times 9=45$, trừ
$\tfrac{3\cdot 9}{2}=13{,}5$ (góc trên -- trái), $\tfrac{2\cdot 6}{2}=6$ (góc trên -- phải),
$\tfrac{5\cdot 3}{2}=7{,}5$ (góc dưới -- phải): còn $18$. Vậy $S_{PBC}=54-18-18=18$. Đúng.}
\ans{a) $54$. \ b) $P(5;3)$. \ c) $13$ tam giác, ứng với $11$ vị trí của $P$.}
}

% ---------------------------------------------------------------- 47
\begin{probbox}{47}{\stars{1}}
\textit{(Mở rộng Bài 25.)} Ta lại xếp các hình chữ nhật $4\times 2$ thành ``cầu thang'' như Bài 25, nhưng lần này
đỉnh trên -- trái của mỗi hình (từ hình thứ hai) đặt tại điểm trên cạnh dưới của hình ngay phía trên, cách đỉnh
dưới -- trái của hình đó một khoảng $s$, với $s$ là $1$, $2$ hoặc $3$. (Bài 25 ứng với $s=2$.)
\begin{figC}\begin{scope}[x=0.38cm,y=0.38cm]
\foreach \i in {0,1,2,3}{\draw[Ink,line width=0.9pt] (\i,-2*\i) rectangle (\i+4,-2*\i+2);}
\draw[tmBlue,<->] (0,-0.5)--(1,-0.5);\node[below,font=\footnotesize] at (0.5,-0.5){$s$};
\node[above] at (2,2){$4$};\node[left] at (0,1){$2$};
\foreach \k in {0,1,2}{\fill[Ink] (5.5+0.5*\k,-6.6-0.5*\k) circle (2.3pt);}
\end{scope}\end{figC}
\begin{enumerate}[label=\alph*),leftmargin=1.6em,itemsep=1pt]
\item Với $s=1$ và $10$ hình chữ nhật, tính chu vi của cả hình.
\item Tìm công thức tính chu vi theo số hình $n$ và khoảng $s$.
\item Có thể chọn $s\in\{1;2;3\}$ và số hình $n$ để chu vi đúng bằng $100$ không? Tìm tất cả các cách.
\end{enumerate}
\end{probbox}
\probsol{
\textbf{Phân tích đề.} Ở Bài 25 ta đã thấy chu vi ``cầu thang'' bằng chu vi hình chữ nhật nhỏ nhất bao quanh
nó. Lập luận đó vẫn đúng với mọi $s$ (vì đường viền trên -- phải chỉ đi sang phải và đi xuống; đường viền dưới --
trái cũng vậy). Vậy chỉ cần tính bề ngang và chiều cao của cả hình.
\begin{recallbox}Hình ``cầu thang'' (đường viền chỉ đi sang phải/xuống ở mỗi phía) có chu vi bằng chu vi hình chữ
nhật bao quanh nó: $2\times(\text{bề ngang}+\text{chiều cao})$.\end{recallbox}

\textbf{b)} (làm trước vì câu a) suy ra ngay)
\step{1}{Bề ngang và chiều cao}
Mỗi hình mới lùi sang phải $s$. Hình đầu chiếm bề ngang $4$, mỗi hình sau thêm $s$:
bề ngang $=4+(n-1)s$. Mỗi hình mới thấp xuống $2$: chiều cao $=2n$.
\step{2}{Công thức}
\[ C=2\big(4+(n-1)s+2n\big). \]

\textbf{a)}
\step{3}{Thay $s=1$, $n=10$}
Bề ngang $4+9\cdot 1=13$, chiều cao $20$; $C=2(13+20)=66$.

\textbf{c)}
\step{4}{Xét từng giá trị của $s$}
\begin{itemize}[leftmargin=1.4em,itemsep=1pt]
\item $s=1$: $C=2(4+n-1+2n)=2(3n+3)=6n+6$. $6n+6=100\Rightarrow 6n=94$: không chia hết, loại.
\item $s=2$: $C=2(4+2n-2+2n)=2(4n+2)=8n+4$. $8n+4=100\Rightarrow n=12$. Nhận.
\item $s=3$: $C=2(4+3n-3+2n)=2(5n+1)=10n+2$. $10n+2=100\Rightarrow 10n=98$: loại.
\end{itemize}
Chỉ có một cách: $s=2$ và $n=12$.
\begin{warnbox}Nhớ rằng chiều cao $2n$ không phụ thuộc $s$; chỉ bề ngang thay đổi. Nhiều bạn cộng nhầm cả
những đoạn bị che bên trong.\end{warnbox}
\verify{Câu a) đếm trực tiếp: cạnh đứng $10\cdot 4=40$; cạnh ngang phía trên $4+9\cdot 1=13$; phía dưới
$9\cdot 1+4=13$. Tổng $66$. Đúng. Câu c): $s=2$, $n=12$: $8\cdot 12+4=100$. Đúng.}
\ans{a) $66$. \ b) $C=2\big(4+(n-1)s+2n\big)$. \ c) Chỉ có $s=2$, $n=12$.}
}

% ---------------------------------------------------------------- 48
\begin{probbox}{48}{\stars{1}}
\textit{(Mở rộng Bài 27.)} Hình dưới đây gồm $15$ hình bình hành nhỏ giống nhau, xếp thành $3$ hàng, mỗi hàng
$5$ hình. Hình bình hành nhỏ ở chính giữa được tô màu.
\begin{figC}\begin{scope}[x=0.5cm,y=0.5cm]
\fill[tmOrange!45] (5.4,1)--(7.6,1)--(8.6,2)--(6.4,2)--cycle;
\foreach \k in {0,...,5}{\draw[Ink,line width=0.9pt] (2.2*\k,0)--(2.2*\k+3,3);}
\foreach \j in {0,...,3}{\draw[Ink,line width=0.9pt] (\j,\j)--(\j+11,\j);}
\end{scope}\end{figC}
\begin{enumerate}[label=\alph*),leftmargin=1.6em,itemsep=1pt]
\item Trong hình có tất cả bao nhiêu hình bình hành?
\item Có bao nhiêu hình bình hành chứa hình bình hành nhỏ được tô màu?
\item Có bao nhiêu hình bình hành được ghép từ \emph{đúng} $4$ hình bình hành nhỏ?
\end{enumerate}
\end{probbox}
\probsol{
\textbf{Phân tích đề.} Hình có $4$ đường ngang và $6$ đường xiên. Như Bài 27, mỗi hình bình hành ứng với một
cách chọn $2$ đường ngang (cạnh trên, cạnh dưới) và $2$ đường xiên (cạnh trái, cạnh phải).
\begin{recallbox}Số cách chọn $2$ đường từ $k$ đường là $\dfrac{k(k-1)}{2}$.\end{recallbox}

\textbf{a)}
\step{1}{Đếm}
Chọn $2$ trong $4$ đường ngang: $\tfrac{4\cdot 3}{2}=6$ cách; chọn $2$ trong $6$ đường xiên: $\tfrac{6\cdot5}{2}=15$
cách. Số hình bình hành: $6\times 15=90$.

\textbf{b)}
\step{2}{Điều kiện để chứa ô tô màu}
Đánh số các đường ngang từ dưới lên là $0,1,2,3$; ô tô màu nằm giữa đường $1$ và đường $2$. Hình bình hành chứa ô
này khi cạnh dưới là đường $0$ hoặc $1$ ($2$ cách) và cạnh trên là đường $2$ hoặc $3$ ($2$ cách).
Đánh số các đường xiên từ trái sang là $0,1,\dots,5$; ô tô màu nằm giữa đường $2$ và đường $3$. Cạnh trái là đường
$0$, $1$ hoặc $2$ ($3$ cách); cạnh phải là đường $3$, $4$ hoặc $5$ ($3$ cách).
\step{3}{Nhân lại}
$2\times 2\times 3\times 3=36$ hình bình hành.
\begin{proofbox}Vì sao được nhân? Bốn lựa chọn (dưới, trên, trái, phải) độc lập với nhau: chọn thế nào cũng được
đúng một hình bình hành chứa ô tô màu, và hình khác nhau ứng với lựa chọn khác nhau.\end{proofbox}

\textbf{c)}
\step{4}{Các ``kích thước'' gồm đúng $4$ ô}
Một hình bình hành cao $h$ hàng, rộng $w$ cột gồm $h\times w$ ô. Cần $h\times w=4$ với $h\le 3$, $w\le 5$:
\begin{itemize}[leftmargin=1.4em,itemsep=1pt]
\item $h=1$, $w=4$: chọn hàng ($3$ cách), chọn $4$ cột liền nhau trong $5$ cột ($2$ cách): $3\cdot 2=6$;
\item $h=2$, $w=2$: chọn $2$ hàng liền nhau ($2$ cách), chọn $2$ cột liền nhau ($4$ cách): $2\cdot 4=8$;
\item $h=4$, $w=1$: không được vì chỉ có $3$ hàng.
\end{itemize}
Tổng: $6+8=14$.
\begin{warnbox}Đừng quên kiểu $2\times 2$! Và hình cao $4$ hàng thì không có.\end{warnbox}
\verify{Câu a) đếm theo hàng: hình cao $1$ hàng: $3\cdot 15=45$; cao $2$ hàng: $2\cdot 15=30$; cao $3$ hàng: $15$
(mỗi dải có $5+4+3+2+1=15$ cách chọn cột). Tổng $90$. Khớp.}
\ans{a) $90$. \ b) $36$. \ c) $14$.}
}

% ---------------------------------------------------------------- 49
\begin{probbox}{49}{\stars{2}}
\textit{(Mở rộng Bài 28, 29.)} Một hình hộp chữ nhật kích thước $4\times 5\times 6$ được ghép từ $120$ khối lập
phương nhỏ cạnh $1$. Người ta sơn tất cả các mặt ngoài của hình hộp, rồi tách rời các khối nhỏ.
\begin{enumerate}[label=\alph*),leftmargin=1.6em,itemsep=1pt]
\item Có bao nhiêu khối nhỏ được sơn đúng $3$ mặt? đúng $2$ mặt? đúng $1$ mặt? không được sơn mặt nào?
\item Tổng số mặt nhỏ (hình vuông $1\times 1$) được sơn là bao nhiêu? Giải thích vì sao kết quả bằng diện tích
toàn phần của hình hộp.
\item Một hình lập phương $n\times n\times n$ ($n\ge 3$) cũng được sơn và tách như vậy. Tìm $n$ để số khối nhỏ được
sơn đúng $1$ mặt bằng số khối nhỏ không được sơn mặt nào.
\end{enumerate}
\end{probbox}
\probsol{
\textbf{Phân tích đề.} Số mặt được sơn của một khối nhỏ phụ thuộc vào \emph{vị trí} của nó: ở đỉnh, trên cạnh,
trên mặt hay nằm hẳn bên trong. Quan sát mặt trước (kích thước $6\times 4$) để thấy rõ:
\begin{figC}\begin{scope}[x=0.5cm,y=0.5cm]
\foreach \x in {0,...,5}{\foreach \z in {0,...,3}{
  \pgfmathtruncatemacro\ex{(\x==0||\x==5)?1:0}\pgfmathtruncatemacro\ez{(\z==0||\z==3)?1:0}
  \pgfmathtruncatemacro\t{\ex+\ez}
  \ifnum\t=2 \fill[tmOrange!60] (\x,\z) rectangle (\x+1,\z+1);\fi
  \ifnum\t=1 \fill[tmGreen!40] (\x,\z) rectangle (\x+1,\z+1);\fi
  \ifnum\t=0 \fill[tmBlue!30] (\x,\z) rectangle (\x+1,\z+1);\fi}}
\draw[Ink] (0,0) grid[step=1] (6,4);
\foreach \x in {0,...,6}{\draw[Ink] (\x,4)--(\x+2.5,5.75);}
\foreach \d in {0,...,5}{\draw[Ink] (0.5*\d,4+0.35*\d)--(6+0.5*\d,4+0.35*\d);}
\foreach \z in {0,...,4}{\draw[Ink] (6,\z)--(8.5,\z+1.75);}
\foreach \d in {0,...,5}{\draw[Ink] (6+0.5*\d,0.35*\d)--(6+0.5*\d,4+0.35*\d);}
\draw[Ink,line width=1.1pt] (0,0)--(6,0)--(8.5,1.75)--(8.5,5.75)--(2.5,5.75)--(0,4)--cycle (0,4)--(6,4)--(6,0) (6,4)--(8.5,5.75);
\node[below,font=\footnotesize] at (3,0){$6$};\node[left,font=\footnotesize] at (0,2){$4$};\node[right,font=\footnotesize] at (7.3,0.6){$5$};
\begin{scope}[xshift=5.2cm,yshift=0.4cm,font=\footnotesize]
\fill[tmOrange!60] (0,3) rectangle (0.6,3.6);\node[right] at (0.7,3.3){đỉnh: $3$ mặt sơn};
\fill[tmGreen!40] (0,2) rectangle (0.6,2.6);\node[right] at (0.7,2.3){trên cạnh: $2$ mặt};
\fill[tmBlue!30] (0,1) rectangle (0.6,1.6);\node[right] at (0.7,1.3){giữa một mặt: $1$ mặt};
\end{scope}
\end{scope}\end{figC}

\textbf{a)}
\step{1}{Khối ở đỉnh: $3$ mặt sơn}
Hình hộp có $8$ đỉnh, mỗi đỉnh có đúng một khối: $8$ khối.
\step{2}{Khối trên cạnh (trừ hai đầu): $2$ mặt sơn}
Hình hộp có $12$ cạnh: $4$ cạnh dài $4$, $4$ cạnh dài $5$, $4$ cạnh dài $6$. Trên một cạnh dài $L$ có $L$ khối, bỏ
$2$ khối ở hai đầu (đã đếm ở đỉnh) còn $L-2$ khối:
\[ 4\cdot(4-2)+4\cdot(5-2)+4\cdot(6-2)=8+12+16=36. \]
\step{3}{Khối ở giữa một mặt: $1$ mặt sơn}
Trên một mặt $p\times q$, bỏ đi viền ngoài còn $(p-2)(q-2)$ ô. Hình hộp có $2$ mặt $4\times5$, $2$ mặt $4\times 6$,
$2$ mặt $5\times 6$:
\[ 2\cdot(2\cdot 3)+2\cdot(2\cdot 4)+2\cdot(3\cdot 4)=12+16+24=52. \]
\step{4}{Khối nằm hẳn bên trong: không sơn}
Bóc đi ``lớp vỏ'' dày $1$ ở mọi phía, còn lại một hình hộp $(4-2)\times(5-2)\times(6-2)=2\times 3\times 4=24$ khối.
\verify{$8+36+52+24=120$ -- đúng bằng tổng số khối.}

\textbf{b)}
\step{5}{Đếm mặt nhỏ được sơn}
$8\cdot 3+36\cdot 2+52\cdot 1+24\cdot 0=24+72+52=148$.
\begin{recallbox}Diện tích toàn phần hình hộp $a\times b\times c$ là $2(ab+bc+ca)$.\end{recallbox}
$2(4\cdot5+5\cdot6+6\cdot4)=2(20+30+24)=148$. Hai kết quả bằng nhau vì mỗi ô vuông $1\times 1$ trên mặt ngoài
hình hộp chính là đúng một mặt nhỏ được sơn của đúng một khối -- ta chỉ đếm cùng một tập hợp theo hai cách.

\textbf{c)}
\step{6}{Lập phương trình}
Với hình lập phương $n\times n\times n$: mỗi mặt có $(n-2)^2$ ô ở giữa, $6$ mặt cho $6(n-2)^2$ khối sơn đúng $1$
mặt; phần bên trong là hình lập phương $(n-2)^3$ khối. Cần:
\[ (n-2)^3=6(n-2)^2. \]
\step{7}{Giải}
Vì $n\ge 3$ nên $(n-2)^2>0$; chia hai vế cho $(n-2)^2$: $n-2=6$, tức $n=8$.
\begin{warnbox}Điều kiện $n\ge 3$ quan trọng: với $n=2$ thì cả hai số đều bằng $0$ (``bằng nhau'' một cách tầm
thường), và ta không được chia cho $0$.\end{warnbox}
\verify{$n=8$: $6\cdot 6^2=216$ và $6^3=216$. Đúng.}
\ans{a) $8$; $36$; $52$; $24$ khối. \ b) $148$. \ c) $n=8$.}
}

% ---------------------------------------------------------------- 50
\begin{probbox}{50}{\stars{2}}
\textit{(Mở rộng Bài 19, 25.)} Ghép $10$ ô vuông cạnh $1$ thành một hình phẳng sao cho: hai ô kề nhau thì ghép
trùng khít một cạnh, và hình ``liền một khối'' (từ ô bất kì có thể đi đến ô bất kì khác qua các cạnh chung).
Chu vi của hình là tổng độ dài các cạnh ô không bị ô khác che.
\begin{figC}\begin{scope}[x=0.38cm,y=0.38cm]
\foreach \c in {(0,0),(1,0),(2,0),(2,1),(3,1),(4,1),(4,2),(1,1),(4,0),(5,1)}{\fill[tmBlue!30,draw=Ink] \c rectangle ++(1,1);}
\node[font=\footnotesize] at (3,-0.8){một ví dụ};
\end{scope}\end{figC}
\begin{enumerate}[label=\alph*),leftmargin=1.6em,itemsep=1pt]
\item Chứng tỏ chu vi luôn là một số chẵn.
\item Chu vi lớn nhất có thể là bao nhiêu?
\item Chu vi nhỏ nhất có thể là bao nhiêu?
\item Những số nào có thể là chu vi của một hình như vậy?
\end{enumerate}
\end{probbox}
\probsol{
\textbf{Phân tích đề.} Gọi $k$ là số \emph{cạnh chung} (cạnh mà hai ô ghép vào nhau). Ý chính: mỗi cạnh chung
làm ``mất'' hai cạnh ô khỏi chu vi.

\textbf{a)}
\step{1}{Công thức chu vi}
$10$ ô có tất cả $10\cdot 4=40$ cạnh ô. Mỗi cạnh chung là cạnh của hai ô và không nằm trên chu vi; mọi cạnh ô
còn lại nằm trên chu vi. Do đó
\[ C=40-2k. \]
$40$ và $2k$ đều chẵn nên $C$ chẵn.

\textbf{b)}
\step{2}{Hình liền khối có ít nhất $9$ cạnh chung}
\begin{proofbox}Ta dựng lại hình bằng cách đặt từng ô một: bắt đầu từ một ô bất kì, mỗi lần đặt thêm một ô kề với
một ô đã đặt (luôn làm được vì hình liền khối). Mỗi ô thêm vào tạo ra ít nhất $1$ cạnh chung mới. Có $9$ lần đặt
thêm nên $k\ge 9$.\end{proofbox}
\step{3}{Kết luận}
$C=40-2k\le 40-18=22$. Dấu bằng xảy ra với hàng thẳng $1\times 10$: chu vi $2(10+1)=22$. Vậy chu vi lớn nhất là $22$.

\textbf{c)}
\step{4}{So sánh với hình chữ nhật bao quanh}
Gọi hình chữ nhật nhỏ nhất (cạnh theo lưới) chứa cả hình là $a\times b$ ($a$ cột, $b$ hàng).
\begin{proofbox}Mỗi cột của hình chữ nhật bao có ít nhất một ô. Ô cao nhất trong cột có cạnh trên nằm trên chu vi,
ô thấp nhất có cạnh dưới nằm trên chu vi: mỗi cột góp ít nhất $2$ đơn vị cạnh ngang. Tương tự mỗi hàng góp ít nhất
$2$ đơn vị cạnh đứng. Vậy $C\ge 2a+2b=2(a+b)$.\end{proofbox}
\step{5}{Chặn $a+b$}
Hình chữ nhật bao chứa $10$ ô nên $ab\ge 10$. Nếu $a+b\le 6$ thì $ab\le 9$ (các cặp có tổng $6$: $1\cdot5=5$,
$2\cdot4=8$, $3\cdot3=9$; tổng nhỏ hơn còn cho tích nhỏ hơn) -- vô lí. Vậy $a+b\ge 7$ và $C\ge 14$.
\step{6}{Có hình đạt $14$}
Hình chữ nhật $2\times 5$: chu vi $2(2+5)=14$. Vậy chu vi nhỏ nhất là $14$.

\textbf{d)}
\step{7}{Các giá trị có thể}
Từ a), b), c): $C$ chẵn và $14\le C\le 22$, nên $C\in\{14;16;18;20;22\}$. Mỗi giá trị đều đạt được bằng các hình
``hai hàng'' sau (dùng ý của Bài 25: đường viền là cầu thang nên chu vi bằng chu vi hình chữ nhật bao quanh):
\begin{figC}\begin{scope}[x=0.3cm,y=0.3cm]
\foreach \top/\bot/\xs/\lab in {5/5/0/14,6/4/6.5/16,7/3/14/18,8/2/22.5/20}{
  \foreach \i in {1,...,\top}{\fill[tmBlue!30,draw=Ink] (\xs+\i-1,1) rectangle ++(1,1);}
  \foreach \i in {1,...,\bot}{\fill[tmBlue!30,draw=Ink] (\xs+\i-1,0) rectangle ++(1,1);}
  \node[font=\footnotesize] at (\xs+\top/2,-0.9){$C=\lab$};}
\foreach \i in {1,...,10}{\fill[tmGreen!35,draw=Ink] (\i-1,-4) rectangle ++(1,1);}
\node[font=\footnotesize,anchor=west] at (10.5,-3.5){$C=22$};
\end{scope}\end{figC}
Hàng trên $5,6,7,8$ ô, hàng dưới $5,4,3,2$ ô: hình chữ nhật bao là $5\times2$, $6\times 2$, $7\times 2$, $8\times 2$,
chu vi lần lượt $14,16,18,20$; và hàng thẳng cho $22$.
\begin{warnbox}Không phải cứ ``gọn'' là chu vi nhỏ: hình $3\times 3$ thêm một ô lồi ra có chu vi $12+2=14$, bằng
hình $2\times 5$, chứ không nhỏ hơn.\end{warnbox}
\verify{Hình $6+4$: $k=5+3+4=12$, $C=40-24=16$. Hình $8+2$: $k=7+1+2=10$, $C=40-20=20$. Đúng.}
\ans{a) $C=40-2k$ chẵn. \ b) $22$. \ c) $14$. \ d) $14$, $16$, $18$, $20$, $22$.}
}

\end{document}
